%% =====================================================================
%%  Kubernetes From First Principles
%%  A Technical Learning Guide to "Kubernetes The Hard Way"
%%
%%  House style: Charter / Helvetica, navy-steel-teal-brick palette,
%%  tcolorbox callouts, booktabs tables, fancyhdr. pdflatex, no shell-escape.
%% =====================================================================
\documentclass[11pt,oneside]{report}

%% --- encoding and fonts ---------------------------------------------
\usepackage[T1]{fontenc}
\usepackage[utf8]{inputenc}
\usepackage{textcomp}
\usepackage{charter}                    % Bitstream Charter (serif body)
\usepackage[scaled=0.92]{helvet}        % Helvetica (headings, labels)
\usepackage{courier}                    % Courier (code)
\usepackage{microtype}

%% --- page geometry ---------------------------------------------------
\usepackage[letterpaper,
            top=1.05in, bottom=1.15in, left=1.15in, right=1.15in,
            headheight=15pt, headsep=20pt, footskip=30pt]{geometry}

%% --- colour palette --------------------------------------------------
\usepackage{xcolor}
\definecolor{navy}  {HTML}{102A52}
\definecolor{steel} {HTML}{465F7D}
\definecolor{teal}  {HTML}{0D6E6E}
\definecolor{brick} {HTML}{962C26}
\definecolor{rule}  {HTML}{C9D2DE}
\definecolor{codebg}{HTML}{F5F7FA}
\definecolor{outbg} {HTML}{F0F1F2}
\definecolor{navytint}{HTML}{EEF2F8}
\definecolor{tealtint}{HTML}{ECF5F4}
\definecolor{bricktint}{HTML}{FBEFEE}
\definecolor{steeltint}{HTML}{F1F3F7}

%% --- tables ----------------------------------------------------------
\usepackage{calc}
\usepackage{array}
\usepackage{booktabs}
\usepackage{longtable}
\usepackage{tabularx}
\usepackage{ragged2e}
\setlength{\tabcolsep}{6pt}
\renewcommand{\arraystretch}{1.18}
\usepackage[font=small,labelfont={bf,sf},labelsep=period]{caption}
\usepackage{etoolbox}
\AtBeginEnvironment{longtable}{\footnotesize\setlength{\parskip}{0pt}}
% encourage breaks inside long paths and flag names
\sloppy
\hfuzz=1pt
% Ragged table cells and wrapped listing lines are loose by construction;
% silence the cosmetic badness reports so a genuine overfull still stands out.
\hbadness=10000
\vbadness=10000

%% --- lists -----------------------------------------------------------
\usepackage{enumitem}
\setlist{itemsep=2pt, parsep=0pt, topsep=4pt}
\providecommand{\tightlist}{%
  \setlength{\itemsep}{0pt}\setlength{\parskip}{0pt}}

%% --- code listings ---------------------------------------------------
\usepackage{listings}
\usepackage{upquote}
\lstset{
  basicstyle=\footnotesize\ttfamily,
  breaklines=true,
  breakatwhitespace=false,
  breakindent=2.2em,
  postbreak=\mbox{\textcolor{steel}{\raisebox{0pt}[0pt][0pt]{$\hookrightarrow$}}\space},
  columns=fullflexible,
  keepspaces=true,
  showstringspaces=false,
  upquote=true,
  xleftmargin=10pt,
  xrightmargin=4pt,
  framexleftmargin=10pt,
  aboveskip=10pt,
  belowskip=10pt,
  frame=leftline,
  framerule=2pt,
}

\lstdefinestyle{cmdstyle}{
  alsoletter={-},
  backgroundcolor=\color{codebg},
  rulecolor=\color{navy},
  basicstyle=\footnotesize\ttfamily\color{black!88},
  commentstyle=\color{steel}\itshape,
  morecomment=[l]{\#},
  keywordstyle=\color{navy}\bfseries,
  morekeywords={ssh,scp,sudo,systemctl,journalctl,kubectl,openssl,etcdctl,
    curl,wget,tar,chmod,mkdir,mv,cp,rm,cat,echo,grep,sed,cut,for,do,done,
    while,read,export,apt-get,git,swapoff,swapon,modprobe,sysctl,hostnamectl,
    ssh-keygen,ssh-copy-id,envsubst,ip,iptables,ss,lsmod,crictl,pwd,ls,
    hexdump,head,dpkg,base64,stat,getent,file,ping,dmesg}
}

\lstdefinestyle{outstyle}{
  backgroundcolor=\color{outbg},
  rulecolor=\color{rule},
  basicstyle=\scriptsize\ttfamily\color{black!75}
}

\lstnewenvironment{cmdlst}{\lstset{style=cmdstyle}}{}
\lstnewenvironment{outlst}{\lstset{style=outstyle}}{}

%% --- callout boxes ---------------------------------------------------
\usepackage[most]{tcolorbox}

\newtcolorbox{whybox}{
  enhanced, breakable, sharp corners=downhill,
  colback=navytint, colframe=navy, boxrule=0pt, leftrule=3pt,
  left=9pt, right=9pt, top=7pt, bottom=7pt,
  fontupper=\small,
  borderline west={3pt}{0pt}{navy}}

\newtcolorbox{tipbox}{
  enhanced, breakable, sharp corners=downhill,
  colback=tealtint, colframe=teal, boxrule=0pt,
  left=9pt, right=9pt, top=7pt, bottom=7pt,
  fontupper=\small,
  borderline west={3pt}{0pt}{teal}}

\newtcolorbox{notebox}{
  enhanced, breakable, sharp corners=downhill,
  colback=steeltint, colframe=steel, boxrule=0pt,
  left=9pt, right=9pt, top=7pt, bottom=7pt,
  fontupper=\small,
  borderline west={3pt}{0pt}{steel}}

\newtcolorbox{riskbox}{
  enhanced, breakable, sharp corners=downhill,
  colback=bricktint, colframe=brick, boxrule=0pt,
  left=9pt, right=9pt, top=7pt, bottom=7pt,
  fontupper=\small,
  borderline west={3pt}{0pt}{brick}}

\newtcolorbox{beyondbox}{
  enhanced, breakable, sharp corners=downhill,
  colback=white, colframe=teal, boxrule=0.4pt,
  left=9pt, right=9pt, top=7pt, bottom=7pt,
  fontupper=\small,
  borderline west={3pt}{0pt}{teal}}

% Short-form macros, matching the repo's callout family.
\newcommand{\why}[1]{\begin{whybox}#1\end{whybox}}
\newcommand{\tip}[1]{\begin{tipbox}#1\end{tipbox}}
\newcommand{\note}[1]{\begin{notebox}#1\end{notebox}}
\newcommand{\risknote}[1]{\begin{riskbox}#1\end{riskbox}}
\newcommand{\beyond}[1]{\begin{beyondbox}#1\end{beyondbox}}

\newcommand{\calloutlabel}[2]{{\sffamily\bfseries\footnotesize\color{#1}#2\quad}}

%% --- sectioning ------------------------------------------------------
\usepackage{titlesec}

\titleformat{\chapter}[display]
  {\sffamily}
  {\raggedright\color{steel}\footnotesize\bfseries\MakeUppercase{\chaptertitlename\ \thechapter}}
  {6pt}
  {\color{navy}\LARGE\bfseries\raggedright}
  [\vspace{4pt}{\color{rule}\titlerule[1.2pt]}]
\titlespacing*{\chapter}{0pt}{-14pt}{26pt}

\titleformat{\section}
  {\sffamily\large\bfseries\color{navy}}{\thesection}{0.7em}{}
\titlespacing*{\section}{0pt}{16pt}{6pt}

\titleformat{\subsection}
  {\sffamily\normalsize\bfseries\color{steel}}{\thesubsection}{0.7em}{}
\titlespacing*{\subsection}{0pt}{12pt}{4pt}

\setcounter{secnumdepth}{2}
\setcounter{tocdepth}{1}

%% --- headers and footers ---------------------------------------------
\usepackage{fancyhdr}
\pagestyle{fancy}
\fancyhf{}
\fancyhead[L]{\sffamily\footnotesize\color{steel}\nouppercase{\leftmark}}
\fancyhead[R]{\sffamily\footnotesize\color{steel}Kubernetes From First Principles}
\fancyfoot[C]{\sffamily\footnotesize\color{steel}\thepage}
\renewcommand{\headrule}{\color{rule}\hrule height 0.6pt}
\renewcommand{\footrulewidth}{0pt}
\fancypagestyle{plain}{%
  \fancyhf{}%
  \fancyfoot[C]{\sffamily\footnotesize\color{steel}\thepage}%
  \renewcommand{\headrule}{}}
\renewcommand{\chaptermark}[1]{\markboth{#1}{}}

%% --- paragraph shape -------------------------------------------------
\setlength{\parindent}{0pt}
\setlength{\parskip}{6pt plus 2pt minus 1pt}
\setlength{\emergencystretch}{3em}
\widowpenalty=10000
\clubpenalty=10000

%% --- links -----------------------------------------------------------
\usepackage[
  colorlinks=true,
  linkcolor=navy,
  urlcolor=teal,
  citecolor=navy,
  pdftitle={Kubernetes From First Principles},
  pdfsubject={A Technical Learning Guide to Kubernetes The Hard Way},
  pdfkeywords={Kubernetes, etcd, containerd, CNI, PKI, RBAC, kubelet},
  pdfdisplaydoctitle=true]{hyperref}
\usepackage{bookmark}
\urlstyle{same}

%% --- small helpers ----------------------------------------------------
\newcommand{\upstreamlab}[1]{%
  {\sffamily\footnotesize\color{steel}Upstream lab:\quad\texttt{#1}}\par\vspace{2pt}%
  {\color{rule}\hrule height 0.4pt}\vspace{8pt}}

\newcommand{\meta}[2]{{\sffamily\footnotesize\color{steel}#1}\quad{#2}}

\begin{document}

%% =====================  TITLE PAGE  =================================
\begin{titlepage}
\thispagestyle{empty}
\sffamily
\vspace*{0.35in}

{\color{navy}\hrule height 3pt}
\vspace{18pt}

{\color{steel}\large\bfseries\MakeUppercase{Technical Learning Guide}}

\vspace{14pt}
{\color{navy}\fontsize{34}{40}\selectfont\bfseries Kubernetes\\From First Principles}

\vspace{16pt}
{\color{steel}\large A companion to \textit{Kubernetes The Hard Way}\\
building a cluster from individual binaries, by hand}

\vspace{14pt}
{\color{rule}\hrule height 1pt}

\vspace{28pt}
\rmfamily\normalsize
\begin{tabular}{@{}ll@{}}
\toprule
\multicolumn{2}{@{}l}{\sffamily\bfseries\color{navy}Target cluster}\\
\midrule
Kubernetes      & v1.32.3 \\
containerd      & v2.1.0-beta.0 \\
etcd            & v3.6.0-rc.3 \\
CNI plugins     & v1.6.2 \\
runc            & v1.3.0-rc.1 \\
crictl          & v1.32.0 \\
\midrule
Hosts           & 4 $\times$ Debian 12 (bookworm), ARM64 or AMD64 \\
Topology        & 1 jumpbox, 1 control plane node, 2 workers \\
Cluster CIDR    & \texttt{10.200.0.0/16} \\
Service CIDR    & \texttt{10.32.0.0/24} \\
\bottomrule
\end{tabular}

\vfill

\begin{beyondbox}
\calloutlabel{teal}{ATTRIBUTION}
This guide is a derivative learning companion to \textit{Kubernetes The Hard
Way} by Kelsey Hightower. The upstream documentation is licensed
\href{http://creativecommons.org/licenses/by-nc-sa/4.0/}{CC BY-NC-SA 4.0};
the repository's code and configuration are licensed Apache 2.0. Every
command, path, IP address, and version pin reproduced here comes from that
source. The explanations, failure analysis, diagnostic tables, and exercises
are additions, and material that does not originate upstream is marked
\textsf{BEYOND THE SOURCE} throughout.
\end{beyondbox}

\vspace{10pt}
{\color{navy}\hrule height 3pt}
\end{titlepage}

%% =====================  FRONT MATTER  ===============================
\pagenumbering{roman}
\setcounter{page}{1}
\pagestyle{fancy}

{\hypersetup{linkcolor=navy}
\renewcommand{\contentsname}{Contents}
\tableofcontents}

\cleardoublepage
\pagenumbering{arabic}
\setcounter{page}{1}
\setcounter{chapter}{-1}

\hypertarget{how-this-guide-works}{%
\chapter*{How This
Guide
Works}\label{how-this-guide-works}}
\addcontentsline{toc}{chapter}{How
This Guide Works}

\markboth{How This Guide Works}{}

\emph{Kubernetes The
Hard Way} is a build
script written in
prose. It tells you
what to type, in the
right order, and it
works. What it
deliberately does
not do is explain
why each step exists
or what the failure
looks like when you
get it wrong ---
because Kelsey's
stated goal is to
make you type the
steps and discover
that for yourself.

This guide sits
alongside it. Every
chapter maps
one-to-one onto one
upstream lab, in the
same order, with the
same commands.
What's added is the
layer underneath:
the concept each
step encodes, the
specific failure
mode when it's
skipped, and how to
confirm success
before moving on.

\hypertarget{who-this-is-for}{%
\section*{Who this
is
for}\label{who-this-is-for}}
\addcontentsline{toc}{section}{Who
this is for}

Engineers
comfortable with
Linux, systemd, TLS,
and basic IP
networking, who have
used Kubernetes as a
consumer
(\texttt{kubectl\ apply})
but have never
assembled one. No
prior cluster
operations
experience is
assumed. If you have
never issued an
X.509 certificate or
written a systemd
unit, you can still
follow along, but
budget extra time in
Chapters 4 and 7.

\hypertarget{what-you-need-before-starting}{%
\section*{What you
need before
starting}\label{what-you-need-before-starting}}
\addcontentsline{toc}{section}{What
you need before
starting}

Four machines
running
\textbf{Debian 12
(bookworm)}, all
ARM64 or all AMD64,
on the same L2
network with static
or reserved IPs,
each reachable from
the others:

\begin{longtable}[]{@{}
  >{\raggedright\arraybackslash}p{(\columnwidth - 8\tabcolsep) * \real{0.1667}}
  >{\raggedright\arraybackslash}p{(\columnwidth - 8\tabcolsep) * \real{0.4444}}
  >{\raggedright\arraybackslash}p{(\columnwidth - 8\tabcolsep) * \real{0.0926}}
  >{\raggedright\arraybackslash}p{(\columnwidth - 8\tabcolsep) * \real{0.1296}}
  >{\raggedright\arraybackslash}p{(\columnwidth - 8\tabcolsep) * \real{0.1667}}@{}}
\toprule\noalign{}
\begin{minipage}[b]{\linewidth}\raggedright
Name
\end{minipage} &
\begin{minipage}[b]{\linewidth}\raggedright
Role
\end{minipage} &
\begin{minipage}[b]{\linewidth}\raggedright
CPU
\end{minipage} &
\begin{minipage}[b]{\linewidth}\raggedright
RAM
\end{minipage} &
\begin{minipage}[b]{\linewidth}\raggedright
Storage
\end{minipage} \\
\midrule\noalign{}
\endhead
\bottomrule\noalign{}
\endlastfoot
jumpbox &
Administration host
& 1 & 512MB &
10GB \\
server & Kubernetes
server & 1 & 2GB &
20GB \\
node-0 & Kubernetes
worker node & 1 &
2GB & 20GB \\
node-1 & Kubernetes
worker node & 1 &
2GB & 20GB \\
\end{longtable}

You also need
roughly 500 MB of
download bandwidth
and outbound HTTPS
to
\texttt{dl.k8s.io}
and
\texttt{github.com}.

\hypertarget{time-budget}{%
\section*{Time
budget}\label{time-budget}}
\addcontentsline{toc}{section}{Time
budget}

\begin{longtable}[]{@{}
  >{\raggedright\arraybackslash}p{(\columnwidth - 4\tabcolsep) * \real{0.3333}}
  >{\raggedright\arraybackslash}p{(\columnwidth - 4\tabcolsep) * \real{0.3333}}
  >{\raggedright\arraybackslash}p{(\columnwidth - 4\tabcolsep) * \real{0.3333}}@{}}
\toprule\noalign{}
\begin{minipage}[b]{\linewidth}\raggedright
Chapter
\end{minipage} &
\begin{minipage}[b]{\linewidth}\raggedright
Topic
\end{minipage} &
\begin{minipage}[b]{\linewidth}\raggedright
Time
\end{minipage} \\
\midrule\noalign{}
\endhead
\bottomrule\noalign{}
\endlastfoot
1 & Prerequisites &
5 min (plus VM
provisioning) \\
2 & Jumpbox & 20 min
(mostly
downloading) \\
3 & Compute
resources & 25
min \\
4 & CA and TLS
certificates & 30
min \\
5 & Kubeconfigs & 20
min \\
6 & Encryption
config & 5 min \\
7 & etcd & 15 min \\
8 & Control plane &
40 min \\
9 & Worker nodes &
45 min \\
10 & Remote kubectl
& 10 min \\
11 & Pod network
routes & 15 min \\
12 & Smoke test & 25
min \\
13 & Cleanup & 5
min \\
\end{longtable}

Total: roughly 4--5
hours of hands-on
time for a first
pass, assuming
nothing breaks.
Something will
break. That is the
point.

\hypertarget{the-three-machine-contexts}{%
\section*{The three
machine
contexts}\label{the-three-machine-contexts}}
\addcontentsline{toc}{section}{The
three machine
contexts}

Almost every mistake
in this build comes
from running a
command on the wrong
host. There are
exactly three
contexts:

\begin{itemize}
\tightlist
\item
  \textbf{jumpbox}
  --- where you
  generate
  certificates,
  kubeconfigs, and
  push files out.
  Working directory
  is always
  \texttt{/}\allowbreak{}\texttt{root/}\allowbreak{}\texttt{kubernetes-}\allowbreak{}\texttt{the-}\allowbreak{}\texttt{hard-}\allowbreak{}\texttt{way}.
\item
  \textbf{server}
  --- the control
  plane host. You
  \texttt{ssh\ root@server}
  to get here.
\item
  \textbf{node-0 /
  node-1} --- the
  workers. You
  \texttt{ssh\ root@node-0}
  to get here, and
  you do the whole
  chapter 9 sequence
  \textbf{twice}.
\end{itemize}

Every code block in
this guide is
annotated with its
host. If a block has
no annotation, it
runs on the jumpbox.

\hypertarget{recovering-from-a-failed-step}{%
\section*{Recovering
from a failed
step}\label{recovering-from-a-failed-step}}
\addcontentsline{toc}{section}{Recovering
from a failed step}

This build has no
rollback. What it
does have is
idempotence in most
places and a short
dependency chain in
the rest. The
general recovery
procedure:

\begin{enumerate}
\def\labelenumi{\arabic{enumi}.}
\tightlist
\item
  \textbf{Identify
  the layer.}
  Certificate
  problem,
  kubeconfig
  problem, service
  problem, or
  network problem?
  The four layers
  fail differently
  and the diagnosis
  in each chapter
  tells you which
  symptom belongs
  where.
\item
  \textbf{Read the
  unit's logs.}
  Every component in
  this build runs
  under systemd.
  \texttt{journalctl\ -}\allowbreak{}\texttt{u\ \textless{}unit\textgreater{}\ -\/-no-}\allowbreak{}\texttt{pager\ -}\allowbreak{}\texttt{n\ 50}
  on the relevant
  host is the first
  command, always.
\item
  \textbf{Re-run the
  generating step on
  the jumpbox, then
  re-distribute.}
  Certificates and
  kubeconfigs are
  cheap to
  regenerate. If you
  regenerate a cert,
  you must also
  regenerate any
  kubeconfig that
  embedded it
  (\texttt{-\/-embed-}\allowbreak{}\texttt{certs=true}
  copies the cert
  \emph{into} the
  kubeconfig --- it
  does not reference
  it).
\item
  \textbf{Restart
  the consuming
  service.}
  \texttt{systemctl\ restart\ kube-}\allowbreak{}\texttt{apiserver}
  and friends. None
  of these
  components watch
  their credential
  files for changes.
\item
  \textbf{When in
  real doubt,
  destroy
  \texttt{server}
  and rebuild it.}
  Everything on the
  control plane host
  is reconstructable
  from the jumpbox
  in about ten
  minutes. Nothing
  you care about
  lives there.
\end{enumerate}

\begin{tipbox}\calloutlabel{teal}{RULE OF THUMB}

If you find yourself
editing a file on
\texttt{server} or a
worker directly,
stop. The jumpbox is
the source of truth;
edit there and
re-push.

\end{tipbox}

\hypertarget{chap-0}{%
\chapter{Architecture
Orientation}\label{chap-0}}

\hypertarget{what-youre-building}{%
\section{What you're
building}\label{what-youre-building}}

A
single-control-plane
Kubernetes cluster
with two worker
nodes, assembled
from individual
binaries running as
systemd services,
with a hand-rolled
PKI and hand-written
routing. No
installer, no cloud
provider, no
kubeadm, no
container images for
the control plane.
Every trust
relationship in the
cluster is one you
will create by hand.

\hypertarget{why-its-necessary}{%
\section{Why it's
necessary}\label{why-its-necessary}}

\begin{whybox}

Managed Kubernetes
hides two things
that account for
most real production
incidents:
\textbf{who is
allowed to talk to
whom, and how
packets get from one
pod to another.}
This build makes
both explicit. By
the end you will
have issued eight
certificates,
written five
kubeconfigs, and
added four routes
--- and you will
know exactly which
one is broken when
something stops
working.

\end{whybox}

\hypertarget{key-concepts}{%
\section{Key
concepts}\label{key-concepts}}

\textbf{1. The
control plane is
four processes, not
a monolith.}

\begin{longtable}[]{@{}
  >{\raggedright\arraybackslash}p{(\columnwidth - 4\tabcolsep) * \real{0.3333}}
  >{\raggedright\arraybackslash}p{(\columnwidth - 4\tabcolsep) * \real{0.3333}}
  >{\raggedright\arraybackslash}p{(\columnwidth - 4\tabcolsep) * \real{0.3333}}@{}}
\toprule\noalign{}
\begin{minipage}[b]{\linewidth}\raggedright
Process
\end{minipage} &
\begin{minipage}[b]{\linewidth}\raggedright
Job
\end{minipage} &
\begin{minipage}[b]{\linewidth}\raggedright
Talks to
\end{minipage} \\
\midrule\noalign{}
\endhead
\bottomrule\noalign{}
\endlastfoot
\texttt{etcd} &
Stores all cluster
state as key-value
data & Only the API
server \\
\texttt{kube-apiserver}
& The only component
that reads/writes
etcd; authenticates
and authorizes every
request & etcd,
kubelets, all
clients \\
\texttt{kube-}\allowbreak{}\texttt{controller-}\allowbreak{}\texttt{manager}
& Runs
reconciliation loops
(deployments,
replicasets, service
accounts, node
lifecycle) & API
server only \\
\texttt{kube-scheduler}
& Assigns
unscheduled pods to
nodes & API server
only \\
\end{longtable}

The controller
manager and
scheduler have no
special access. They
are ordinary API
clients with
certificates,
exactly like
\texttt{kubectl}.
This is the single
most important
structural fact
about Kubernetes:
\textbf{the API
server is the hub,
and everything else
is a spoke.}

\textbf{2. The
worker is three
processes.}

\begin{longtable}[]{@{}
  >{\raggedright\arraybackslash}p{(\columnwidth - 2\tabcolsep) * \real{0.5000}}
  >{\raggedright\arraybackslash}p{(\columnwidth - 2\tabcolsep) * \real{0.5000}}@{}}
\toprule\noalign{}
\begin{minipage}[b]{\linewidth}\raggedright
Process
\end{minipage} &
\begin{minipage}[b]{\linewidth}\raggedright
Job
\end{minipage} \\
\midrule\noalign{}
\endhead
\bottomrule\noalign{}
\endlastfoot
\texttt{containerd}
& Runs containers
(via \texttt{runc}),
pulls images,
manages snapshots \\
\texttt{kubelet} &
Watches the API
server for pods
assigned to this
node; tells
containerd what to
run; reports
node/pod status
back \\
\texttt{kube-proxy}
& Programs
\texttt{iptables}
rules so that
Service ClusterIPs
and NodePorts
resolve to pod
IPs \\
\end{longtable}

The kubelet does not
route traffic.
\texttt{kube-proxy}
does not run
containers. They are
independent
consumers of the
same API.

\textbf{3. Three
address spaces, and
they must not
overlap.}

\begin{longtable}[]{@{}
  >{\raggedright\arraybackslash}p{(\columnwidth - 4\tabcolsep) * \real{0.3333}}
  >{\raggedright\arraybackslash}p{(\columnwidth - 4\tabcolsep) * \real{0.3333}}
  >{\raggedright\arraybackslash}p{(\columnwidth - 4\tabcolsep) * \real{0.3333}}@{}}
\toprule\noalign{}
\begin{minipage}[b]{\linewidth}\raggedright
Range
\end{minipage} &
\begin{minipage}[b]{\linewidth}\raggedright
Purpose
\end{minipage} &
\begin{minipage}[b]{\linewidth}\raggedright
Where it's
configured
\end{minipage} \\
\midrule\noalign{}
\endhead
\bottomrule\noalign{}
\endlastfoot
Your LAN
(e.g.~\texttt{192.168.x.0/24})
& Node IPs & Your
hypervisor / DHCP
reservations \\
\texttt{10.200.0.0/16}
& Pod network
(cluster CIDR) &
\texttt{kube-}\allowbreak{}\texttt{controller-}\allowbreak{}\texttt{manager\ -\/-cluster-}\allowbreak{}\texttt{cidr},
\texttt{kube-}\allowbreak{}\texttt{proxy-}\allowbreak{}\texttt{config.}\allowbreak{}\texttt{yaml} \\
\texttt{10.32.0.0/24}
& Service network
(ClusterIPs) &
\texttt{kube-}\allowbreak{}\texttt{controller-}\allowbreak{}\texttt{manager\ -\/-service-}\allowbreak{}\texttt{cluster-}\allowbreak{}\texttt{ip-}\allowbreak{}\texttt{range} \\
\end{longtable}

The pod CIDR is
subdivided per node:
\texttt{node-0} gets
\texttt{10.200.0.0/24},
\texttt{node-1} gets
\texttt{10.200.1.0/24}.
The service range's
first address,
\texttt{10.32.0.1},
is reserved for the
\texttt{kubernetes}
Service that fronts
the API server ---
which is why it
appears as a SAN in
the API server's
certificate.

\textbf{4. Two
independent auth
systems meet at the
API server.}

\emph{Authentication}
answers ``who are
you?'' --- here,
always via X.509
client certificate.
The certificate's
\textbf{Common Name
becomes the
username} and its
\textbf{Organization
becomes the group}.
\emph{Authorization}
answers ``may you do
this?'' --- here via
two modes,
\texttt{Node,RBAC},
evaluated in order.

This is why the
certificate subjects
in \texttt{ca.conf}
are not cosmetic.
\texttt{CN=admin,\ O=system:}\allowbreak{}\texttt{masters}
grants cluster-admin
because Kubernetes
ships a built-in
ClusterRoleBinding
for the
\texttt{system:masters}
group.
\texttt{CN=system:}\allowbreak{}\texttt{node:}\allowbreak{}\texttt{node-}\allowbreak{}\texttt{0,\ O=system:}\allowbreak{}\texttt{nodes}
is what the Node
authorizer looks
for. Get a subject
wrong and you get a
\texttt{403\ Forbidden}
from a component
that appears, by
every other measure,
to be correctly
configured.

\hypertarget{the-topology}{%
\section{The
topology}\label{the-topology}}

\begin{outlst}
                        +-------------------------------+
                        |          jumpbox              |
                        |  (no cluster role at runtime) |
                        |                               |
                        |  * openssl -> CA + 8 certs    |
                        |  * kubectl -> 5 kubeconfigs   |
                        |  * downloads/ (all binaries)  |
                        |  * machines.txt (source of    |
                        |    truth for IPs and subnets) |
                        +---------------+---------------+
                                        | ssh / scp (root, key-based)
                +-----------------------+-----------------------+
                |                       |                       |
                v                       v                       v
   +------------------------+  +------------------+  +------------------+
   |        server          |  |      node-0      |  |      node-1      |
   |                        |  |                  |  |                  |
   |  etcd  :2379 (loopback)|  |  containerd      |  |  containerd      |
   |    ^                   |  |    ^             |  |    ^             |
   |    | http://127.0.0.1  |  |    | CRI (unix)  |  |    | CRI (unix)  |
   |  kube-apiserver :6443 -+--+-> kubelet :10250 |  |  kubelet :10250  |
   |    ^          ^        |  |    |             |  |    |             |
   |    |          |        |  |  kube-proxy      |  |  kube-proxy      |
   |  kube-        kube-    |  |    | iptables    |  |    | iptables    |
   |  controller-  scheduler|  |    v             |  |    v             |
   |  manager               |  |  cni0 bridge     |  |  cni0 bridge     |
   |                        |  |  10.200.0.0/24   |  |  10.200.1.0/24   |
   +------------------------+  +------------------+  +------------------+
            |                          ^                      ^
            +--------------------------+----------------------+
                   static ip routes (Chapter 11) carry
                   pod-to-pod traffic between nodes
\end{outlst}

Note the arrow
directions. The API
server initiates
connections
\emph{to} kubelets
(for \texttt{logs},
\texttt{exec},
\texttt{port-forward}),
and kubelets
initiate connections
\emph{to} the API
server (to watch for
pods and report
status). This
bidirectionality is
why the kubelet
needs both a server
certificate and a
client certificate,
and why Chapter 8
ends by applying an
RBAC ClusterRole for
the API server's own
identity.

\hypertarget{the-trust-graph}{%
\section{The trust
graph}\label{the-trust-graph}}

\begin{outlst}
  ca.crt / ca.key  --  self-signed root, 3653 days
     |
     +-- kube-api-server.crt   CN=kubernetes
     |     SANs: 127.0.0.1, 10.32.0.1, kubernetes,
     |           kubernetes.default, kubernetes.default.svc,
     |           kubernetes.default.svc.cluster,
     |           kubernetes.svc.cluster.local,
     |           server.kubernetes.local, api-server.kubernetes.local
     |
     +-- admin.crt             CN=admin
     |                         O=system:masters
     +-- kube-controller-manager.crt
     |                         CN=O=system:kube-controller-manager
     +-- kube-scheduler.crt    CN=system:kube-scheduler
     |                         O=system:system:kube-scheduler
     +-- kube-proxy.crt        CN=system:kube-proxy
     |                         O=system:node-proxier
     +-- node-0.crt            CN=system:node:node-0
     |                         O=system:nodes
     +-- node-1.crt            CN=system:node:node-1
     |                         O=system:nodes
     |
     +-- service-accounts.crt/.key   signing key pair, no client identity
\end{outlst}

One CA signs
everything. There is
no intermediate, no
separate etcd CA, no
separate front-proxy
CA. Production
clusters typically
use three or more
CAs to limit blast
radius; this build
uses one for
clarity.

\hypertarget{check-your-understanding}{%
\section{Check your
understanding}\label{check-your-understanding}}

\textbf{Q1. The
scheduler decides
which node a pod
runs on. Does the
scheduler contact
the kubelet on that
node?}

No.~The scheduler's
only action is to
write a binding to
the API server,
setting
\texttt{pod.}\allowbreak{}\texttt{spec.}\allowbreak{}\texttt{nodeName}.
The kubelet on that
node is
independently
watching the API
server for pods
assigned to itself,
sees the new
assignment, and
acts. No component
in Kubernetes pushes
work to another
component;
everything
reconciles against
API server state.
This is why a
cluster with a dead
scheduler keeps
running existing
pods perfectly well.

\textbf{Q2. Why does
the API server's
certificate list
\texttt{10.32.0.1}
as a Subject
Alternative Name,
given that the API
server process binds
to
\texttt{0.0.0.0:6443}
on the host
network?}

Because pods reach
the API server
through the
in-cluster
\texttt{kubernetes}
Service, whose
ClusterIP is the
first address of the
service range ---
\texttt{10.32.0.1}.
A pod connecting to
\texttt{https:}\allowbreak{}\texttt{/}\allowbreak{}\texttt{/}\allowbreak{}\texttt{10.}\allowbreak{}\texttt{32.}\allowbreak{}\texttt{0.}\allowbreak{}\texttt{1}
performs TLS
verification against
that IP, so the IP
must be in the
certificate's SANs
or the handshake
fails with a
hostname-verification
error.

\textbf{Q3. If you
gave the
\texttt{kube-scheduler}
certificate
\texttt{O=system:}\allowbreak{}\texttt{masters}
by accident, would
the cluster work?}

Yes --- and that is
the danger. It would
work perfectly,
while silently
granting the
scheduler
unrestricted
cluster-admin.
Certificate subjects
are credentials, and
over-broad ones fail
open, not closed.
This is worth
internalizing before
Chapter 4.

\hypertarget{chap-1}{%
\chapter{Prerequisites}\label{chap-1}}

\upstreamlab{docs/01-prerequisites.md}

\hypertarget{what-youre-building-1}{%
\section{What you're
building}\label{what-youre-building-1}}

Nothing yet. You are
confirming that four
machines exist, run
the right OS, and
meet minimum specs.
This is the only
chapter with no
cluster artifacts.

\hypertarget{why-its-necessary-1}{%
\section{Why it's
necessary}\label{why-its-necessary-1}}

\begin{whybox}

The build depends on
Debian 12 specifics
in three places you
will hit later:
\texttt{dpkg\ -\/-print-}\allowbreak{}\texttt{architecture}
is used to select
the correct binary
manifest (Chapter
2), \texttt{apt-get}
package names for
\texttt{socat\ conntrack\ ipset\ kmod}
are Debian-specific
(Chapter 9), and
\texttt{sshd\_config}
with
\texttt{PermitRootLogin}
defaulting to
disabled is a Debian
behaviour (Chapter
3). On a different
distribution these
steps silently
diverge. On a
different
architecture mix ---
say, an ARM64
jumpbox and AMD64
workers --- the
binaries you
download on the
jumpbox will not
execute on the
nodes, and you will
not find out until
\texttt{systemctl\ start\ kubelet}
fails with
\texttt{Exec\ format\ error}
in Chapter 9.

\end{whybox}

The RAM figures also
matter more than
they look.
\texttt{server} runs
four Go processes,
one of which
(\texttt{kube-apiserver})
will comfortably use
400--600 MB at idle.
Below 2 GB the API
server gets
OOM-killed
intermittently and
the failure presents
as random
\texttt{connection\ refused}
from kubectl.

\hypertarget{key-concepts-1}{%
\section{Key
concepts}\label{key-concepts-1}}

\textbf{Architecture
must be uniform.}
The download
manifests
(\texttt{downloads-}\allowbreak{}\texttt{amd64.}\allowbreak{}\texttt{txt},
\texttt{downloads-}\allowbreak{}\texttt{arm64.}\allowbreak{}\texttt{txt})
are selected once,
on the jumpbox, and
the resulting
binaries are pushed
to every machine.
Mixed architectures
require running the
build twice with two
download sets, which
this tutorial does
not do.

\textbf{Swap must be
off on workers.} Not
here --- Chapter 9
handles it --- but
it is worth knowing
now that the kubelet
config in this build
sets
\texttt{failSwapOn:}\allowbreak{}\texttt{\ false}
and
\texttt{memorySwap.}\allowbreak{}\texttt{swapBehavior:}\allowbreak{}\texttt{\ NoSwap},
which means the
kubelet will
\emph{start} with
swap present but
will not let
containers use it.
The tutorial still
has you disable
swap, because pod
memory accounting is
unreliable
otherwise.

\hypertarget{the-work}{%
\section{The
work}\label{the-work}}

Run on \textbf{each
of the four
machines}:

\begin{cmdlst}
# on: jumpbox, server, node-0, node-1
cat /etc/os-release
\end{cmdlst}

\hypertarget{verification}{%
\section{Verification}\label{verification}}

You should see
something similar to
the following
output:

\begin{outlst}
PRETTY_NAME="Debian GNU/Linux 12 (bookworm)"
NAME="Debian GNU/Linux"
VERSION_ID="12"
VERSION="12 (bookworm)"
VERSION_CODENAME=bookworm
ID=debian
\end{outlst}

\begin{beyondbox}\calloutlabel{teal}{BEYOND THE SOURCE}

two additional
checks are worth
thirty seconds each,
because both failure
modes are expensive
to discover in
Chapter 9. Run
\texttt{dpkg\ -\/-print-}\allowbreak{}\texttt{architecture}
on all four machines
and confirm they
match. Run
\texttt{ping\ -c1}
between every pair
of machines and
confirm they are
mutually reachable.

\end{beyondbox}

\hypertarget{when-it-goes-wrong}{%
\section{When it
goes
wrong}\label{when-it-goes-wrong}}

\begin{longtable}[]{@{}
  >{\raggedright\arraybackslash}p{(\columnwidth - 4\tabcolsep) * \real{0.3200}}
  >{\raggedright\arraybackslash}p{(\columnwidth - 4\tabcolsep) * \real{0.2400}}
  >{\raggedright\arraybackslash}p{(\columnwidth - 4\tabcolsep) * \real{0.4400}}@{}}
\toprule\noalign{}
\begin{minipage}[b]{\linewidth}\raggedright
Symptom
\end{minipage} &
\begin{minipage}[b]{\linewidth}\raggedright
Cause
\end{minipage} &
\begin{minipage}[b]{\linewidth}\raggedright
Fix
\end{minipage} \\
\midrule\noalign{}
\endhead
\bottomrule\noalign{}
\endlastfoot
\texttt{VERSION\_ID="11"}
or \texttt{"13"} &
Wrong Debian release
& Rebuild the VM.
Package names and
\texttt{sshd\_config}
defaults differ
across releases. \\
\texttt{ID=ubuntu} &
Ubuntu, not Debian &
Mostly compatible,
but
\texttt{PermitRootLogin}
handling and default
package sets differ.
Expect to adapt
Chapter 3. \\
Mixed output from
\texttt{dpkg\ -\/-print-}\allowbreak{}\texttt{architecture}
& Mixed ARM64/AMD64
fleet & Rebuild so
all four match.
There is no graceful
fix later. \\
One machine
unreachable from
another & VLAN,
firewall, or
host-only networking
& Resolve now.
Chapters 7--11
assume unrestricted
connectivity between
all four. \\
\end{longtable}

\hypertarget{check-your-understanding-1}{%
\section{Check your
understanding}\label{check-your-understanding-1}}

\textbf{Q1. Why does
the jumpbox need
only 512 MB when the
others need 2 GB?}

The jumpbox runs no
cluster components.
It is an
administrative
workstation:
\texttt{openssl},
\texttt{kubectl},
\texttt{ssh},
\texttt{wget}. Its
10 GB of storage is
sized for the
\textasciitilde500
MB of downloaded
binaries plus
extraction headroom,
not for running
anything.

\textbf{Q2. The
cluster has three
members but the
tutorial requires
four machines. Why
isn't the jumpbox
part of the
cluster?}

Separating the
administrative host
from cluster members
makes the trust
boundary visible.
The jumpbox holds
\texttt{ca.key} ---
the private key that
can mint a
cluster-admin
credential for
anyone. Keeping it
off a machine that
runs workloads is
the correct
instinct, even in a
lab.

\hypertarget{chap-2}{%
\chapter{Setting Up
the
Jumpbox}\label{chap-2}}

\upstreamlab{docs/02-jumpbox.md}

\hypertarget{what-youre-building-2}{%
\section{What you're
building}\label{what-youre-building-2}}

The administrative
workstation:
command-line tools,
a clone of the
tutorial repository
(which supplies
every config
template and systemd
unit you will use),
and a local cache of
all cluster binaries
sorted into
\texttt{client/},
\texttt{controller/},
\texttt{worker/},
and
\texttt{cni-plugins/}
directories.

\hypertarget{why-its-necessary-2}{%
\section{Why it's
necessary}\label{why-its-necessary-2}}

\begin{whybox}

Downloading
\textasciitilde500
MB once and
distributing over
the LAN, rather than
pulling on each of
three machines, is
the stated reason.
The structural
reason matters more:
the jumpbox becomes
the \textbf{single
source of truth} for
the entire build.
Certificates,
kubeconfigs, machine
inventory, and
binaries all
originate here and
are pushed outward.
Every later
chapter's
\texttt{scp} command
assumes this
directory layout
exists, with these
exact paths. If you
clone the repo
somewhere else or
skip the sort into
\texttt{downloads/}\allowbreak{}\texttt{\{client,cni-}\allowbreak{}\texttt{plugins,controller,worker\}},
roughly a dozen
later commands break
with
\texttt{No\ such\ file\ or\ directory}.

\end{whybox}

\hypertarget{key-concepts-2}{%
\section{Key
concepts}\label{key-concepts-2}}

\textbf{1.
Everything runs as
root.} The tutorial
is explicit that
this is for
convenience. It
removes
\texttt{sudo} from
every command and
lets
\texttt{ssh\ root@host}
work uniformly. It
is also why
\texttt{ca.key} will
live at
\texttt{/}\allowbreak{}\texttt{root/}\allowbreak{}\texttt{kubernetes-}\allowbreak{}\texttt{the-}\allowbreak{}\texttt{hard-}\allowbreak{}\texttt{way/}\allowbreak{}\texttt{ca.}\allowbreak{}\texttt{key}
with no additional
protection ---
acceptable in a
disposable lab,
unacceptable
anywhere else.

\textbf{2. The
binaries are release
artifacts, not
packages.} There is
no
\texttt{apt\ install\ kubernetes}.
Each component is a
statically linked Go
binary fetched from
\texttt{dl.k8s.io}
or a GitHub release,
made executable, and
dropped into
\texttt{/usr/local/bin}.
This is why the
cluster has no
package manager, no
automatic upgrades,
and no service
dependencies beyond
what you write into
the systemd units
yourself.

\textbf{3. Version
pins are pre-release
in places.} The
manifest pins
\texttt{containerd-}\allowbreak{}\texttt{2.}\allowbreak{}\texttt{1.}\allowbreak{}\texttt{0-}\allowbreak{}\texttt{beta.}\allowbreak{}\texttt{0},
\texttt{etcd-}\allowbreak{}\texttt{v3.}\allowbreak{}\texttt{6.}\allowbreak{}\texttt{0-}\allowbreak{}\texttt{rc.}\allowbreak{}\texttt{3},
and
\texttt{runc\ v1.}\allowbreak{}\texttt{3.}\allowbreak{}\texttt{0-}\allowbreak{}\texttt{rc.}\allowbreak{}\texttt{1}.
These are beta and
release-candidate
builds. That is
upstream's choice
and this guide
reproduces it
exactly, but be
aware you are not
building on GA
components for the
container runtime or
the datastore.

\hypertarget{the-work-1}{%
\section{The
work}\label{the-work-1}}

Log in to the
\texttt{jumpbox}:

\begin{cmdlst}
ssh root@jumpbox
\end{cmdlst}

All commands will be
run as the
\texttt{root} user.

Install the command
line utilities:

\begin{cmdlst}
# on: jumpbox
{
  apt-get update
  apt-get -y install wget curl vim openssl git
}
\end{cmdlst}

The brace block is a
shell grouping ---
it runs the commands
as a unit and is
used throughout the
tutorial so a
copy-paste of the
whole block behaves
as one step.

Clone the
repository:

\begin{cmdlst}
# on: jumpbox
git clone --depth 1 \
  https://github.com/kelseyhightower/kubernetes-the-hard-way.git
\end{cmdlst}

\texttt{-\/-depth\ 1}
fetches only the
latest commit ---
you need the files,
not the history.

Change into the
directory.
\textbf{This is the
working directory
for the rest of the
tutorial.}

\begin{cmdlst}
# on: jumpbox
cd kubernetes-the-hard-way
\end{cmdlst}

If you ever get
lost, verify where
you are:

\begin{cmdlst}
# on: jumpbox
pwd
\end{cmdlst}

\begin{outlst}
/root/kubernetes-the-hard-way
\end{outlst}

Review the download
manifest for your
architecture:

\begin{cmdlst}
# on: jumpbox
cat downloads-$(dpkg --print-architecture).txt
\end{cmdlst}

Download the
binaries into a
\texttt{downloads}
directory:

\begin{cmdlst}
# on: jumpbox
wget -q --show-progress \
  --https-only \
  --timestamping \
  -P downloads \
  -i downloads-$(dpkg --print-architecture).txt
\end{cmdlst}

\texttt{-\/-https-only}
refuses to follow a
redirect to plain
HTTP --- a
meaningful
protection when you
are downloading
binaries that will
run as root on every
machine.
\texttt{-\/-timestamping}
makes the command
resumable:
re-running it skips
files already
present and current,
which matters
because this
download takes a
while and may be
interrupted.

List what arrived:

\begin{cmdlst}
# on: jumpbox
ls -oh downloads
\end{cmdlst}

Extract and organize
the component
binaries:

\begin{cmdlst}
# on: jumpbox
{
  ARCH=$(dpkg --print-architecture)
  mkdir -p downloads/{client,cni-plugins,controller,worker}
  tar -xvf downloads/crictl-v1.32.0-linux-${ARCH}.tar.gz \
    -C downloads/worker/
  tar -xvf downloads/containerd-2.1.0-beta.0-linux-${ARCH}.tar.gz \
    --strip-components 1 \
    -C downloads/worker/
  tar -xvf downloads/cni-plugins-linux-${ARCH}-v1.6.2.tgz \
    -C downloads/cni-plugins/
  tar -xvf downloads/etcd-v3.6.0-rc.3-linux-${ARCH}.tar.gz \
    -C downloads/ \
    --strip-components 1 \
    etcd-v3.6.0-rc.3-linux-${ARCH}/etcdctl \
    etcd-v3.6.0-rc.3-linux-${ARCH}/etcd
  mv downloads/{etcdctl,kubectl} downloads/client/
  mv downloads/{etcd,kube-apiserver,kube-controller-manager,kube-scheduler} \
    downloads/controller/
  mv downloads/{kubelet,kube-proxy} downloads/worker/
  mv downloads/runc.${ARCH} downloads/worker/runc
}
\end{cmdlst}

Two details worth
understanding.
\texttt{-\/-strip-}\allowbreak{}\texttt{components\ 1}
on the containerd
archive discards its
leading
\texttt{bin/}
directory so the
binaries land
directly in
\texttt{downloads/}\allowbreak{}\texttt{worker/}\allowbreak{}.
On the etcd archive,
the two trailing
arguments select
only
\texttt{etcdctl} and
\texttt{etcd} from
an archive that
contains
documentation and
other files you do
not need. The final
\texttt{mv} renames
\texttt{runc.amd64}
(or
\texttt{runc.arm64})
to plain
\texttt{runc},
because the systemd
unit and containerd
both expect that
name.

The sort is not
cosmetic. Chapter 9
runs
\texttt{scp\ downloads/}\allowbreak{}\texttt{worker/}\allowbreak{}\texttt{*\ root@\$\{HOST\}:}\allowbreak{}\texttt{\textasciitilde{}/}\allowbreak{}
--- it relies on
\texttt{worker/}
containing exactly
the files a worker
needs and nothing
else.

Remove the archives:

\begin{cmdlst}
# on: jumpbox
rm -rf downloads/*gz
\end{cmdlst}

Make the binaries
executable:

\begin{cmdlst}
# on: jumpbox
{
  chmod +x downloads/{client,cni-plugins,controller,worker}/*
}
\end{cmdlst}

Install
\texttt{kubectl} on
the jumpbox:

\begin{cmdlst}
# on: jumpbox
{
  cp downloads/client/kubectl /usr/local/bin/
}
\end{cmdlst}

Note this is
\texttt{cp}, not
\texttt{mv} --- the
copy in
\texttt{downloads/}\allowbreak{}\texttt{client/}\allowbreak{}
is still needed,
because Chapters 8
and 9 push it to
\texttt{server} and
both workers.

\hypertarget{verification-1}{%
\section{Verification}\label{verification-1}}

\begin{cmdlst}
# on: jumpbox
kubectl version --client
\end{cmdlst}

\begin{outlst}
Client Version: v1.32.3
Kustomize Version: v5.5.0
\end{outlst}

Also confirm the
sort worked:

\begin{beyondbox}\calloutlabel{teal}{BEYOND THE SOURCE}

\texttt{ls\ downloads/*/}
should show
\texttt{client/}
with
\texttt{etcdctl} and
\texttt{kubectl};
\texttt{controller/}
with \texttt{etcd},
\texttt{kube-apiserver},
\texttt{kube-}\allowbreak{}\texttt{controller-}\allowbreak{}\texttt{manager},
\texttt{kube-scheduler};
\texttt{worker/}
with
\texttt{containerd},
\texttt{containerd-}\allowbreak{}\texttt{shim-}\allowbreak{}\texttt{runc-}\allowbreak{}\texttt{v2},
\texttt{containerd-}\allowbreak{}\texttt{stress},
\texttt{crictl},
\texttt{kube-proxy},
\texttt{kubelet},
\texttt{runc}; and
\texttt{cni-plugins/}
with roughly a dozen
plugin binaries
including
\texttt{bridge},
\texttt{host-local},
and
\texttt{loopback}.

\end{beyondbox}

\hypertarget{when-it-goes-wrong-1}{%
\section{When it
goes
wrong}\label{when-it-goes-wrong-1}}

\begin{longtable}[]{@{}
  >{\raggedright\arraybackslash}p{(\columnwidth - 4\tabcolsep) * \real{0.3200}}
  >{\raggedright\arraybackslash}p{(\columnwidth - 4\tabcolsep) * \real{0.2400}}
  >{\raggedright\arraybackslash}p{(\columnwidth - 4\tabcolsep) * \real{0.4400}}@{}}
\toprule\noalign{}
\begin{minipage}[b]{\linewidth}\raggedright
Symptom
\end{minipage} &
\begin{minipage}[b]{\linewidth}\raggedright
Cause
\end{minipage} &
\begin{minipage}[b]{\linewidth}\raggedright
Diagnosis
\end{minipage} \\
\midrule\noalign{}
\endhead
\bottomrule\noalign{}
\endlastfoot
\texttt{tar:}\allowbreak{}\texttt{\ downloads/}\allowbreak{}\texttt{containerd-}\allowbreak{}\texttt{2.}\allowbreak{}\texttt{1.}\allowbreak{}\texttt{0-}\allowbreak{}\texttt{beta.}\allowbreak{}\texttt{0-}\allowbreak{}\texttt{linux-}\allowbreak{}\texttt{amd64.}\allowbreak{}\texttt{tar.}\allowbreak{}\texttt{gz:}\allowbreak{}\texttt{\ Cannot\ open}
& The download was
incomplete or the
architecture suffix
does not match &
Re-run the
\texttt{wget}
command;
\texttt{-\/-timestamping}
makes this safe.
Check
\texttt{dpkg\ -\/-print-}\allowbreak{}\texttt{architecture}
matches the
filename. \\
\texttt{mv:}\allowbreak{}\texttt{\ cannot\ stat\ \textquotesingle{}downloads/}\allowbreak{}\texttt{runc.}\allowbreak{}\texttt{amd64\textquotesingle{}}
& \texttt{runc}
failed to download
--- it is a bare
binary, not an
archive, and is easy
to miss in a partial
transfer &
\texttt{ls\ downloads/}\allowbreak{}\texttt{\ \textbackslash{}\textbar{}\ grep\ runc},
then re-run
\texttt{wget}. \\
\texttt{kubectl:}\allowbreak{}\texttt{\ command\ not\ found}
after install &
\texttt{/usr/local/bin}
not in PATH for this
shell &
\texttt{echo\ \$PATH};
re-login or use the
absolute path. \\
\texttt{wget} fails
on
\texttt{dl.k8s.io} &
Egress filtering or
DNS &
\texttt{curl\ -}\allowbreak{}\texttt{I\ https:}\allowbreak{}\texttt{/}\allowbreak{}\texttt{/}\allowbreak{}\texttt{dl.}\allowbreak{}\texttt{k8s.}\allowbreak{}\texttt{io/}\allowbreak{}
to isolate DNS from
TLS from policy. \\
Download extremely
slow & The 500 MB
figure is real & Let
it run;
\texttt{-\/-timestamping}
means an interrupted
download resumes
cleanly. \\
\end{longtable}

\hypertarget{check-your-understanding-2}{%
\section{Check your
understanding}\label{check-your-understanding-2}}

\textbf{Q1. Why does
the tutorial copy
\texttt{kubectl} to
\texttt{/usr/local/bin}
rather than move
it?}

Because
\texttt{downloads/}\allowbreak{}\texttt{client/}\allowbreak{}\texttt{kubectl}
is a distribution
source, not just a
local install.
Chapter 8 copies it
to \texttt{server}
and Chapter 9 copies
it to both workers.
Moving it here would
break those later
\texttt{scp}
commands.

\textbf{Q2.
\texttt{crictl} is
installed on workers
but never used in
any lab. Why is it
there?}

\texttt{crictl} is
the CRI-level
debugging client ---
it talks directly to
containerd's CRI
socket, bypassing
Kubernetes entirely.
It is included
because when a pod
is stuck in
\texttt{ContainerCreating}
and the kubelet logs
are ambiguous,
\texttt{crictl\ ps\ -a}
and
\texttt{crictl\ logs}
on the node are how
you find out whether
the container
runtime ever
received the
request. It is
diagnostic
equipment, staged in
advance.

\textbf{Q3. What
would break if you
cloned the
repository to
\texttt{/opt/kthw}
instead?}

Nothing immediately
--- but every later
chapter's relative
paths
(\texttt{configs/},
\texttt{units/},
\texttt{downloads/},
\texttt{ca.conf},
\texttt{machines.txt})
assume you are
\emph{in} the
repository directory
when you run them.
As long as you
\texttt{cd} there
first, the location
is arbitrary. The
failure mode is
subtle: running a
\texttt{scp\ configs/...}
command from the
wrong directory
gives
\texttt{No\ such\ file\ or\ directory}
on a file you know
exists.

\hypertarget{chap-3}{%
\chapter{Provisioning
Compute
Resources}\label{chap-3}}

\upstreamlab{docs/03-compute-resources.md}

\hypertarget{what-youre-building-3}{%
\section{What you're
building}\label{what-youre-building-3}}

A machine inventory
file, key-based root
SSH from the jumpbox
to all three cluster
machines, correct
hostnames on each
machine, and a
shared
\texttt{/etc/hosts}
entry set so that
all four machines
can resolve each
other by name.

\hypertarget{why-its-necessary-3}{%
\section{Why it's
necessary}\label{why-its-necessary-3}}

\begin{whybox}

Kubernetes
identifies nodes by
\textbf{hostname},
not IP address. The
kubelet registers
itself under
\texttt{hostname},
the API server
certificate has
\texttt{server.}\allowbreak{}\texttt{kubernetes.}\allowbreak{}\texttt{local}
in its SANs, and
every kubeconfig
points at
\texttt{https:}\allowbreak{}\texttt{/}\allowbreak{}\texttt{/}\allowbreak{}\texttt{server.}\allowbreak{}\texttt{kubernetes.}\allowbreak{}\texttt{local:}\allowbreak{}\texttt{6443}.
If name resolution
is wrong or
inconsistent, the
failures land far
from the cause:

\end{whybox}

\begin{itemize}
\tightlist
\item
  Hostname not set
  on a worker
  \textrightarrow{}
  the kubelet
  registers as
  \texttt{debian} or
  \texttt{localhost},
  its client
  certificate says
  \texttt{CN=system:}\allowbreak{}\texttt{node:}\allowbreak{}\texttt{node-}\allowbreak{}\texttt{0},
  and the Node
  authorizer rejects
  its requests with
  \texttt{403\ Forbidden}
  because the
  username does not
  match the node it
  is trying to act
  on.
\item
  \texttt{server.}\allowbreak{}\texttt{kubernetes.}\allowbreak{}\texttt{local}
  unresolvable on a
  worker
  \textrightarrow{}
  the kubelet cannot
  reach the API
  server at all,
  logs
  \texttt{dial\ tcp:}\allowbreak{}\texttt{\ lookup\ server.}\allowbreak{}\texttt{kubernetes.}\allowbreak{}\texttt{local:}\allowbreak{}\texttt{\ no\ such\ host},
  and the node never
  appears in
  \texttt{kubectl\ get\ nodes}.
\item
  \texttt{server.}\allowbreak{}\texttt{kubernetes.}\allowbreak{}\texttt{local}
  resolvable but not
  in the API server
  certificate SANs
  \textrightarrow{}
  TLS handshake
  fails with
  \texttt{certificate\ is\ valid\ for\ .}\allowbreak{}\texttt{.}\allowbreak{}\texttt{.}\allowbreak{}\texttt{,\ not\ server.}\allowbreak{}\texttt{kubernetes.}\allowbreak{}\texttt{local}.
\end{itemize}

There is no DNS
server in this
build.
\texttt{/etc/hosts}
is the entire name
resolution layer.

\hypertarget{key-concepts-3}{%
\section{Key
concepts}\label{key-concepts-3}}

\textbf{1.
\texttt{machines.txt}
is the cluster's
inventory database.}
Four
whitespace-separated
columns:

\begin{outlst}
IPV4_ADDRESS FQDN HOSTNAME POD_SUBNET
\end{outlst}

It is read by
\texttt{while\ read\ IP\ FQDN\ HOST\ SUBNET}
loops in this
chapter, and again
in Chapter 9 (to
template the
per-node CNI subnet)
and Chapter 11 (to
build the routing
table). Get a column
wrong here and the
error surfaces six
chapters later as a
pod networking
failure.

\textbf{2.
\texttt{server} has
no POD\_SUBNET.}
Look closely at the
example:

\begin{outlst}
XXX.XXX.XXX.XXX server.kubernetes.local server
XXX.XXX.XXX.XXX node-0.kubernetes.local node-0 10.200.0.0/24
XXX.XXX.XXX.XXX node-1.kubernetes.local node-1 10.200.1.0/24
\end{outlst}

The control plane
host runs no pods,
so it gets no pod
subnet --- the
fourth field is
simply absent on
that line. The
\texttt{while\ read}
loops handle this
because
\texttt{SUBNET}
becomes an empty
string. This is
intentional, not an
omission in the
example.

\textbf{3. Root SSH
is a deliberate,
acknowledged
tradeoff.} The
tutorial says so
plainly:
\emph{``Security is
a tradeoff, and in
this case, we are
optimizing for
convenience.''}
Enabling
\texttt{PermitRootLogin\ yes}
on an
internet-reachable
machine with
password auth is how
machines get taken.
In a lab on a
private network,
with key-based auth,
the risk is bounded.
Do not carry this
pattern anywhere
else.

\textbf{4.
\texttt{127.0.1.1}
is a Debian
convention.} The
hostname-setting
loop rewrites the
\texttt{127.0.1.1}
line rather than
\texttt{127.0.0.1}.
Debian uses
\texttt{127.0.1.1}
for the machine's
own FQDN
specifically so that
\texttt{hostname\ -\/-fqdn}
resolves even with
no network. Editing
\texttt{127.0.0.1}
instead would break
loopback in
confusing ways.

\hypertarget{the-work-2}{%
\section{The
work}\label{the-work-2}}

Create
\texttt{machines.txt}
on the jumpbox with
your actual IP
addresses:

\begin{cmdlst}
# on: jumpbox
cat machines.txt
\end{cmdlst}

\begin{outlst}
XXX.XXX.XXX.XXX server.kubernetes.local server
XXX.XXX.XXX.XXX node-0.kubernetes.local node-0 10.200.0.0/24
XXX.XXX.XXX.XXX node-1.kubernetes.local node-1 10.200.1.0/24
\end{outlst}

Use single spaces as
separators and no
trailing whitespace
--- the
\texttt{cut\ -d\ "\ "\ -f\ 4}
commands in Chapter
11 are literal about
this.

\textbf{Enable root
SSH access.} Skip
this section if root
SSH already works.
On \textbf{each of
\texttt{server},
\texttt{node-0},
\texttt{node-1}},
log in with your
normal account and
switch to root:

\begin{cmdlst}
# on: server, node-0, node-1
su - root
\end{cmdlst}

\begin{cmdlst}
# on: server, node-0, node-1
sed -i \
  's/^#*PermitRootLogin.*/PermitRootLogin yes/' \
  /etc/ssh/sshd_config
\end{cmdlst}

The \texttt{\^{}\#*}
matches the line
whether it is
commented out or
already present.

\begin{cmdlst}
# on: server, node-0, node-1
systemctl restart sshd
\end{cmdlst}

\textbf{Generate and
distribute SSH
keys.} Back on the
jumpbox:

\begin{cmdlst}
# on: jumpbox
ssh-keygen
\end{cmdlst}

\begin{outlst}
Generating public/private rsa key pair.
Enter file in which to save the key (/root/.ssh/id_rsa):
Enter passphrase (empty for no passphrase):
Enter same passphrase again:
Your identification has been saved in /root/.ssh/id_rsa
Your public key has been saved in /root/.ssh/id_rsa.pub
\end{outlst}

Accept the defaults.
An empty passphrase
is what makes the
unattended
\texttt{while\ read}
loops in later
chapters work
without an agent.

\begin{cmdlst}
# on: jumpbox
while read IP FQDN HOST SUBNET; do
  ssh-copy-id root@${IP}
done < machines.txt
\end{cmdlst}

You will be prompted
for the root
password of each
machine once.

Verify key-based
access:

\begin{cmdlst}
# on: jumpbox
while read IP FQDN HOST SUBNET; do
  ssh -n root@${IP} hostname
done < machines.txt
\end{cmdlst}

\begin{outlst}
server
node-0
node-1
\end{outlst}

The \texttt{-n} flag
redirects stdin from
\texttt{/dev/null}.
Without it,
\texttt{ssh} inside
a
\texttt{while\ read}
loop consumes the
remaining lines of
\texttt{machines.txt}
and the loop runs
exactly once. This
is the single most
common shell bug in
this entire build,
and it appears in
nearly every loop
from here on.

\textbf{Set
hostnames.} From the
jumpbox:

\begin{cmdlst}
# on: jumpbox
while read IP FQDN HOST SUBNET; do
    CMD="sed -i 's/^127.0.1.1.*/127.0.1.1\t${FQDN} ${HOST}/' /etc/hosts"
    ssh -n root@${IP} "$CMD"
    ssh -n root@${IP} hostnamectl set-hostname ${HOST}
    ssh -n root@${IP} systemctl restart systemd-hostnamed
done < machines.txt
\end{cmdlst}

Three operations per
machine: rewrite the
FQDN line in
\texttt{/etc/hosts},
set the transient
and static hostname
via
\texttt{hostnamectl},
and restart
\texttt{systemd-}\allowbreak{}\texttt{hostnamed}
so the change is
picked up
immediately.

Verify:

\begin{cmdlst}
# on: jumpbox
while read IP FQDN HOST SUBNET; do
  ssh -n root@${IP} hostname --fqdn
done < machines.txt
\end{cmdlst}

\begin{outlst}
server.kubernetes.local
node-0.kubernetes.local
node-1.kubernetes.local
\end{outlst}

\texttt{hostname\ -\/-fqdn}
(not plain
\texttt{hostname})
is the check that
matters --- it
proves the
\texttt{/etc/hosts}
edit and the
\texttt{hostnamectl}
call agree.

\textbf{Build the
shared hosts file:}

\begin{cmdlst}
# on: jumpbox
echo "" > hosts
echo "# Kubernetes The Hard Way" >> hosts
\end{cmdlst}

\begin{cmdlst}
# on: jumpbox
while read IP FQDN HOST SUBNET; do
    ENTRY="${IP} ${FQDN} ${HOST}"
    echo $ENTRY >> hosts
done < machines.txt
\end{cmdlst}

\begin{cmdlst}
# on: jumpbox
cat hosts
\end{cmdlst}

\begin{outlst}

# Kubernetes The Hard Way
XXX.XXX.XXX.XXX server.kubernetes.local server
XXX.XXX.XXX.XXX node-0.kubernetes.local node-0
XXX.XXX.XXX.XXX node-1.kubernetes.local node-1
\end{outlst}

\textbf{Append to
the jumpbox's own
\texttt{/etc/hosts}:}

\begin{cmdlst}
# on: jumpbox
cat hosts >> /etc/hosts
\end{cmdlst}

\begin{cmdlst}
# on: jumpbox
cat /etc/hosts
\end{cmdlst}

\begin{outlst}
127.0.0.1       localhost
127.0.1.1       jumpbox

# The following lines are desirable for IPv6 capable hosts
::1     localhost ip6-localhost ip6-loopback
ff02::1 ip6-allnodes
ff02::2 ip6-allrouters

# Kubernetes The Hard Way
XXX.XXX.XXX.XXX server.kubernetes.local server
XXX.XXX.XXX.XXX node-0.kubernetes.local node-0
XXX.XXX.XXX.XXX node-1.kubernetes.local node-1
\end{outlst}

Note this is
\texttt{\textgreater{}\textgreater{}},
appending. Running
it twice produces
duplicate entries
--- harmless for
resolution, but
confusing when you
\texttt{cat\ /etc/hosts}
later to diagnose
something.

Confirm
hostname-based SSH
now works:

\begin{cmdlst}
# on: jumpbox
for host in server node-0 node-1
   do ssh root@${host} hostname
done
\end{cmdlst}

\begin{outlst}
server
node-0
node-1
\end{outlst}

\textbf{Distribute
the hosts file to
the cluster
machines:}

\begin{cmdlst}
# on: jumpbox
while read IP FQDN HOST SUBNET; do
  scp hosts root@${HOST}:~/
  ssh -n \
    root@${HOST} "cat hosts >> /etc/hosts"
done < machines.txt
\end{cmdlst}

This loop uses
\texttt{\$\{HOST\}}
rather than
\texttt{\$\{IP\}}
--- it works only
because the previous
step made hostnames
resolvable on the
jumpbox.

\hypertarget{verification-2}{%
\section{Verification}\label{verification-2}}

Every machine should
resolve every other
machine by short
name and FQDN. From
the jumpbox:

\begin{beyondbox}\calloutlabel{teal}{BEYOND THE SOURCE}

\texttt{for\ h\ in\ server\ node-}\allowbreak{}\texttt{0\ node-}\allowbreak{}\texttt{1;\ do\ echo\ "-\/-\/-\ \$h";\ ssh\ -}\allowbreak{}\texttt{n\ root@\$h\ "getent\ hosts\ server\ node-}\allowbreak{}\texttt{0\ node-}\allowbreak{}\texttt{1";\ done}
--- three lines of
output per machine
confirms the full
mesh before you
depend on it in
Chapter 7.

\end{beyondbox}

\hypertarget{when-it-goes-wrong-2}{%
\section{When it
goes
wrong}\label{when-it-goes-wrong-2}}

\begin{longtable}[]{@{}
  >{\raggedright\arraybackslash}p{(\columnwidth - 4\tabcolsep) * \real{0.3200}}
  >{\raggedright\arraybackslash}p{(\columnwidth - 4\tabcolsep) * \real{0.2400}}
  >{\raggedright\arraybackslash}p{(\columnwidth - 4\tabcolsep) * \real{0.4400}}@{}}
\toprule\noalign{}
\begin{minipage}[b]{\linewidth}\raggedright
Symptom
\end{minipage} &
\begin{minipage}[b]{\linewidth}\raggedright
Cause
\end{minipage} &
\begin{minipage}[b]{\linewidth}\raggedright
Diagnosis
\end{minipage} \\
\midrule\noalign{}
\endhead
\bottomrule\noalign{}
\endlastfoot
A
\texttt{while\ read}
loop runs once and
stops & Missing
\texttt{-n} on
\texttt{ssh}; the
SSH session ate
stdin & Add
\texttt{-n}. This is
the classic bug ---
check it first in
any loop that
terminates early. \\
\texttt{Permission\ denied\ (publickey,password)}
&
\texttt{PermitRootLogin}
not applied, or
\texttt{sshd} not
restarted & On the
target:
\texttt{sshd\ -}\allowbreak{}\texttt{T\ \textbackslash{}\textbar{}\ grep\ permitrootlogin}
shows the effective
value. \\
\texttt{hostname\ -\/-fqdn}
returns the short
name only & The
\texttt{/etc/hosts}
\texttt{127.0.1.1}
line was not
rewritten &
\texttt{grep\ 127.}\allowbreak{}\texttt{0.}\allowbreak{}\texttt{1.}\allowbreak{}\texttt{1\ /}\allowbreak{}\texttt{etc/}\allowbreak{}\texttt{hosts}
on the target. The
\texttt{sed} may
have failed if the
line is absent
entirely. \\
\texttt{hostname\ -\/-fqdn}:
\texttt{Name\ or\ service\ not\ known}
& No
\texttt{127.0.1.1}
line at all & Add it
manually:
\texttt{127.}\allowbreak{}\texttt{0.}\allowbreak{}\texttt{1.}\allowbreak{}\texttt{1\textless{}TAB\textgreater{}node-}\allowbreak{}\texttt{0.}\allowbreak{}\texttt{kubernetes.}\allowbreak{}\texttt{local\ node-}\allowbreak{}\texttt{0}. \\
Duplicate blocks in
\texttt{/etc/hosts}
& The append step
ran twice &
Harmless, but edit
them out to keep
later diagnosis
readable. \\
\texttt{cut\ -d\ "\ "\ -f\ 4}
returns empty in
Chapter 11 & Tabs or
multiple spaces in
\texttt{machines.txt}
& Rewrite the file
with exactly one
space between
fields. \\
\end{longtable}

\hypertarget{check-your-understanding-3}{%
\section{Check your
understanding}\label{check-your-understanding-3}}

\textbf{Q1. The
tutorial sets both a
hostname and an
\texttt{/etc/hosts}
entry on each
machine. Why isn't
one enough?}

They solve different
problems.
\texttt{hostnamectl\ set-}\allowbreak{}\texttt{hostname}
determines what the
machine calls
\emph{itself} ---
this is what the
kubelet uses to
register as a node.
The
\texttt{/etc/hosts}
entries determine
how machines resolve
\emph{each other}
--- this is what
lets a kubelet on
\texttt{node-0}
connect to
\texttt{https:}\allowbreak{}\texttt{/}\allowbreak{}\texttt{/}\allowbreak{}\texttt{server.}\allowbreak{}\texttt{kubernetes.}\allowbreak{}\texttt{local:}\allowbreak{}\texttt{6443}.
Setting only the
first gives you
correctly-named
nodes that cannot
find the API server;
setting only the
second gives you
reachable machines
that register under
the wrong node name.

\textbf{Q2. Why is
\texttt{server.}\allowbreak{}\texttt{kubernetes.}\allowbreak{}\texttt{local}
used in kubeconfigs
rather than the
server's IP
address?}

Three reasons. The
API server
certificate's SANs
include that name,
so TLS verification
succeeds against it.
The name is stable
across IP changes
--- re-addressing
the server means
editing
\texttt{/etc/hosts}
rather than
regenerating five
kubeconfigs. And it
makes the
kubeconfigs
identical across
nodes, which is what
allows the
\texttt{for\ host\ in\ node-}\allowbreak{}\texttt{0\ node-}\allowbreak{}\texttt{1}
loops to work
without per-host
special casing.

\textbf{Q3. What
happens in Chapter
11 if you put a tab
between fields in
\texttt{machines.txt}?}

Chapter 11 extracts
subnets with
\texttt{cut\ -d\ "\ "\ -f\ 4},
which splits on a
literal single
space. With tabs,
field 4 comes back
empty,
\texttt{ip\ route\ add\ \ via}
is a syntax error,
and pod-to-pod
cross-node traffic
silently never
works. The
\texttt{while\ read}
loops in this
chapter would still
work fine, because
\texttt{read} splits
on any whitespace
--- so the
inconsistency
survives undetected
for eight chapters.

\hypertarget{chap-4}{%
\chapter{Provisioning
a CA and Generating
TLS
Certificates}\label{chap-4}}

\upstreamlab{docs/04-certificate-authority.md}

\hypertarget{what-youre-building-4}{%
\section{What you're
building}\label{what-youre-building-4}}

A self-signed
certificate
authority, and eight
certificate/key
pairs signed by it:
\texttt{admin},
\texttt{node-0},
\texttt{node-1},
\texttt{kube-proxy},
\texttt{kube-scheduler},
\texttt{kube-}\allowbreak{}\texttt{controller-}\allowbreak{}\texttt{manager},
\texttt{kube-api-server},
and
\texttt{service-}\allowbreak{}\texttt{accounts}.
Some are then
distributed to the
machines that need
them.

\hypertarget{why-its-necessary-4}{%
\section{Why it's
necessary}\label{why-its-necessary-4}}

\begin{whybox}

Every authenticated
connection in this
cluster is mutual
TLS, and there is no
other authentication
mechanism configured
--- no tokens, no
OIDC, no bootstrap
tokens. The
certificates are not
transport encryption
with identity bolted
on; \textbf{the
certificate subject
\emph{is} the
identity}. When
\texttt{kube-scheduler}
connects to the API
server, the API
server reads
\texttt{CN=system:}\allowbreak{}\texttt{kube-}\allowbreak{}\texttt{scheduler}
out of the client
certificate and that
string becomes the
username it
authorizes against
RBAC.

\end{whybox}

This is why the
failure mode is so
distinctive. A wrong
certificate subject
does not produce a
TLS error. The
handshake succeeds,
the connection is
established, and
then every API
request returns
\texttt{403\ Forbidden}
with a message
naming a user you
did not expect. If
you see
\texttt{User\ "system:}\allowbreak{}\texttt{kube-}\allowbreak{}\texttt{scheduler"\ cannot\ list\ resource\ "pods"},
the certificate is
fine and RBAC is the
problem. If you see
\texttt{User\ ""\ cannot.}\allowbreak{}\texttt{.}\allowbreak{}\texttt{.}\allowbreak{}
or an anonymous-user
error, the
certificate never
got presented.

\hypertarget{key-concepts-4}{%
\section{Key
concepts}\label{key-concepts-4}}

\textbf{1. CN
becomes the
username; O becomes
the group.} This
mapping is hardcoded
in the Kubernetes
X.509 authenticator.
Kubernetes ships
built-in
ClusterRoleBindings
that expect specific
values:

\begin{longtable}[]{@{}
  >{\raggedright\arraybackslash}p{(\columnwidth - 6\tabcolsep) * \real{0.2500}}
  >{\raggedright\arraybackslash}p{(\columnwidth - 6\tabcolsep) * \real{0.2500}}
  >{\raggedright\arraybackslash}p{(\columnwidth - 6\tabcolsep) * \real{0.2500}}
  >{\raggedright\arraybackslash}p{(\columnwidth - 6\tabcolsep) * \real{0.2500}}@{}}
\toprule\noalign{}
\begin{minipage}[b]{\linewidth}\raggedright
Certificate
\end{minipage} &
\begin{minipage}[b]{\linewidth}\raggedright
CN
(\textrightarrow{}
username)
\end{minipage} &
\begin{minipage}[b]{\linewidth}\raggedright
O (\textrightarrow{}
group)
\end{minipage} &
\begin{minipage}[b]{\linewidth}\raggedright
What authorizes it
\end{minipage} \\
\midrule\noalign{}
\endhead
\bottomrule\noalign{}
\endlastfoot
\texttt{admin} &
\texttt{admin} &
\texttt{system:masters}
& Built-in binding
grants
\texttt{cluster-admin}
to the
\texttt{system:masters}
\textbf{group} \\
\texttt{node-0} &
\texttt{system:}\allowbreak{}\texttt{node:}\allowbreak{}\texttt{node-}\allowbreak{}\texttt{0}
&
\texttt{system:nodes}
& The \textbf{Node
authorizer}, which
requires exactly
this username
format \\
\texttt{kube-proxy}
&
\texttt{system:}\allowbreak{}\texttt{kube-}\allowbreak{}\texttt{proxy}
&
\texttt{system:}\allowbreak{}\texttt{node-}\allowbreak{}\texttt{proxier}
& Built-in
\texttt{system:}\allowbreak{}\texttt{node-}\allowbreak{}\texttt{proxier}
binding, on the
\textbf{group} \\
\texttt{kube-scheduler}
&
\texttt{system:}\allowbreak{}\texttt{kube-}\allowbreak{}\texttt{scheduler}
&
\texttt{system:}\allowbreak{}\texttt{system:}\allowbreak{}\texttt{kube-}\allowbreak{}\texttt{scheduler}
& Built-in
\texttt{system:}\allowbreak{}\texttt{kube-}\allowbreak{}\texttt{scheduler}
binding, on the
\textbf{user} \\
\texttt{kube-}\allowbreak{}\texttt{controller-}\allowbreak{}\texttt{manager}
&
\texttt{system:}\allowbreak{}\texttt{kube-}\allowbreak{}\texttt{controller-}\allowbreak{}\texttt{manager}
&
\texttt{system:}\allowbreak{}\texttt{kube-}\allowbreak{}\texttt{controller-}\allowbreak{}\texttt{manager}
& Built-in binding,
on the
\textbf{user} \\
\texttt{kube-api-server}
&
\texttt{kubernetes}
& \emph{(none)} &
Used as a server
cert, and as a
client cert to
kubelets --- see
Chapter 8's RBAC
step \\
\end{longtable}

Note the scheduler's
\texttt{O\ =\ system:}\allowbreak{}\texttt{system:}\allowbreak{}\texttt{kube-}\allowbreak{}\texttt{scheduler}.
That doubled prefix
is in the upstream
\texttt{ca.conf} as
written. It is
harmless here
because the
scheduler is
authorized by its
\textbf{username},
not its group ---
but it is a good
illustration of how
a subject field can
be wrong without
anything failing,
which is exactly
what makes
certificate subjects
dangerous.

\textbf{2. The Node
authorizer needs
both fields.} From
\texttt{ca.conf}'s
own comment:
\emph{Kubelets must
use a credential
that identifies them
as being in the
\texttt{system:nodes}
group, with a
username of
\texttt{system:}\allowbreak{}\texttt{node:}\allowbreak{}\texttt{\textless{}nodeName\textgreater{}}.}
The Node authorizer
then restricts each
kubelet to reading
only the objects
related to pods
scheduled on
\emph{its own} node.
This is a dedicated
authorization mode
precisely because
writing equivalent
RBAC rules per node
would be
unmanageable.

\textbf{3. The
\texttt{service-}\allowbreak{}\texttt{accounts}
pair is a signing
key, not a login.}
It is used by the
API server as
\texttt{-\/-service-}\allowbreak{}\texttt{account-}\allowbreak{}\texttt{key-}\allowbreak{}\texttt{file}
and
\texttt{-\/-service-}\allowbreak{}\texttt{account-}\allowbreak{}\texttt{signing-}\allowbreak{}\texttt{key-}\allowbreak{}\texttt{file},
and by the
controller manager
as
\texttt{-\/-service-}\allowbreak{}\texttt{account-}\allowbreak{}\texttt{private-}\allowbreak{}\texttt{key-}\allowbreak{}\texttt{file}.
Its job is to sign
and verify the JWTs
mounted into pods as
ServiceAccount
tokens. The
certificate around
it is essentially
ceremonial --- what
matters is that both
components have the
same key.

\textbf{4. The API
server's SANs are
the full set of
names anyone might
use to reach it.}
From
\texttt{ca.conf}:

\begin{outlst}
IP.0  = 127.0.0.1
IP.1  = 10.32.0.1
DNS.0 = kubernetes
DNS.1 = kubernetes.default
DNS.2 = kubernetes.default.svc
DNS.3 = kubernetes.default.svc.cluster
DNS.4 = kubernetes.svc.cluster.local
DNS.5 = server.kubernetes.local
DNS.6 = api-server.kubernetes.local
\end{outlst}

\texttt{127.0.0.1}
for
\texttt{admin.}\allowbreak{}\texttt{kubeconfig}
used on the server
itself.
\texttt{10.32.0.1}
for in-cluster pods
reaching the
\texttt{kubernetes}
Service. The
\texttt{kubernetes.*}
names for in-cluster
DNS resolution.
\texttt{server.}\allowbreak{}\texttt{kubernetes.}\allowbreak{}\texttt{local}
for everything
external. A name
missing from this
list is a name that
cannot be used,
ever, without
regenerating the
certificate and
restarting the API
server.

\textbf{5.
\texttt{-}\allowbreak{}\texttt{copy\_extensions\ copyall}
is load-bearing.} By
default,
\texttt{openssl\ x509\ -}\allowbreak{}\texttt{req}
\textbf{discards}
the extensions
requested in the CSR
--- including
\texttt{subjectAltName}.
Without this flag
the API server
certificate would be
issued with no SANs
at all, and every
client would fail
TLS verification. It
is one flag between
a working cluster
and a completely
broken one.

\hypertarget{the-work-3}{%
\section{The
work}\label{the-work-3}}

All commands in this
chapter run on the
\textbf{jumpbox}, in
\texttt{/}\allowbreak{}\texttt{root/}\allowbreak{}\texttt{kubernetes-}\allowbreak{}\texttt{the-}\allowbreak{}\texttt{hard-}\allowbreak{}\texttt{way}.

Review the OpenSSL
configuration:

\begin{cmdlst}
# on: jumpbox
cat ca.conf
\end{cmdlst}

You do not need to
understand every
line, but do read
the
\texttt{{[}admin{]}},
\texttt{{[}node-0{]}},
and
\texttt{{[}kube-}\allowbreak{}\texttt{api-}\allowbreak{}\texttt{server{]}}
sections against the
table above --- they
are the three that
explain the whole
scheme.

\textbf{Generate the
CA:}

\begin{cmdlst}
# on: jumpbox
{
  openssl genrsa -out ca.key 4096
  openssl req -x509 -new -sha512 -noenc \
    -key ca.key -days 3653 \
    -config ca.conf \
    -out ca.crt
}
\end{cmdlst}

Results:

\begin{outlst}
ca.crt ca.key
\end{outlst}

\texttt{-x509} makes
this self-signed
rather than a CSR.
\texttt{-noenc}
leaves the private
key unencrypted on
disk, which is what
lets the unattended
loop below use it
without a passphrase
prompt --- and is
also why
\texttt{ca.key} is
now the most
sensitive file in
the build.
\texttt{3653} days
is ten years plus
leap days;
certificate expiry
will not be what
breaks your lab.

\textbf{Generate the
eight component
certificates:}

\begin{cmdlst}
# on: jumpbox
certs=(
  "admin" "node-0" "node-1"
  "kube-proxy" "kube-scheduler"
  "kube-controller-manager"
  "kube-api-server"
  "service-accounts"
)
\end{cmdlst}

\begin{cmdlst}
# on: jumpbox
for i in ${certs[*]}; do
  openssl genrsa -out "${i}.key" 4096

  openssl req -new -key "${i}.key" -sha256 \
    -config "ca.conf" -section ${i} \
    -out "${i}.csr"

  openssl x509 -req -days 3653 -in "${i}.csr" \
    -copy_extensions copyall \
    -sha256 -CA "ca.crt" \
    -CAkey "ca.key" \
    -CAcreateserial \
    -out "${i}.crt"
done
\end{cmdlst}

Three steps per
certificate:
generate a 4096-bit
RSA key, build a CSR
using the matching
\texttt{{[}section{]}}
of \texttt{ca.conf},
then sign it with
the CA. The array
element names
\textbf{must} match
the section names in
\texttt{ca.conf} ---
\texttt{-section\ \$\{i\}}
is a direct lookup,
and a typo produces
an OpenSSL error
rather than a
silently wrong
certificate.

List the results:

\begin{cmdlst}
# on: jumpbox
ls -1 *.crt *.key *.csr
\end{cmdlst}

\textbf{Distribute
to the worker
nodes:}

\begin{cmdlst}
# on: jumpbox
for host in node-0 node-1; do
  ssh root@${host} mkdir /var/lib/kubelet/

  scp ca.crt root@${host}:/var/lib/kubelet/

  scp ${host}.crt \
    root@${host}:/var/lib/kubelet/kubelet.crt

  scp ${host}.key \
    root@${host}:/var/lib/kubelet/kubelet.key
done
\end{cmdlst}

Each node gets three
files: the CA
certificate (so the
kubelet can verify
clients that connect
to \emph{it}, per
\texttt{clientCAFile}
in
\texttt{kubelet-}\allowbreak{}\texttt{config.}\allowbreak{}\texttt{yaml}),
and its own
certificate and key
\textbf{renamed} to
\texttt{kubelet.crt}
/
\texttt{kubelet.key}.
The rename is what
lets one identical
\texttt{kubelet-}\allowbreak{}\texttt{config.}\allowbreak{}\texttt{yaml}
work on both nodes.

\textbf{Distribute
to the control
plane:}

\begin{cmdlst}
# on: jumpbox
scp \
  ca.key ca.crt \
  kube-api-server.key kube-api-server.crt \
  service-accounts.key service-accounts.crt \
  root@server:~/
\end{cmdlst}

\begin{riskbox}\calloutlabel{brick}{RISK}

\texttt{ca.key} goes
to \texttt{server}
because
\texttt{kube-}\allowbreak{}\texttt{controller-}\allowbreak{}\texttt{manager}
uses it as
\texttt{-\/-cluster-}\allowbreak{}\texttt{signing-}\allowbreak{}\texttt{key-}\allowbreak{}\texttt{file}
--- it signs
certificates for the
CSR API. That means
the control plane
host can mint any
credential in the
cluster, which is
worth registering as
a fact about
Kubernetes
generally, not just
this lab.

\end{riskbox}

\begin{notebox}\calloutlabel{steel}{FROM THE SOURCE}

The
\texttt{kube-proxy},
\texttt{kube-}\allowbreak{}\texttt{controller-}\allowbreak{}\texttt{manager},
\texttt{kube-scheduler},
and \texttt{kubelet}
client certificates
will be used to
generate client
authentication
configuration files
in the next lab.

\end{notebox}

That is, those four
certificates stay on
the jumpbox for now.
Chapter 5 embeds
them into
kubeconfigs, and it
is the kubeconfigs
--- not the raw
certificates ---
that get
distributed.

\hypertarget{verification-3}{%
\section{Verification}\label{verification-3}}

Confirm a
certificate's
subject and
extensions came out
as intended:

\begin{beyondbox}\calloutlabel{teal}{BEYOND THE SOURCE}

\texttt{openssl\ x509\ -}\allowbreak{}\texttt{in\ node-}\allowbreak{}\texttt{0.}\allowbreak{}\texttt{crt\ -}\allowbreak{}\texttt{noout\ -}\allowbreak{}\texttt{subject\ -}\allowbreak{}\texttt{ext\ subjectAltName}
should print
\texttt{subject=CN=system:}\allowbreak{}\texttt{node:}\allowbreak{}\texttt{node-}\allowbreak{}\texttt{0,\ O=system:}\allowbreak{}\texttt{nodes,\ C=US,\ ST=Washington,\ L=Seattle}
and a SAN line
containing
\texttt{DNS:}\allowbreak{}\texttt{node-}\allowbreak{}\texttt{0,\ IP\ Address:}\allowbreak{}\texttt{127.}\allowbreak{}\texttt{0.}\allowbreak{}\texttt{0.}\allowbreak{}\texttt{1}.
If the SAN line is
missing,
\texttt{-}\allowbreak{}\texttt{copy\_extensions\ copyall}
did not take effect
and every
certificate needs
regenerating. Check
\texttt{kube-}\allowbreak{}\texttt{api-}\allowbreak{}\texttt{server.}\allowbreak{}\texttt{crt}
the same way --- it
should list all nine
SANs.

\end{beyondbox}

\hypertarget{when-it-goes-wrong-3}{%
\section{When it
goes
wrong}\label{when-it-goes-wrong-3}}

\begin{longtable}[]{@{}
  >{\raggedright\arraybackslash}p{(\columnwidth - 4\tabcolsep) * \real{0.3200}}
  >{\raggedright\arraybackslash}p{(\columnwidth - 4\tabcolsep) * \real{0.2400}}
  >{\raggedright\arraybackslash}p{(\columnwidth - 4\tabcolsep) * \real{0.4400}}@{}}
\toprule\noalign{}
\begin{minipage}[b]{\linewidth}\raggedright
Symptom
\end{minipage} &
\begin{minipage}[b]{\linewidth}\raggedright
Cause
\end{minipage} &
\begin{minipage}[b]{\linewidth}\raggedright
Diagnosis
\end{minipage} \\
\midrule\noalign{}
\endhead
\bottomrule\noalign{}
\endlastfoot
\texttt{Error\ opening\ config\ file\ ca.}\allowbreak{}\texttt{conf}
& Wrong working
directory &
\texttt{pwd} should
be
\texttt{/}\allowbreak{}\texttt{root/}\allowbreak{}\texttt{kubernetes-}\allowbreak{}\texttt{the-}\allowbreak{}\texttt{hard-}\allowbreak{}\texttt{way}. \\
\texttt{Unable\ to\ get\ section\ \textless{}name\textgreater{}}
& Array element does
not match a
\texttt{ca.conf}
section & Compare
\texttt{certs=(...)}
against
\texttt{grep\ \textquotesingle{}\^{}\textbackslash{}{[}\textquotesingle{}\ ca.}\allowbreak{}\texttt{conf}. \\
Later:
\texttt{x509:}\allowbreak{}\texttt{\ certificate\ is\ valid\ for\ X,\ not\ server.}\allowbreak{}\texttt{kubernetes.}\allowbreak{}\texttt{local}
& SANs missing from
the API server cert
&
\texttt{openssl\ x509\ -}\allowbreak{}\texttt{in\ kube-}\allowbreak{}\texttt{api-}\allowbreak{}\texttt{server.}\allowbreak{}\texttt{crt\ -}\allowbreak{}\texttt{noout\ -}\allowbreak{}\texttt{ext\ subjectAltName}.
Empty means
\texttt{-}\allowbreak{}\texttt{copy\_extensions}
was lost;
regenerate. \\
Later:
\texttt{Unauthorized}
from a component &
Certificate not
presented, or signed
by a different CA &
On \texttt{server}:
\texttt{journalctl\ -}\allowbreak{}\texttt{u\ kube-}\allowbreak{}\texttt{apiserver\ \textbackslash{}\textbar{}\ grep\ -}\allowbreak{}\texttt{i\ unauthorized}.
Verify with
\texttt{openssl\ verify\ -}\allowbreak{}\texttt{CAfile\ ca.}\allowbreak{}\texttt{crt\ \textless{}cert\textgreater{}.}\allowbreak{}\texttt{crt}. \\
Later:
\texttt{403\ Forbidden},
user named in the
message &
Certificate
presented fine;
subject is wrong for
the intended role &
Read the username in
the error and
compare to the table
above. \\
\texttt{mkdir:}\allowbreak{}\texttt{\ cannot\ create\ directory\ \textquotesingle{}/}\allowbreak{}\texttt{var/}\allowbreak{}\texttt{lib/}\allowbreak{}\texttt{kubelet/}\allowbreak{}\texttt{\textquotesingle{}:}\allowbreak{}\texttt{\ File\ exists}
& Re-running the
distribution loop &
Harmless. The
\texttt{scp}
commands after it
still execute. \\
\end{longtable}

\hypertarget{check-your-understanding-4}{%
\section{Check your
understanding}\label{check-your-understanding-4}}

\textbf{Q1. You
regenerate
\texttt{kube-}\allowbreak{}\texttt{scheduler.}\allowbreak{}\texttt{crt}
to fix a subject
typo, and copy it to
\texttt{server}. The
scheduler still
fails to
authenticate. Why?}

Because the
scheduler does not
read
\texttt{kube-}\allowbreak{}\texttt{scheduler.}\allowbreak{}\texttt{crt}.
It reads
\texttt{/}\allowbreak{}\texttt{var/}\allowbreak{}\texttt{lib/}\allowbreak{}\texttt{kubernetes/}\allowbreak{}\texttt{kube-}\allowbreak{}\texttt{scheduler.}\allowbreak{}\texttt{kubeconfig},
which --- thanks to
\texttt{-\/-embed-}\allowbreak{}\texttt{certs=true}
in Chapter 5 ---
contains a base64
copy of the
\emph{old}
certificate baked
inside it.
Regenerating a
certificate always
means regenerating
every kubeconfig
that embedded it,
then re-distributing
and restarting the
service.

\textbf{Q2.
\texttt{ca.key} is
copied to
\texttt{server}.
What is the
practical
consequence?}

Anyone with root on
\texttt{server} can
issue a certificate
with
\texttt{CN=admin,\ O=system:}\allowbreak{}\texttt{masters}
and hold permanent
cluster-admin. The
\texttt{kube-}\allowbreak{}\texttt{controller-}\allowbreak{}\texttt{manager}
needs the key for
the CSR signing
controller, so this
is not a mistake in
the tutorial --- it
is the actual design
of Kubernetes, and
it is why control
plane hosts are
treated as the
highest-trust
machines in any
production cluster.

\textbf{Q3. Why does
each worker get
\texttt{ca.crt} at
\texttt{/}\allowbreak{}\texttt{var/}\allowbreak{}\texttt{lib/}\allowbreak{}\texttt{kubelet/}\allowbreak{}\texttt{ca.}\allowbreak{}\texttt{crt}
when it already
trusts the API
server through its
kubeconfig?}

Two different
directions. The
kubeconfig's
embedded CA lets the
kubelet verify the
\emph{API server}
when it connects
outbound.
\texttt{/}\allowbreak{}\texttt{var/}\allowbreak{}\texttt{lib/}\allowbreak{}\texttt{kubelet/}\allowbreak{}\texttt{ca.}\allowbreak{}\texttt{crt}
is referenced by
\texttt{authentication.}\allowbreak{}\texttt{x509.}\allowbreak{}\texttt{clientCAFile}
in
\texttt{kubelet-}\allowbreak{}\texttt{config.}\allowbreak{}\texttt{yaml},
and lets the kubelet
verify
\emph{inbound}
clients ---
specifically the API
server connecting to
port 10250 for
\texttt{logs},
\texttt{exec}, and
\texttt{port-forward}.
Same CA, opposite
roles.

\hypertarget{chap-5}{%
\chapter{Generating
Kubernetes
Configuration Files
for
Authentication}\label{chap-5}}

\upstreamlab{docs/05-kubernetes-configuration-files.md}

\hypertarget{what-youre-building-5}{%
\section{What you're
building}\label{what-youre-building-5}}

Six kubeconfig files
--- \texttt{node-0},
\texttt{node-1},
\texttt{kube-proxy},
\texttt{kube-}\allowbreak{}\texttt{controller-}\allowbreak{}\texttt{manager},
\texttt{kube-scheduler},
and \texttt{admin}
--- each bundling
three things: which
API server to
contact, which CA to
trust, and which
client certificate
to present. Then
distributing them to
the machines that
need them.

\hypertarget{why-its-necessary-5}{%
\section{Why it's
necessary}\label{why-its-necessary-5}}

\begin{whybox}

A kubeconfig is a
credential bundle.
Every component in
this cluster except
the API server
itself is an API
\emph{client}, and
every API client
needs the same three
answers. Rather than
pass six flags to
each service,
Kubernetes
standardizes on a
single YAML file
with a
\texttt{-\/-kubeconfig}
flag.

\end{whybox}

Skipping or
misgenerating one
produces a component
that starts cleanly
and then fails on
every request.
\texttt{kube-}\allowbreak{}\texttt{controller-}\allowbreak{}\texttt{manager}
with a bad
kubeconfig will
start, log
\texttt{Unauthorized}
in a loop, and no
deployment will ever
create a pod ---
while
\texttt{systemctl\ is-}\allowbreak{}\texttt{active}
cheerfully reports
\texttt{active}.

\hypertarget{key-concepts-5}{%
\section{Key
concepts}\label{key-concepts-5}}

\textbf{1. A
kubeconfig has three
parts and a join.}
\texttt{set-cluster}
defines \emph{where
and what to trust}.
\texttt{set-credentials}
defines \emph{who
you are}.
\texttt{set-context}
\textbf{joins} a
cluster to a user.
\texttt{use-context}
selects the active
join. Omit
\texttt{use-context}
and the file has a
context defined but
none selected, and
clients fail with
\texttt{no\ context\ chosen}.

\textbf{2.
\texttt{-\/-embed-}\allowbreak{}\texttt{certs=true}
copies, it does not
reference.} The
certificate and key
are base64-encoded
directly into the
file. This is what
makes a kubeconfig
portable --- you
\texttt{scp} one
file, not three. It
is also why
regenerating a
certificate
invalidates every
kubeconfig built
from it.
Consequence:
\textbf{the
kubeconfigs are as
sensitive as private
keys}, because they
contain private
keys.

\textbf{3. The two
\texttt{admin.}\allowbreak{}\texttt{kubeconfig}
files are different
files.} This chapter
generates
\texttt{admin.}\allowbreak{}\texttt{kubeconfig}
pointing at
\texttt{https:}\allowbreak{}\texttt{/}\allowbreak{}\texttt{/}\allowbreak{}\texttt{127.}\allowbreak{}\texttt{0.}\allowbreak{}\texttt{0.}\allowbreak{}\texttt{1:}\allowbreak{}\texttt{6443}
and ships it to
\texttt{server}.
Chapter 10 generates
\texttt{\textasciitilde{}/.kube/config}
on the jumpbox
pointing at
\texttt{https:}\allowbreak{}\texttt{/}\allowbreak{}\texttt{/}\allowbreak{}\texttt{server.}\allowbreak{}\texttt{kubernetes.}\allowbreak{}\texttt{local:}\allowbreak{}\texttt{6443}.
Both authenticate as
\texttt{admin}; they
differ only in
server address,
because one is used
\emph{on} the
control plane host
and one is used
remotely. This is
the single most
confusing detail in
the tutorial, and it
explains why Chapter
8's verification
command carries an
explicit
\texttt{-\/-kubeconfig\ admin.}\allowbreak{}\texttt{kubeconfig}
while Chapter 10's
does not.

\textbf{4. Each
kubelet gets its own
kubeconfig.}
\texttt{system:}\allowbreak{}\texttt{node:}\allowbreak{}\texttt{node-}\allowbreak{}\texttt{0}
and
\texttt{system:}\allowbreak{}\texttt{node:}\allowbreak{}\texttt{node-}\allowbreak{}\texttt{1}
are distinct
usernames, so the
files cannot be
shared. The
\texttt{for\ host\ in\ node-}\allowbreak{}\texttt{0\ node-}\allowbreak{}\texttt{1}
loop generates them
from the per-host
certificates
produced in Chapter
4.

\hypertarget{the-work-4}{%
\section{The
work}\label{the-work-4}}

All generation runs
on the
\textbf{jumpbox}, in
the same directory
as the certificates
from Chapter 4.

\textbf{Kubelet
kubeconfigs:}

\begin{cmdlst}
# on: jumpbox
for host in node-0 node-1; do
  kubectl config set-cluster kubernetes-the-hard-way \
    --certificate-authority=ca.crt \
    --embed-certs=true \
    --server=https://server.kubernetes.local:6443 \
    --kubeconfig=${host}.kubeconfig

  kubectl config set-credentials system:node:${host} \
    --client-certificate=${host}.crt \
    --client-key=${host}.key \
    --embed-certs=true \
    --kubeconfig=${host}.kubeconfig

  kubectl config set-context default \
    --cluster=kubernetes-the-hard-way \
    --user=system:node:${host} \
    --kubeconfig=${host}.kubeconfig

  kubectl config use-context default \
    --kubeconfig=${host}.kubeconfig
done
\end{cmdlst}

Results:

\begin{outlst}
node-0.kubeconfig
node-1.kubeconfig
\end{outlst}

The
\texttt{-\/-kubeconfig}
flag on every
subcommand is what
keeps
\texttt{kubectl}
from writing to
\texttt{\textasciitilde{}/.kube/config}.
Omit it once and you
silently corrupt
your default context
instead of building
the file you wanted.

\textbf{kube-proxy:}

\begin{cmdlst}
# on: jumpbox
{
  kubectl config set-cluster kubernetes-the-hard-way \
    --certificate-authority=ca.crt \
    --embed-certs=true \
    --server=https://server.kubernetes.local:6443 \
    --kubeconfig=kube-proxy.kubeconfig

  kubectl config set-credentials system:kube-proxy \
    --client-certificate=kube-proxy.crt \
    --client-key=kube-proxy.key \
    --embed-certs=true \
    --kubeconfig=kube-proxy.kubeconfig

  kubectl config set-context default \
    --cluster=kubernetes-the-hard-way \
    --user=system:kube-proxy \
    --kubeconfig=kube-proxy.kubeconfig

  kubectl config use-context default \
    --kubeconfig=kube-proxy.kubeconfig
}
\end{cmdlst}

Results:

\begin{outlst}
kube-proxy.kubeconfig
\end{outlst}

\textbf{kube-controller-manager:}

\begin{cmdlst}
# on: jumpbox
{
  kubectl config set-cluster kubernetes-the-hard-way \
    --certificate-authority=ca.crt \
    --embed-certs=true \
    --server=https://server.kubernetes.local:6443 \
    --kubeconfig=kube-controller-manager.kubeconfig

  kubectl config set-credentials system:kube-controller-manager \
    --client-certificate=kube-controller-manager.crt \
    --client-key=kube-controller-manager.key \
    --embed-certs=true \
    --kubeconfig=kube-controller-manager.kubeconfig

  kubectl config set-context default \
    --cluster=kubernetes-the-hard-way \
    --user=system:kube-controller-manager \
    --kubeconfig=kube-controller-manager.kubeconfig

  kubectl config use-context default \
    --kubeconfig=kube-controller-manager.kubeconfig
}
\end{cmdlst}

Results:

\begin{outlst}
kube-controller-manager.kubeconfig
\end{outlst}

\textbf{kube-scheduler:}

\begin{cmdlst}
# on: jumpbox
{
  kubectl config set-cluster kubernetes-the-hard-way \
    --certificate-authority=ca.crt \
    --embed-certs=true \
    --server=https://server.kubernetes.local:6443 \
    --kubeconfig=kube-scheduler.kubeconfig

  kubectl config set-credentials system:kube-scheduler \
    --client-certificate=kube-scheduler.crt \
    --client-key=kube-scheduler.key \
    --embed-certs=true \
    --kubeconfig=kube-scheduler.kubeconfig

  kubectl config set-context default \
    --cluster=kubernetes-the-hard-way \
    --user=system:kube-scheduler \
    --kubeconfig=kube-scheduler.kubeconfig

  kubectl config use-context default \
    --kubeconfig=kube-scheduler.kubeconfig
}
\end{cmdlst}

Results:

\begin{outlst}
kube-scheduler.kubeconfig
\end{outlst}

\textbf{admin:}

\begin{cmdlst}
# on: jumpbox
{
  kubectl config set-cluster kubernetes-the-hard-way \
    --certificate-authority=ca.crt \
    --embed-certs=true \
    --server=https://127.0.0.1:6443 \
    --kubeconfig=admin.kubeconfig

  kubectl config set-credentials admin \
    --client-certificate=admin.crt \
    --client-key=admin.key \
    --embed-certs=true \
    --kubeconfig=admin.kubeconfig

  kubectl config set-context default \
    --cluster=kubernetes-the-hard-way \
    --user=admin \
    --kubeconfig=admin.kubeconfig

  kubectl config use-context default \
    --kubeconfig=admin.kubeconfig
}
\end{cmdlst}

Results:

\begin{outlst}
admin.kubeconfig
\end{outlst}

Note
\texttt{-\/-server=https:}\allowbreak{}\texttt{/}\allowbreak{}\texttt{/}\allowbreak{}\texttt{127.}\allowbreak{}\texttt{0.}\allowbreak{}\texttt{0.}\allowbreak{}\texttt{1:}\allowbreak{}\texttt{6443}
here, unlike every
other kubeconfig
above. This file is
destined for
\texttt{server},
where the API server
is local.
\texttt{127.0.0.1}
is in the API server
certificate's SANs,
so TLS verification
succeeds.

\textbf{Distribute
to the workers:}

\begin{cmdlst}
# on: jumpbox
for host in node-0 node-1; do
  ssh root@${host} "mkdir -p /var/lib/{kube-proxy,kubelet}"

  scp kube-proxy.kubeconfig \
    root@${host}:/var/lib/kube-proxy/kubeconfig \

  scp ${host}.kubeconfig \
    root@${host}:/var/lib/kubelet/kubeconfig
done
\end{cmdlst}

Both files are
\textbf{renamed to
\texttt{kubeconfig}}
at their
destination. This is
what the systemd
units and config
files expect:
\texttt{kubelet.service}
passes
\texttt{-\/-kubeconfig=/}\allowbreak{}\texttt{var/}\allowbreak{}\texttt{lib/}\allowbreak{}\texttt{kubelet/}\allowbreak{}\texttt{kubeconfig},
and
\texttt{kube-}\allowbreak{}\texttt{proxy-}\allowbreak{}\texttt{config.}\allowbreak{}\texttt{yaml}
sets
\texttt{clientConnection.}\allowbreak{}\texttt{kubeconfig:}\allowbreak{}\texttt{\ "/}\allowbreak{}\texttt{var/}\allowbreak{}\texttt{lib/}\allowbreak{}\texttt{kube-}\allowbreak{}\texttt{proxy/}\allowbreak{}\texttt{kubeconfig"}.
The rename is what
makes the identical
config work on both
nodes.

\textbf{Distribute
to the control
plane:}

\begin{cmdlst}
# on: jumpbox
scp admin.kubeconfig \
  kube-controller-manager.kubeconfig \
  kube-scheduler.kubeconfig \
  root@server:~/
\end{cmdlst}

These land in root's
home directory on
\texttt{server} and
are moved into place
in Chapter 8.

\hypertarget{verification-4}{%
\section{Verification}\label{verification-4}}

Inspect a generated
file before trusting
it:

\begin{beyondbox}\calloutlabel{teal}{BEYOND THE SOURCE}

\texttt{kubectl\ config\ view\ -\/-kubeconfig=node-}\allowbreak{}\texttt{0.}\allowbreak{}\texttt{kubeconfig}
prints the structure
with credentials
redacted. Check
three things:
\texttt{current-}\allowbreak{}\texttt{context:}\allowbreak{}\texttt{\ default}
is present, the
\texttt{server:} URL
is correct for the
file's destination,
and the
\texttt{user:} name
matches the
certificate's CN. A
file missing
\texttt{current-context}
means
\texttt{use-context}
did not run.

\end{beyondbox}

You can also confirm
the embedded
certificate is the
one you expect:

\begin{beyondbox}\calloutlabel{teal}{BEYOND THE SOURCE}

\texttt{kubectl\ config\ view\ -\/-kubeconfig=node-}\allowbreak{}\texttt{0.}\allowbreak{}\texttt{kubeconfig\ -\/-raw\ -}\allowbreak{}\texttt{o\ jsonpath=\textquotesingle{}\{.}\allowbreak{}\texttt{users{[}0{]}.}\allowbreak{}\texttt{user.}\allowbreak{}\texttt{client-}\allowbreak{}\texttt{certificate-}\allowbreak{}\texttt{data\}\textquotesingle{}\ \textbar{}\ base64\ -}\allowbreak{}\texttt{d\ \textbar{}\ openssl\ x509\ -}\allowbreak{}\texttt{noout\ -}\allowbreak{}\texttt{subject}
should print
\texttt{subject=CN=system:}\allowbreak{}\texttt{node:}\allowbreak{}\texttt{node-}\allowbreak{}\texttt{0,\ .}\allowbreak{}\texttt{.}\allowbreak{}\texttt{.}\allowbreak{}.

\end{beyondbox}

\hypertarget{when-it-goes-wrong-4}{%
\section{When it
goes
wrong}\label{when-it-goes-wrong-4}}

\begin{longtable}[]{@{}
  >{\raggedright\arraybackslash}p{(\columnwidth - 4\tabcolsep) * \real{0.3200}}
  >{\raggedright\arraybackslash}p{(\columnwidth - 4\tabcolsep) * \real{0.2400}}
  >{\raggedright\arraybackslash}p{(\columnwidth - 4\tabcolsep) * \real{0.4400}}@{}}
\toprule\noalign{}
\begin{minipage}[b]{\linewidth}\raggedright
Symptom
\end{minipage} &
\begin{minipage}[b]{\linewidth}\raggedright
Cause
\end{minipage} &
\begin{minipage}[b]{\linewidth}\raggedright
Diagnosis
\end{minipage} \\
\midrule\noalign{}
\endhead
\bottomrule\noalign{}
\endlastfoot
Your own
\texttt{kubectl}
stops working & A
\texttt{-\/-kubeconfig}
flag was omitted and
the command wrote to
\texttt{\textasciitilde{}/.kube/config}
&
\texttt{kubectl\ config\ view}
on the jumpbox.
Re-run Chapter 10 to
rebuild it. \\
Component logs
\texttt{Unauthorized}
in a loop & Embedded
certificate is stale
or signed by a
different CA &
Extract and inspect
it with the command
above; regenerate
the kubeconfig. \\
\texttt{error:}\allowbreak{}\texttt{\ no\ context\ chosen}
/
\texttt{context\ "default"\ does\ not\ exist}
&
\texttt{use-context}
did not run & Re-run
the full block for
that file. \\
\texttt{x509:}\allowbreak{}\texttt{\ certificate\ signed\ by\ unknown\ authority}
& \texttt{ca.crt}
changed after the
kubeconfig was built
& Regenerate every
kubeconfig. A new CA
invalidates all of
them at once. \\
kubelet:
\texttt{connection\ refused}
to
\texttt{server.}\allowbreak{}\texttt{kubernetes.}\allowbreak{}\texttt{local:}\allowbreak{}\texttt{6443}
& Correct config,
API server not up
yet & Expected until
Chapter 8
completes. \\
\texttt{scp:}\allowbreak{}\texttt{\ /}\allowbreak{}\texttt{var/}\allowbreak{}\texttt{lib/}\allowbreak{}\texttt{kubelet/}\allowbreak{}\texttt{kubeconfig:}\allowbreak{}\texttt{\ No\ such\ file\ or\ directory}
& The
\texttt{mkdir\ -p}
step did not run &
Re-run the
distribution loop
from the top. \\
\end{longtable}

\hypertarget{check-your-understanding-5}{%
\section{Check your
understanding}\label{check-your-understanding-5}}

\textbf{Q1. Why does
\texttt{kube-}\allowbreak{}\texttt{proxy.}\allowbreak{}\texttt{kubeconfig}
work identically on
both nodes while the
kubelet kubeconfigs
cannot be shared?}

\texttt{kube-proxy}
authenticates as a
single cluster-wide
identity,
\texttt{system:}\allowbreak{}\texttt{kube-}\allowbreak{}\texttt{proxy},
authorized through
the
\texttt{system:}\allowbreak{}\texttt{node-}\allowbreak{}\texttt{proxier}
group --- its
permissions do not
vary by node.
Kubelets
authenticate as
\texttt{system:}\allowbreak{}\texttt{node:}\allowbreak{}\texttt{\textless{}nodeName\textgreater{}},
and the Node
authorizer uses that
exact name to scope
what each one may
read. Sharing a
kubelet kubeconfig
would make
\texttt{node-1}
claim to be
\texttt{node-0}, and
the Node authorizer
would deny access to
\texttt{node-1}'s
own pods.

\textbf{Q2. You
change the server's
IP address. Which
files must you
regenerate?}

None of the
kubeconfigs --- they
reference
\texttt{server.}\allowbreak{}\texttt{kubernetes.}\allowbreak{}\texttt{local},
not an IP. You
update
\texttt{/etc/hosts}
on all four machines
and
\texttt{machines.txt}
on the jumpbox. The
one exception is
\texttt{admin.}\allowbreak{}\texttt{kubeconfig},
which points at
\texttt{127.0.0.1}
and is therefore
address-independent
by construction.
This is exactly the
portability that
Chapter 3's naming
scheme bought.

\textbf{Q3. Could
you skip the
kubeconfigs entirely
and pass
\texttt{-\/-client-}\allowbreak{}\texttt{certificate}
and
\texttt{-\/-client-key}
flags to each
component?}

Not for these
components.
\texttt{kubelet.service},
\texttt{kube-}\allowbreak{}\texttt{controller-}\allowbreak{}\texttt{manager.}\allowbreak{}\texttt{service},
and
\texttt{kube-}\allowbreak{}\texttt{scheduler.}\allowbreak{}\texttt{yaml}
all accept a
kubeconfig path and
no equivalent
individual flags for
client credentials.
The kubeconfig is
the interface.
\texttt{kubectl} is
the only client here
that accepts both
forms.

\hypertarget{chap-6}{%
\chapter{Generating
the Data Encryption
Config and
Key}\label{chap-6}}

\upstreamlab{docs/06-data-encryption-keys.md}

\hypertarget{what-youre-building-6}{%
\section{What you're
building}\label{what-youre-building-6}}

A 32-byte random key
and an
\texttt{EncryptionConfiguration}
file that tells the
API server to
encrypt Secret
objects before
writing them to
etcd.

\hypertarget{why-its-necessary-6}{%
\section{Why it's
necessary}\label{why-its-necessary-6}}

\begin{whybox}

By default,
Kubernetes Secrets
are stored in etcd
\textbf{base64-encoded,
not encrypted}.
Base64 is an
encoding, not a
cipher --- anyone
with read access to
etcd's data
directory, a backup,
or a disk snapshot
can decode every
Secret in the
cluster in one
command. Encryption
at rest closes that
gap.

\end{whybox}

The gap is wider in
this specific build
than in most,
because Chapter 7's
etcd listens on
plain HTTP with no
authentication.
Anything with access
to
\texttt{127.0.0.1:2379}
on \texttt{server}
can read the entire
datastore.
Encryption at rest
is the only thing
standing between
that access and your
Secrets, which is
why Chapter 12's
first smoke test is
a hexdump proving it
works.

\hypertarget{key-concepts-6}{%
\section{Key
concepts}\label{key-concepts-6}}

\textbf{1. The
provider order in
the file is the
policy.}

\begin{outlst}
providers:
  - aescbc:
      keys:
        - name: key1
          secret: ${ENCRYPTION_KEY}
  - identity: {}
\end{outlst}

For \textbf{writes},
the API server
always uses the
\emph{first}
provider ---
\texttt{aescbc}. For
\textbf{reads}, it
tries each provider
in order until one
succeeds.
\texttt{identity}
means ``no
encryption,'' so
listing it second
lets the API server
still read
pre-existing
plaintext Secrets.
Reverse the order
and you get the
opposite policy:
writes in plaintext,
reads that can
decrypt old
ciphertext. That
reversal is a common
and quiet
misconfiguration,
and nothing in the
cluster will
complain about it.

\textbf{2. Only
\texttt{secrets} are
encrypted here.} The
\texttt{resources:}
list names
\texttt{secrets} and
nothing else.
ConfigMaps, pod
specs, and
everything else
remain plaintext in
etcd. That is
upstream's scoping
choice; production
configurations often
extend the list.

\textbf{3. The key
is generated but
never persisted.}
\texttt{export\ ENCRYPTION\_KEY=.}\allowbreak{}\texttt{.}\allowbreak{}\texttt{.}\allowbreak{}
lives in your
shell's environment.
It is substituted
into the YAML and
then forgotten.
\textbf{If you lose
\texttt{encryption-}\allowbreak{}\texttt{config.}\allowbreak{}\texttt{yaml},
every encrypted
Secret in etcd is
permanently
unrecoverable.}
There is no escrow,
no backup, no
recovery path.
Losing this file is
worse than losing
the CA key --- you
can reissue
certificates, but
you cannot decrypt
without the key.

\textbf{4.
\texttt{envsubst}
needs the variable
exported.} The
template contains
\texttt{\$\{ENCRYPTION\_KEY\}}.
\texttt{envsubst}
reads from the
\emph{environment},
so a plain
\texttt{ENCRYPTION\_KEY=.}\allowbreak{}\texttt{.}\allowbreak{}\texttt{.}\allowbreak{}
assignment without
\texttt{export}
produces a file with
a literal empty
value and the API
server fails to
start.

\hypertarget{the-work-5}{%
\section{The
work}\label{the-work-5}}

Run on the
\textbf{jumpbox}.

Generate the
encryption key:

\begin{cmdlst}
# on: jumpbox
export ENCRYPTION_KEY=$(head -c 32 /dev/urandom | base64)
\end{cmdlst}

32 bytes is the
AES-256 key size
required by the
\texttt{aescbc}
provider.
\texttt{/dev/urandom}
is the correct
source --- it is
non-blocking and
cryptographically
suitable on any
modern Linux kernel.

Create the
encryption config
file:

\begin{cmdlst}
# on: jumpbox
envsubst < configs/encryption-config.yaml \
  > encryption-config.yaml
\end{cmdlst}

\texttt{envsubst}
ships in the
\texttt{gettext-base}
package, present by
default on Debian
12. Note the input
is
\texttt{configs/}\allowbreak{}\texttt{encryption-}\allowbreak{}\texttt{config.}\allowbreak{}\texttt{yaml}
(the template) and
the output is
\texttt{encryption-}\allowbreak{}\texttt{config.}\allowbreak{}\texttt{yaml}
in the current
directory (the
rendered file) ---
two different paths,
one character apart
in appearance.

Copy it to the
control plane:

\begin{cmdlst}
# on: jumpbox
scp encryption-config.yaml root@server:~/
\end{cmdlst}

\hypertarget{verification-5}{%
\section{Verification}\label{verification-5}}

Confirm the
substitution
actually happened:

\begin{beyondbox}\calloutlabel{teal}{BEYOND THE SOURCE}

\texttt{grep\ secret:}\allowbreak{}\texttt{\ encryption-}\allowbreak{}\texttt{config.}\allowbreak{}\texttt{yaml}
should show a
44-character base64
string, not the
literal text
\texttt{\$\{ENCRYPTION\_KEY\}}
and not an empty
value. This is a
five-second check
that prevents a
confusing API server
startup failure in
Chapter 8.

\end{beyondbox}

Full verification
--- that encryption
is actually applied
--- happens in
Chapter 12's hexdump
test, which is the
first thing the
smoke test does.

\hypertarget{when-it-goes-wrong-5}{%
\section{When it
goes
wrong}\label{when-it-goes-wrong-5}}

\begin{longtable}[]{@{}
  >{\raggedright\arraybackslash}p{(\columnwidth - 4\tabcolsep) * \real{0.3200}}
  >{\raggedright\arraybackslash}p{(\columnwidth - 4\tabcolsep) * \real{0.2400}}
  >{\raggedright\arraybackslash}p{(\columnwidth - 4\tabcolsep) * \real{0.4400}}@{}}
\toprule\noalign{}
\begin{minipage}[b]{\linewidth}\raggedright
Symptom
\end{minipage} &
\begin{minipage}[b]{\linewidth}\raggedright
Cause
\end{minipage} &
\begin{minipage}[b]{\linewidth}\raggedright
Diagnosis
\end{minipage} \\
\midrule\noalign{}
\endhead
\bottomrule\noalign{}
\endlastfoot
\texttt{secret:} is
empty or literal
\texttt{\$\{ENCRYPTION\_KEY\}}
& Variable not
exported, or
\texttt{envsubst}
missing &
\texttt{echo\ \$ENCRYPTION\_KEY}
in the same shell;
\texttt{which\ envsubst}.
Re-export and
re-render. \\
\texttt{envsubst:}\allowbreak{}\texttt{\ command\ not\ found}
&
\texttt{gettext-base}
not installed &
\texttt{apt-}\allowbreak{}\texttt{get\ -}\allowbreak{}\texttt{y\ install\ gettext-}\allowbreak{}\texttt{base},
then re-run. \\
Chapter 8: API
server exits
immediately after
start & Malformed
encryption config &
\texttt{journalctl\ -}\allowbreak{}\texttt{u\ kube-}\allowbreak{}\texttt{apiserver\ -}\allowbreak{}\texttt{n\ 30}
on \texttt{server}.
Errors about the
encryption provider
config point
here. \\
Chapter 8:
\texttt{error\ while\ parsing\ file:}\allowbreak{}\texttt{\ illegal\ base64\ data}
& The key was
truncated or has
stray whitespace &
Regenerate the key
and re-render the
file. \\
Chapter 12: hexdump
shows readable
\texttt{mydata} &
The API server is
not using this
config, or
\texttt{identity} is
first & On
\texttt{server}:
confirm
\texttt{/}\allowbreak{}\texttt{var/}\allowbreak{}\texttt{lib/}\allowbreak{}\texttt{kubernetes/}\allowbreak{}\texttt{encryption-}\allowbreak{}\texttt{config.}\allowbreak{}\texttt{yaml}
exists and the API
server unit has
\texttt{-\/-encryption-}\allowbreak{}\texttt{provider-}\allowbreak{}\texttt{config}
pointing at it. \\
\end{longtable}

\hypertarget{check-your-understanding-6}{%
\section{Check your
understanding}\label{check-your-understanding-6}}

\textbf{Q1. You
rotate the
encryption key by
generating a new one
and replacing the
file. What happens
to existing
Secrets?}

They become
unreadable --- the
API server will try
\texttt{aescbc} with
the new key, fail to
decrypt, then fall
through to
\texttt{identity},
which cannot decrypt
ciphertext either.
Correct rotation
requires listing the
\emph{old} key
alongside the new
one under
\texttt{keys:} (new
first, old second),
restarting the API
server, then
rewriting every
Secret
(\texttt{kubectl\ get\ secrets\ -}\allowbreak{}\texttt{A\ -}\allowbreak{}\texttt{o\ json\ \textbar{}\ kubectl\ replace\ -}\allowbreak{}\texttt{f\ -}\allowbreak{})
so they are
re-encrypted with
the new key, and
only then removing
the old key.

\textbf{Q2. Why does
this chapter come
before etcd is even
running?}

Because
\texttt{encryption-}\allowbreak{}\texttt{config.}\allowbreak{}\texttt{yaml}
is an API server
startup dependency,
and Chapter 8 moves
it into
\texttt{/}\allowbreak{}\texttt{var/}\allowbreak{}\texttt{lib/}\allowbreak{}\texttt{kubernetes/}\allowbreak{}
alongside the
certificates.
Generating it here
keeps all
jumpbox-side
artifact creation in
Chapters 4--6 and
all control plane
installation in
Chapters 7--8. The
ordering is about
clean phases, not a
technical
constraint.

\textbf{Q3. If etcd
were configured with
TLS and client
certificate
authentication,
would encryption at
rest still matter?}

Yes. TLS protects
data \emph{in
transit} between the
API server and etcd.
Encryption at rest
protects data
\emph{on disk} ---
in
\texttt{/var/lib/etcd},
in backups, in
snapshots, on a
decommissioned
drive. They defend
against different
adversaries and
neither substitutes
for the other.

\hypertarget{chap-7}{%
\chapter{Bootstrapping
the etcd
Cluster}\label{chap-7}}

\upstreamlab{docs/07-bootstrapping-etcd.md}

\hypertarget{what-youre-building-7}{%
\section{What you're
building}\label{what-youre-building-7}}

A single-member etcd
cluster on
\texttt{server},
running under
systemd, listening
on loopback only,
with its data
directory at
\texttt{/var/lib/etcd}.

\hypertarget{why-its-necessary-7}{%
\section{Why it's
necessary}\label{why-its-necessary-7}}

\begin{whybox}

etcd holds
\textbf{all} cluster
state. Every
Deployment, Secret,
ConfigMap, node
registration, and
lease is a key in
etcd. The API server
is the only
component that talks
to it, and the API
server will not
start without it ---
it fails its
readiness check and
logs
\texttt{etcdserver:}\allowbreak{}\texttt{\ request\ timed\ out}
or connection
refused in a loop.

\end{whybox}

The ordering is
strict: etcd must be
running and healthy
before Chapter 8
starts the API
server. Nothing else
in the cluster
depends on etcd
directly.

\hypertarget{key-concepts-7}{%
\section{Key
concepts}\label{key-concepts-7}}

\textbf{1. One
member, no quorum,
no fault tolerance.}
The unit sets
\texttt{-\/-initial-}\allowbreak{}\texttt{cluster\ controller=http:}\allowbreak{}\texttt{/}\allowbreak{}\texttt{/}\allowbreak{}\texttt{127.}\allowbreak{}\texttt{0.}\allowbreak{}\texttt{0.}\allowbreak{}\texttt{1:}\allowbreak{}\texttt{2380}
--- a single-member
cluster. Raft with
one member means
every write commits
immediately with no
consensus round
trip, which is fast
and simple. It also
means \textbf{etcd
is a single point of
failure with no
redundancy}: if
\texttt{/var/lib/etcd}
is corrupted or the
host dies, the
cluster's entire
state is gone.
Production runs
three or five
members precisely to
survive this.

\textbf{2. The
member is named
\texttt{controller},
not
\texttt{server}.}
The tutorial says
\emph{``Set the etcd
name to match the
hostname of the
current compute
instance''} but the
shipped unit
hardcodes
\texttt{-\/-name\ controller}.
Since this is a
single-member
cluster reached over
loopback, the name
is only an
identifier and the
mismatch is cosmetic
--- it just means
\texttt{etcdctl\ member\ list}
prints
\texttt{controller}
where you might
expect
\texttt{server}.
Worth knowing so the
output does not look
like an error.

\textbf{3. etcd here
is plain HTTP on
loopback, and this
is a significant
shortcut.} Every URL
in the unit is
\texttt{http:}\allowbreak{}\texttt{/}\allowbreak{}\texttt{/}\allowbreak{}\texttt{127.}\allowbreak{}\texttt{0.}\allowbreak{}\texttt{0.}\allowbreak{}\texttt{1}
--- no TLS, no
client certificate
authentication. The
API server connects
with
\texttt{-\/-etcd-}\allowbreak{}\texttt{servers=http:}\allowbreak{}\texttt{/}\allowbreak{}\texttt{/}\allowbreak{}\texttt{127.}\allowbreak{}\texttt{0.}\allowbreak{}\texttt{0.}\allowbreak{}\texttt{1:}\allowbreak{}\texttt{2379},
also plaintext. Two
consequences worth
stating plainly:

\begin{itemize}
\tightlist
\item
  The only access
  control is that
  etcd binds to
  loopback. Anything
  running as any
  user on
  \texttt{server}
  can read and write
  the entire cluster
  datastore with no
  credentials.
\item
  Production
  clusters use TLS
  with client
  certificates
  between the API
  server and etcd,
  and separate peer
  certificates
  between etcd
  members.
  Replicating that
  is a good
  follow-up
  exercise.
\end{itemize}

\textbf{4.
\texttt{/etc/etcd/}
receives
certificates that
this configuration
does not use.} The
setup step copies
\texttt{ca.crt},
\texttt{kube-}\allowbreak{}\texttt{api-}\allowbreak{}\texttt{server.}\allowbreak{}\texttt{key},
and
\texttt{kube-}\allowbreak{}\texttt{api-}\allowbreak{}\texttt{server.}\allowbreak{}\texttt{crt}
into
\texttt{/etc/etcd/},
but
\texttt{etcd.service}
as shipped
references none of
them --- all its
URLs are
\texttt{http://}.
These files are
vestigial in this
version of the
tutorial, left from
when etcd ran with
TLS. Copy them
anyway (the guide
reproduces the
source exactly), but
do not spend time
wondering what
consumes them:
nothing does. If you
later convert etcd
to TLS, they are
already staged.

\textbf{5.
\texttt{Type=notify}
matters.} etcd
signals systemd via
\texttt{sd\_notify}
once it is actually
ready to serve,
rather than when the
process forks.
\texttt{systemctl\ start\ etcd}
therefore blocks
until etcd is
genuinely up, which
is why there is no
sleep before the
verification step.

\textbf{6.
\texttt{chmod\ 700\ /}\allowbreak{}\texttt{var/}\allowbreak{}\texttt{lib/}\allowbreak{}\texttt{etcd}
is required, not
hygiene.} etcd
refuses to start if
its data directory
has permissions more
permissive than
\texttt{0700}. This
is an explicit check
in etcd, not a
convention.

\hypertarget{the-work-6}{%
\section{The
work}\label{the-work-6}}

\textbf{From the
jumpbox}, copy the
binaries and unit
file:

\begin{cmdlst}
# on: jumpbox
scp \
  downloads/controller/etcd \
  downloads/client/etcdctl \
  units/etcd.service \
  root@server:~/
\end{cmdlst}

Note that
\texttt{etcdctl}
comes from
\texttt{downloads/}\allowbreak{}\texttt{client/}\allowbreak{}
while \texttt{etcd}
comes from
\texttt{downloads/}\allowbreak{}\texttt{controller/}\allowbreak{}
--- the sort from
Chapter 2 pays off
here.

Log in to
\texttt{server}.
\textbf{Everything
below runs there.}

\begin{cmdlst}
# on: jumpbox
ssh root@server
\end{cmdlst}

Install the
binaries:

\begin{cmdlst}
# on: server
{
  mv etcd etcdctl /usr/local/bin/
}
\end{cmdlst}

Configure the etcd
server:

\begin{cmdlst}
# on: server
{
  mkdir -p /etc/etcd /var/lib/etcd
  chmod 700 /var/lib/etcd
  cp ca.crt kube-api-server.key kube-api-server.crt \
    /etc/etcd/
}
\end{cmdlst}

Note this is
\texttt{cp}, not
\texttt{mv} --- the
originals stay in
root's home
directory because
Chapter 8 moves them
to
\texttt{/}\allowbreak{}\texttt{var/}\allowbreak{}\texttt{lib/}\allowbreak{}\texttt{kubernetes/}\allowbreak{}.
Running \texttt{mv}
here breaks Chapter
8.

Install the systemd
unit:

\begin{cmdlst}
# on: server
mv etcd.service /etc/systemd/system/
\end{cmdlst}

For reference, the
unit you just
installed:

\begin{outlst}
[Unit]
Description=etcd
Documentation=https://github.com/etcd-io/etcd

[Service]
Type=notify
ExecStart=/usr/local/bin/etcd \
  --name controller \
  --initial-advertise-peer-urls http://127.0.0.1:2380 \
  --listen-peer-urls http://127.0.0.1:2380 \
  --listen-client-urls http://127.0.0.1:2379 \
  --advertise-client-urls http://127.0.0.1:2379 \
  --initial-cluster-token etcd-cluster-0 \
  --initial-cluster controller=http://127.0.0.1:2380 \
  --initial-cluster-state new \
  --data-dir=/var/lib/etcd
Restart=on-failure
RestartSec=5

[Install]
WantedBy=multi-user.target
\end{outlst}

Two port numbers,
two purposes:
\textbf{2380} is the
peer port
(etcd-to-etcd Raft
traffic --- unused
with one member, but
still bound),
\textbf{2379} is the
client port (what
the API server
connects to).
\texttt{-\/-initial-}\allowbreak{}\texttt{cluster-}\allowbreak{}\texttt{state\ new}
means ``form a new
cluster''; on a
restart of an
existing cluster
this would be
\texttt{existing},
and etcd ignores the
flag once
\texttt{/var/lib/etcd}
contains a member
directory.

Start etcd:

\begin{cmdlst}
# on: server
{
  systemctl daemon-reload
  systemctl enable etcd
  systemctl start etcd
}
\end{cmdlst}

\texttt{daemon-reload}
makes systemd
re-read unit files
from disk ---
required after any
\texttt{mv} into
\texttt{/}\allowbreak{}\texttt{etc/}\allowbreak{}\texttt{systemd/}\allowbreak{}\texttt{system/}\allowbreak{}.
\texttt{enable}
creates the
boot-time symlink;
\texttt{start} runs
it now. Both are
needed; neither
implies the other.

\hypertarget{verification-6}{%
\section{Verification}\label{verification-6}}

List the etcd
cluster members:

\begin{cmdlst}
# on: server
etcdctl member list
\end{cmdlst}

\begin{outlst}
6702b0a34e2cfd39, started, controller, http://127.0.0.1:2380, http://127.0.0.1:2379, false
\end{outlst}

The hex ID is
randomly generated
and will differ.
What matters is
\texttt{started},
the member name
\texttt{controller},
both URLs, and the
trailing
\texttt{false} (not
a learner member).
\texttt{etcdctl}
needs no flags
because it defaults
to
\texttt{http:}\allowbreak{}\texttt{/}\allowbreak{}\texttt{/}\allowbreak{}\texttt{127.}\allowbreak{}\texttt{0.}\allowbreak{}\texttt{0.}\allowbreak{}\texttt{1:}\allowbreak{}\texttt{2379},
which is exactly
where etcd is
listening.

\begin{beyondbox}\calloutlabel{teal}{BEYOND THE SOURCE}

\texttt{etcdctl\ endpoint\ health}
is a more direct
liveness check and
prints
\texttt{127.}\allowbreak{}\texttt{0.}\allowbreak{}\texttt{0.}\allowbreak{}\texttt{1:}\allowbreak{}\texttt{2379\ is\ healthy}.
\texttt{systemctl\ is-}\allowbreak{}\texttt{active\ etcd}
should print
\texttt{active}.

\end{beyondbox}

\hypertarget{when-it-goes-wrong-6}{%
\section{When it
goes
wrong}\label{when-it-goes-wrong-6}}

\begin{longtable}[]{@{}
  >{\raggedright\arraybackslash}p{(\columnwidth - 4\tabcolsep) * \real{0.3200}}
  >{\raggedright\arraybackslash}p{(\columnwidth - 4\tabcolsep) * \real{0.2400}}
  >{\raggedright\arraybackslash}p{(\columnwidth - 4\tabcolsep) * \real{0.4400}}@{}}
\toprule\noalign{}
\begin{minipage}[b]{\linewidth}\raggedright
Symptom
\end{minipage} &
\begin{minipage}[b]{\linewidth}\raggedright
Cause
\end{minipage} &
\begin{minipage}[b]{\linewidth}\raggedright
Diagnosis
\end{minipage} \\
\midrule\noalign{}
\endhead
\bottomrule\noalign{}
\endlastfoot
\texttt{etcdctl:}\allowbreak{}\texttt{\ command\ not\ found}
& The \texttt{mv}
did not run, or you
are on the wrong
host &
\texttt{ls\ -}\allowbreak{}\texttt{l\ /}\allowbreak{}\texttt{usr/}\allowbreak{}\texttt{local/}\allowbreak{}\texttt{bin/}\allowbreak{}\texttt{etcd*}.
Confirm you are on
\texttt{server}. \\
etcd exits at
startup,
\texttt{cannot\ access\ data\ directory}
&
\texttt{/var/lib/etcd}
is more permissive
than \texttt{0700} &
\texttt{stat\ -}\allowbreak{}\texttt{c\ \%a\ /}\allowbreak{}\texttt{var/}\allowbreak{}\texttt{lib/}\allowbreak{}\texttt{etcd}
\textrightarrow{}
must be
\texttt{700}. Fix
and restart. \\
\texttt{member\ list}
hangs, then times
out & etcd not
running, or
listening elsewhere
&
\texttt{systemctl\ status\ etcd};
\texttt{ss\ -}\allowbreak{}\texttt{tlnp\ \textbackslash{}\textbar{}\ grep\ 2379}. \\
\texttt{error\ validating\ peerURLs\ .}\allowbreak{}\texttt{.}\allowbreak{}\texttt{.}\allowbreak{}\texttt{\ member\ count\ is\ unequal}
&
\texttt{/var/lib/etcd}
contains data from a
previous,
differently-configured
run & This is a lab:
\texttt{systemctl\ stop\ etcd;\ rm\ -}\allowbreak{}\texttt{rf\ /}\allowbreak{}\texttt{var/}\allowbreak{}\texttt{lib/}\allowbreak{}\texttt{etcd/}\allowbreak{}\texttt{*;\ systemctl\ start\ etcd}.
Never do this to a
cluster you care
about. \\
Chapter 8: API
server logs
\texttt{etcdserver:}\allowbreak{}\texttt{\ request\ timed\ out}
& etcd is down or
unreachable on 2379
&
\texttt{etcdctl\ endpoint\ health}
on
\texttt{server}. \\
Chapter 8:
\texttt{mv:}\allowbreak{}\texttt{\ cannot\ stat\ \textquotesingle{}ca.}\allowbreak{}\texttt{crt\textquotesingle{}}
& You used
\texttt{mv} instead
of \texttt{cp} in
the configure step &
Re-copy the files
from the jumpbox
with the Chapter 4
\texttt{scp}
command. \\
\end{longtable}

\hypertarget{check-your-understanding-7}{%
\section{Check your
understanding}\label{check-your-understanding-7}}

\textbf{Q1. The API
server connects to
etcd over plaintext
HTTP on loopback.
What single change
would most improve
this, and what would
it require?}

Enabling TLS with
client certificate
authentication on
etcd:
\texttt{-\/-client-}\allowbreak{}\texttt{cert-}\allowbreak{}\texttt{auth},
\texttt{-\/-trusted-}\allowbreak{}\texttt{ca-}\allowbreak{}\texttt{file},
\texttt{-\/-cert-file},
\texttt{-\/-key-file}
on the etcd side,
and
\texttt{-\/-etcd-cafile},
\texttt{-\/-etcd-certfile},
\texttt{-\/-etcd-keyfile}
on the API server
side. The
certificates are
already staged in
\texttt{/etc/etcd/}
from this chapter's
setup step. The
change would mean
that gaining shell
on \texttt{server}
is no longer
sufficient to read
every Secret in the
cluster.

\textbf{Q2. You
restore
\texttt{/var/lib/etcd}
from a backup taken
an hour ago. What is
the state of the
cluster?}

Exactly as it was an
hour ago --- every
object, every
resourceVersion.
Running pods on the
workers are
unaffected in the
moment, because the
kubelets do not
consult etcd. But
the control plane's
view is now stale:
pods created in the
last hour do not
exist as far as it
is concerned, so the
kubelets' status
reports will be for
pods it has no
record of, and the
controller manager
will start
reconciling toward
the old state. etcd
is the cluster's
only durable truth,
and rolling it back
rolls back
everything.

\textbf{Q3. Why does
\texttt{-\/-initial-}\allowbreak{}\texttt{cluster-}\allowbreak{}\texttt{state\ new}
not cause problems
when etcd restarts?}

etcd only honours
the
\texttt{-\/-initial-}\allowbreak{}\texttt{cluster*}
flags when
bootstrapping into
an empty data
directory. Once
\texttt{/var/lib/etcd}
contains a member,
the on-disk cluster
configuration wins
and those flags are
ignored. This is
what makes
\texttt{Restart=on-}\allowbreak{}\texttt{failure}
safe.

\hypertarget{chap-8}{%
\chapter{Bootstrapping
the Kubernetes
Control
Plane}\label{chap-8}}

\upstreamlab{docs/08-bootstrapping-kubernetes-controllers.md}

\hypertarget{what-youre-building-8}{%
\section{What you're
building}\label{what-youre-building-8}}

The three control
plane processes on
\texttt{server} ---
\texttt{kube-apiserver},
\texttt{kube-}\allowbreak{}\texttt{controller-}\allowbreak{}\texttt{manager},
\texttt{kube-scheduler}
--- each as a
systemd service,
plus the RBAC
ClusterRole that
permits the API
server to reach the
kubelet API on
worker nodes.

\hypertarget{why-its-necessary-8}{%
\section{Why it's
necessary}\label{why-its-necessary-8}}

\begin{whybox}

This is the chapter
that turns a pile of
certificates into a
working Kubernetes
API. Everything
before it was
preparation;
everything after it
is a client of what
you build here.

\end{whybox}

The RBAC step at the
end is easy to skip
and produces a
specific, delayed
failure. Without it,
the cluster works
--- you can create
deployments, pods
get scheduled and
run --- but
\texttt{kubectl\ logs},
\texttt{kubectl\ exec},
and
\texttt{kubectl\ port-}\allowbreak{}\texttt{forward}
all fail with
\texttt{Forbidden\ (user=kubernetes,\ verb=get,\ resource=nodes,\ subresource=proxy)}.
That is because
those three commands
are proxied
\emph{through} the
API server to the
kubelet, and the
kubelet is
configured with
\texttt{authorization.}\allowbreak{}\texttt{mode:}\allowbreak{}\texttt{\ Webhook},
which asks the API
server whether the
caller may do that.
The caller, in this
case, is the API
server itself,
presenting
\texttt{kube-}\allowbreak{}\texttt{api-}\allowbreak{}\texttt{server.}\allowbreak{}\texttt{crt}
with
\texttt{CN=kubernetes}
--- an identity with
no permissions until
you grant them.

\hypertarget{key-concepts-8}{%
\section{Key
concepts}\label{key-concepts-8}}

\textbf{1.
\texttt{-\/-authorization-}\allowbreak{}\texttt{mode=Node,RBAC}
is an ordered
chain.} The Node
authorizer is
checked first; it
handles kubelet
requests and scopes
each node to its own
objects. RBAC
handles everything
else. A request is
allowed if
\emph{either}
permits it, denied
only if both
decline.

\textbf{2. Admission
plugins run after
authorization and
can mutate or
reject.} The unit
enables:

\begin{outlst}
NamespaceLifecycle,NodeRestriction,LimitRanger,ServiceAccount,DefaultStorageClass,ResourceQuota
\end{outlst}

\texttt{NodeRestriction}
is the
security-relevant
one here: it
prevents a kubelet
from modifying nodes
other than its own,
or from setting
labels that would
let it attract pods
it should not
receive. It is the
admission-time
complement to the
Node authorizer's
read-time scoping.

\textbf{3. The API
server holds three
distinct key roles.}
Reading the unit
carefully:

\begin{itemize}
\tightlist
\item
  \texttt{-\/-tls-cert-file}
  /
  \texttt{-\/-tls-}\allowbreak{}\texttt{private-}\allowbreak{}\texttt{key-}\allowbreak{}\texttt{file}
  \textrightarrow{}
  its
  \textbf{server}
  identity on :6443
\item
  \texttt{-\/-kubelet-}\allowbreak{}\texttt{client-}\allowbreak{}\texttt{certificate}
  /
  \texttt{-\/-kubelet-}\allowbreak{}\texttt{client-}\allowbreak{}\texttt{key}
  \textrightarrow{}
  its
  \textbf{client}
  identity when
  calling kubelets
  on :10250
\item
  \texttt{-\/-client-}\allowbreak{}\texttt{ca-}\allowbreak{}\texttt{file}
  \textrightarrow{}
  the CA used to
  \textbf{verify}
  incoming client
  certificates
\end{itemize}

The same
\texttt{kube-}\allowbreak{}\texttt{api-}\allowbreak{}\texttt{server.}\allowbreak{}\texttt{crt}
serves the first two
roles, which is why
\texttt{ca.conf}
gives it
\texttt{nsCertType\ =\ client,\ server}
--- the only
certificate in the
build with both.

\textbf{4.
ServiceAccount token
signing spans two
components.} The API
server signs tokens
with
\texttt{-\/-service-}\allowbreak{}\texttt{account-}\allowbreak{}\texttt{signing-}\allowbreak{}\texttt{key-}\allowbreak{}\texttt{file}
and verifies them
with
\texttt{-\/-service-}\allowbreak{}\texttt{account-}\allowbreak{}\texttt{key-}\allowbreak{}\texttt{file};
the controller
manager holds the
same key as
\texttt{-\/-service-}\allowbreak{}\texttt{account-}\allowbreak{}\texttt{private-}\allowbreak{}\texttt{key-}\allowbreak{}\texttt{file}
to create tokens for
new ServiceAccounts.
\texttt{-\/-service-}\allowbreak{}\texttt{account-}\allowbreak{}\texttt{issuer=https:}\allowbreak{}\texttt{/}\allowbreak{}\texttt{/}\allowbreak{}\texttt{server.}\allowbreak{}\texttt{kubernetes.}\allowbreak{}\texttt{local:}\allowbreak{}\texttt{6443}
becomes the
\texttt{iss} claim
in every token. If
these three
references disagree,
pods get tokens the
API server rejects.

\textbf{5. The
controller manager
needs the CA private
key.}
\texttt{-\/-cluster-}\allowbreak{}\texttt{signing-}\allowbreak{}\texttt{cert-}\allowbreak{}\texttt{file}
and
\texttt{-\/-cluster-}\allowbreak{}\texttt{signing-}\allowbreak{}\texttt{key-}\allowbreak{}\texttt{file}
point at
\texttt{ca.crt} and
\texttt{ca.key},
enabling the CSR
signing controller.
\texttt{-\/-use-}\allowbreak{}\texttt{service-}\allowbreak{}\texttt{account-}\allowbreak{}\texttt{credentials=true}
makes each
controller loop
authenticate as its
own ServiceAccount
rather than as the
controller manager
identity --- which
means a bug in one
controller has a
bounded blast
radius.

\textbf{6. The
scheduler is
configured by file,
not flags.}
\texttt{kube-}\allowbreak{}\texttt{scheduler.}\allowbreak{}\texttt{service}
passes only
\texttt{-\/-config=/}\allowbreak{}\texttt{etc/}\allowbreak{}\texttt{kubernetes/}\allowbreak{}\texttt{config/}\allowbreak{}\texttt{kube-}\allowbreak{}\texttt{scheduler.}\allowbreak{}\texttt{yaml},
and that file names
the kubeconfig and
enables leader
election. Leader
election with a
single replica is a
no-op that still
exercises the lease
machinery --- a
useful default that
means adding a
second scheduler
later requires no
config change.

\hypertarget{the-work-7}{%
\section{The
work}\label{the-work-7}}

\textbf{From the
jumpbox}, copy
binaries, units, and
configs:

\begin{cmdlst}
# on: jumpbox
scp \
  downloads/controller/kube-apiserver \
  downloads/controller/kube-controller-manager \
  downloads/controller/kube-scheduler \
  downloads/client/kubectl \
  units/kube-apiserver.service \
  units/kube-controller-manager.service \
  units/kube-scheduler.service \
  configs/kube-scheduler.yaml \
  configs/kube-apiserver-to-kubelet.yaml \
  root@server:~/
\end{cmdlst}

Log in to
\texttt{server}.
\textbf{Everything
below runs there
unless marked
otherwise.}

\begin{cmdlst}
# on: jumpbox
ssh root@server
\end{cmdlst}

Create the
configuration
directory:

\begin{cmdlst}
# on: server
mkdir -p /etc/kubernetes/config
\end{cmdlst}

Install the
binaries:

\begin{cmdlst}
# on: server
{
  mv kube-apiserver \
    kube-controller-manager \
    kube-scheduler kubectl \
    /usr/local/bin/
}
\end{cmdlst}

\textbf{Configure
the API server.}
Move the credentials
into place:

\begin{cmdlst}
# on: server
{
  mkdir -p /var/lib/kubernetes/

  mv ca.crt ca.key \
    kube-api-server.key kube-api-server.crt \
    service-accounts.key service-accounts.crt \
    encryption-config.yaml \
    /var/lib/kubernetes/
}
\end{cmdlst}

This is the
\texttt{mv} that
fails if Chapter 7
used \texttt{mv}
instead of
\texttt{cp} when
populating
\texttt{/etc/etcd/}.

\begin{cmdlst}
# on: server
mv kube-apiserver.service \
  /etc/systemd/system/kube-apiserver.service
\end{cmdlst}

The unit, annotated:

\begin{outlst}
[Service]
ExecStart=/usr/local/bin/kube-apiserver \
  --allow-privileged=true \                       # permit privileged containers (kube-proxy needs this)
  --audit-log-maxage=30 \                         # audit log retention, days
  --audit-log-maxbackup=3 \
  --audit-log-maxsize=100 \                       # MB before rotation
  --audit-log-path=/var/log/audit.log \           # audit logging on, default policy
  --authorization-mode=Node,RBAC \                # ordered authorizer chain
  --bind-address=0.0.0.0 \                        # listen on all interfaces, :6443
  --client-ca-file=/var/lib/kubernetes/ca.crt \   # CA for verifying client certs
  --enable-admission-plugins=NamespaceLifecycle,NodeRestriction,LimitRanger,ServiceAccount,DefaultStorageClass,ResourceQuota \
  --etcd-servers=http://127.0.0.1:2379 \          # plaintext, loopback only
  --event-ttl=1h \                                # Events expire after an hour
  --encryption-provider-config=/var/lib/kubernetes/encryption-config.yaml \
  --kubelet-certificate-authority=/var/lib/kubernetes/ca.crt \   # verify kubelet server certs
  --kubelet-client-certificate=/var/lib/kubernetes/kube-api-server.crt \
  --kubelet-client-key=/var/lib/kubernetes/kube-api-server.key \
  --runtime-config='api/all=true' \               # enable all API groups, incl. alpha/beta
  --service-account-key-file=/var/lib/kubernetes/service-accounts.crt \
  --service-account-signing-key-file=/var/lib/kubernetes/service-accounts.key \
  --service-account-issuer=https://server.kubernetes.local:6443 \
  --service-node-port-range=30000-32767 \         # NodePort allocation range
  --tls-cert-file=/var/lib/kubernetes/kube-api-server.crt \
  --tls-private-key-file=/var/lib/kubernetes/kube-api-server.key \
  --v=2
Restart=on-failure
RestartSec=5
\end{outlst}

\textbf{Configure
the controller
manager:}

\begin{cmdlst}
# on: server
mv kube-controller-manager.kubeconfig /var/lib/kubernetes/
\end{cmdlst}

\begin{cmdlst}
# on: server
mv kube-controller-manager.service /etc/systemd/system/
\end{cmdlst}

Its flags worth
noting:
\texttt{-\/-cluster-}\allowbreak{}\texttt{cidr=10.}\allowbreak{}\texttt{200.}\allowbreak{}\texttt{0.}\allowbreak{}\texttt{0/}\allowbreak{}\texttt{16}
(the pod network it
hands out per-node
ranges from),
\texttt{-\/-service-}\allowbreak{}\texttt{cluster-}\allowbreak{}\texttt{ip-}\allowbreak{}\texttt{range=10.}\allowbreak{}\texttt{32.}\allowbreak{}\texttt{0.}\allowbreak{}\texttt{0/}\allowbreak{}\texttt{24}
(must match the API
server's
understanding of the
service range),
\texttt{-\/-root-}\allowbreak{}\texttt{ca-}\allowbreak{}\texttt{file=/}\allowbreak{}\texttt{var/}\allowbreak{}\texttt{lib/}\allowbreak{}\texttt{kubernetes/}\allowbreak{}\texttt{ca.}\allowbreak{}\texttt{crt}
(injected into every
ServiceAccount's
token so pods can
verify the API
server).

\textbf{Configure
the scheduler:}

\begin{cmdlst}
# on: server
mv kube-scheduler.kubeconfig /var/lib/kubernetes/
\end{cmdlst}

\begin{cmdlst}
# on: server
mv kube-scheduler.yaml /etc/kubernetes/config/
\end{cmdlst}

\begin{cmdlst}
# on: server
mv kube-scheduler.service /etc/systemd/system/
\end{cmdlst}

Note the scheduler's
config file goes to
\texttt{/}\allowbreak{}\texttt{etc/}\allowbreak{}\texttt{kubernetes/}\allowbreak{}\texttt{config/}\allowbreak{}
while its kubeconfig
goes to
\texttt{/}\allowbreak{}\texttt{var/}\allowbreak{}\texttt{lib/}\allowbreak{}\texttt{kubernetes/}\allowbreak{}.
Two different
directories, and the
paths inside
\texttt{kube-}\allowbreak{}\texttt{scheduler.}\allowbreak{}\texttt{yaml}
depend on both being
right.

\textbf{Start all
three:}

\begin{cmdlst}
# on: server
{
  systemctl daemon-reload

  systemctl enable kube-apiserver \
    kube-controller-manager kube-scheduler

  systemctl start kube-apiserver \
    kube-controller-manager kube-scheduler
}
\end{cmdlst}

\begin{notebox}\calloutlabel{steel}{FROM THE SOURCE}

Allow up to 10
seconds for the
Kubernetes API
Server to fully
initialize.

\end{notebox}

Unlike etcd, these
units do not use
\texttt{Type=notify},
so systemd reports
success as soon as
the process starts
--- before the API
server is serving.
The controller
manager and
scheduler will log
connection errors
during this window
and retry; that is
expected and
self-correcting.

Check status:

\begin{cmdlst}
# on: server
systemctl is-active kube-apiserver
\end{cmdlst}

\begin{cmdlst}
# on: server
systemctl status kube-apiserver
\end{cmdlst}

\begin{cmdlst}
# on: server
journalctl -u kube-apiserver
\end{cmdlst}

\hypertarget{verification-7}{%
\section{Verification}\label{verification-7}}

\begin{cmdlst}
# on: server
kubectl cluster-info \
  --kubeconfig admin.kubeconfig
\end{cmdlst}

\begin{outlst}
Kubernetes control plane is running at https://127.0.0.1:6443
\end{outlst}

This works because
\texttt{admin.}\allowbreak{}\texttt{kubeconfig}
--- the one
generated in Chapter
5 --- points at
\texttt{127.0.0.1},
and you are on the
host where the API
server runs.

\textbf{Now apply
the RBAC for kubelet
authorization.}
These commands run
on \texttt{server}
and affect the whole
cluster:

\begin{cmdlst}
# on: jumpbox
ssh root@server
\end{cmdlst}

\begin{cmdlst}
# on: server
kubectl apply -f kube-apiserver-to-kubelet.yaml \
  --kubeconfig admin.kubeconfig
\end{cmdlst}

The manifest creates
a ClusterRole
granting \texttt{*}
on
\texttt{nodes/proxy},
\texttt{nodes/stats},
\texttt{nodes/log},
\texttt{nodes/spec},
and
\texttt{nodes/metrics},
and binds it to the
\textbf{User}
\texttt{kubernetes}
--- which is the CN
of the API server's
own certificate.
Those five
subresources are
precisely what
\texttt{kubectl\ logs},
\texttt{exec}, and
\texttt{port-forward}
traverse.

\textbf{Final
verification, from
the jumpbox:}

\begin{cmdlst}
# on: jumpbox
curl --cacert ca.crt \
  https://server.kubernetes.local:6443/version
\end{cmdlst}

\begin{outlst}
{
  "major": "1",
  "minor": "32",
  "gitVersion": "v1.32.3",
  "gitCommit": "32cc146f75aad04beaaa245a7157eb35063a9f99",
  "gitTreeState": "clean",
  "buildDate": "2025-03-11T19:52:21Z",
  "goVersion": "go1.23.6",
  "compiler": "gc",
  "platform": "linux/arm64"
}
\end{outlst}

\texttt{platform}
reflects your
architecture and
will read
\texttt{linux/amd64}
on an AMD64 build.
This single command
proves four things
at once: name
resolution works,
the API server is
listening, its
certificate chains
to your CA, and
\texttt{server.}\allowbreak{}\texttt{kubernetes.}\allowbreak{}\texttt{local}
is a valid SAN.
\texttt{/version} is
served
unauthenticated,
which is why no
client certificate
is needed.

\hypertarget{when-it-goes-wrong-7}{%
\section{When it
goes
wrong}\label{when-it-goes-wrong-7}}

\begin{longtable}[]{@{}
  >{\raggedright\arraybackslash}p{(\columnwidth - 4\tabcolsep) * \real{0.3200}}
  >{\raggedright\arraybackslash}p{(\columnwidth - 4\tabcolsep) * \real{0.2400}}
  >{\raggedright\arraybackslash}p{(\columnwidth - 4\tabcolsep) * \real{0.4400}}@{}}
\toprule\noalign{}
\begin{minipage}[b]{\linewidth}\raggedright
Symptom
\end{minipage} &
\begin{minipage}[b]{\linewidth}\raggedright
Cause
\end{minipage} &
\begin{minipage}[b]{\linewidth}\raggedright
Diagnosis
\end{minipage} \\
\midrule\noalign{}
\endhead
\bottomrule\noalign{}
\endlastfoot
API server starts
then exits within
seconds & Missing
credential file, or
malformed encryption
config &
\texttt{journalctl\ -}\allowbreak{}\texttt{u\ kube-}\allowbreak{}\texttt{apiserver\ -}\allowbreak{}\texttt{n\ 50\ -\/-no-}\allowbreak{}\texttt{pager}.
The last lines name
the file. \\
\texttt{mv:}\allowbreak{}\texttt{\ cannot\ stat\ \textquotesingle{}ca.}\allowbreak{}\texttt{crt\textquotesingle{}}
& Chapter 7 moved
the files instead of
copying &
Re-\texttt{scp} from
the jumpbox (Chapter
4's control plane
distribution
command). \\
\texttt{etcdserver:}\allowbreak{}\texttt{\ request\ timed\ out}
in API server logs &
etcd down & On
\texttt{server}:
\texttt{etcdctl\ endpoint\ health}. \\
\texttt{curl:}\allowbreak{}\texttt{\ (60)\ SSL\ certificate\ problem}
& Wrong CA, or you
are not in the
directory containing
\texttt{ca.crt} &
\texttt{pwd} on the
jumpbox;
\texttt{openssl\ s\_client\ -}\allowbreak{}\texttt{connect\ server.}\allowbreak{}\texttt{kubernetes.}\allowbreak{}\texttt{local:}\allowbreak{}\texttt{6443\ -}\allowbreak{}\texttt{showcerts\ \textless{}/}\allowbreak{}\texttt{dev/}\allowbreak{}\texttt{null}. \\
\texttt{curl:}\allowbreak{}\texttt{\ (7)\ Failed\ to\ connect}
& API server not
listening, or
\texttt{/etc/hosts}
wrong &
\texttt{ssh\ root@server\ "ss\ -}\allowbreak{}\texttt{tlnp\ \textbackslash{}\textbar{}\ grep\ 6443"}. \\
Controller manager
restarting in a loop
& Bad kubeconfig, or
missing
\texttt{ca.key} &
\texttt{journalctl\ -}\allowbreak{}\texttt{u\ kube-}\allowbreak{}\texttt{controller-}\allowbreak{}\texttt{manager\ -}\allowbreak{}\texttt{n\ 50}.
\texttt{Unauthorized}
means credentials;
\texttt{no\ such\ file}
means a missing
\texttt{mv}. \\
Scheduler:
\texttt{no\ such\ file\ or\ directory}
on its config &
\texttt{kube-}\allowbreak{}\texttt{scheduler.}\allowbreak{}\texttt{yaml}
went to the wrong
directory &
\texttt{ls\ /}\allowbreak{}\texttt{etc/}\allowbreak{}\texttt{kubernetes/}\allowbreak{}\texttt{config/}\allowbreak{}\texttt{kube-}\allowbreak{}\texttt{scheduler.}\allowbreak{}\texttt{yaml}. \\
Later:
\texttt{kubectl\ logs}
returns
\texttt{Forbidden\ (user=kubernetes.}\allowbreak{}\texttt{.}\allowbreak{}\texttt{.}\allowbreak{}\texttt{)}
& The RBAC manifest
was never applied &
Re-run the
\texttt{kubectl\ apply}
step. \\
\end{longtable}

\hypertarget{check-your-understanding-8}{%
\section{Check your
understanding}\label{check-your-understanding-8}}

\textbf{Q1. The
ClusterRoleBinding
grants permissions
to a User named
\texttt{kubernetes}.
Where does that name
come from, and why a
User rather than a
ServiceAccount?}

From
\texttt{ca.conf}'s
\texttt{{[}kube-}\allowbreak{}\texttt{api-}\allowbreak{}\texttt{server\_distinguished\_name{]}}
section:
\texttt{CN\ =\ kubernetes}.
When the API server
connects to a
kubelet it presents
\texttt{kube-}\allowbreak{}\texttt{api-}\allowbreak{}\texttt{server.}\allowbreak{}\texttt{crt},
the kubelet forwards
a
SubjectAccessReview
naming the
certificate's CN as
the user, and RBAC
evaluates that name.
It cannot be a
ServiceAccount
because the API
server authenticates
by certificate, not
by a mounted token
--- it is an
external client to
the kubelet, not a
pod.

\textbf{Q2. You stop
\texttt{kube-scheduler}
and create a
Deployment. What
happens?}

The Deployment
object is created,
the controller
manager creates a
ReplicaSet, and the
ReplicaSet creates a
Pod --- all of which
succeed. The Pod
sits in
\texttt{Pending}
with no
\texttt{nodeName},
forever. Nothing
errors; nothing is
logged as a failure.
Restart the
scheduler and the
pod is bound and
started within a
second. This is the
clearest
demonstration in the
whole build that
Kubernetes
components are
independent
reconcilers against
shared state, not a
pipeline.

\textbf{Q3.
\texttt{-\/-bind-}\allowbreak{}\texttt{address=0.}\allowbreak{}\texttt{0.}\allowbreak{}\texttt{0.}\allowbreak{}\texttt{0}
means the API server
accepts connections
on every interface.
What actually
prevents anonymous
access?}

\texttt{-\/-client-}\allowbreak{}\texttt{ca-}\allowbreak{}\texttt{file}
combined with
\texttt{-\/-authorization-}\allowbreak{}\texttt{mode=Node,RBAC}.
Anyone can complete
a TCP connection and
a TLS handshake, and
anyone can read
unauthenticated
endpoints like
\texttt{/version}
--- which is exactly
what the
\texttt{curl}
verification does.
But any request
against the API
requires a client
certificate signed
by the cluster CA,
and the resulting
username must have
an RBAC grant.
Network reachability
and authorization
are separate layers
here, which is why
the \texttt{curl}
check works without
credentials while
\texttt{kubectl\ get\ nodes}
would not.

\hypertarget{chap-9}{%
\chapter{Bootstrapping
the Kubernetes
Worker
Nodes}\label{chap-9}}

\upstreamlab{docs/09-bootstrapping-kubernetes-workers.md}

\hypertarget{what-youre-building-9}{%
\section{What you're
building}\label{what-youre-building-9}}

On each of
\texttt{node-0} and
\texttt{node-1}: OS
dependencies, swap
disabled,
\texttt{runc}, CNI
plugins,
\texttt{containerd},
\texttt{kubelet},
and
\texttt{kube-proxy}
--- installed,
configured, and
running under
systemd. At the end,
both nodes appear as
\texttt{Ready} in
the API.

This is the longest
chapter and the one
with the most moving
parts. \textbf{You
run the entire
per-node section
twice}, once on each
worker.

\hypertarget{why-its-necessary-9}{%
\section{Why it's
necessary}\label{why-its-necessary-9}}

\begin{whybox}

Everything so far
has been control
plane. This is where
containers actually
run. Four layers
stack up on each
node, and each
depends on the one
below it:

\end{whybox}

\begin{outlst}
    Pod  ---------------------------------------------
      ^   kubelet tells containerd what to run
      |
  kubelet ---- watches API server for pods on this node
      ^   |    reports status back
      |   +--> CNI plugins: allocate pod IP, attach to cni0 bridge
      |
 containerd -- pulls images, manages snapshots
      ^
      |
    runc ----- creates the actual Linux namespaces and cgroups
\end{outlst}

\texttt{kube-proxy}
sits beside this
stack rather than in
it: it programs
iptables so Service
ClusterIPs work, and
it does not
participate in
running pods at all.

\hypertarget{key-concepts-9}{%
\section{Key
concepts}\label{key-concepts-9}}

\textbf{1. Swap is
disabled because pod
memory accounting
requires it.}
Kubernetes schedules
by memory request
and limit. With swap
active, a container
can exceed its
memory limit without
being killed --- it
just gets slower ---
which makes limits
unenforceable and
node pressure
signals meaningless.
Note the tension in
this build:
\texttt{kubelet-}\allowbreak{}\texttt{config.}\allowbreak{}\texttt{yaml}
sets
\texttt{failSwapOn:}\allowbreak{}\texttt{\ false}
(the kubelet will
start even with swap
present) and
\texttt{memorySwap.}\allowbreak{}\texttt{swapBehavior:}\allowbreak{}\texttt{\ NoSwap}
(containers may not
use it), yet the
tutorial still has
you run
\texttt{swapoff\ -a}.
Belt and braces.

\textbf{2.
\texttt{cgroupDriver:}\allowbreak{}\texttt{\ systemd}
must match
containerd.}
\texttt{kubelet-}\allowbreak{}\texttt{config.}\allowbreak{}\texttt{yaml}
sets
\texttt{cgroupDriver:}\allowbreak{}\texttt{\ systemd},
and
\texttt{containerd-}\allowbreak{}\texttt{config.}\allowbreak{}\texttt{toml}
sets
\texttt{SystemdCgroup\ =\ true}.
Both must agree.
When they disagree,
pods start and then
die in confusing
ways --- the two
components create
competing cgroup
hierarchies and
neither can account
for the other's
containers. This
mismatch is one of
the most common
real-world
Kubernetes node
failures.

\textbf{3.
\texttt{br-netfilter}
is what makes
Services work for
bridged traffic.}
Pods attach to the
\texttt{cni0} Linux
bridge. By default,
traffic that stays
within a bridge is
switched at L2 and
never traverses
\texttt{iptables}
--- which means
\texttt{kube-proxy}'s
Service rules would
be invisible to
pod-to-pod traffic
on the same node.
Loading
\texttt{br-netfilter}
and setting
\texttt{net.}\allowbreak{}\texttt{bridge.}\allowbreak{}\texttt{bridge-}\allowbreak{}\texttt{nf-}\allowbreak{}\texttt{call-}\allowbreak{}\texttt{iptables\ =\ 1}
forces bridged
frames through the
netfilter chains.
Without this,
ClusterIP Services
work from the node
but silently fail
from inside a pod on
the same node.

\textbf{4. The CNI
bridge config is
templated per node.}
\texttt{configs/}\allowbreak{}\texttt{10-}\allowbreak{}\texttt{bridge.}\allowbreak{}\texttt{conf}
contains the literal
token
\texttt{SUBNET},
which is replaced
with the node's pod
CIDR from
\texttt{machines.txt}:

\begin{outlst}
{
  "cniVersion": "1.0.0",
  "name": "bridge",
  "type": "bridge",
  "bridge": "cni0",
  "isGateway": true,
  "ipMasq": true,
  "ipam": {
    "type": "host-local",
    "ranges": [
      [{"subnet": "SUBNET"}]
    ],
    "routes": [{"dst": "0.0.0.0/0"}]
  }
}
\end{outlst}

\texttt{host-local}
IPAM means
\textbf{each node
allocates pod IPs
independently from
its own range}, with
no coordination.
This works only
because the ranges
are disjoint by
construction
(\texttt{10.200.0.0/24}
and
\texttt{10.200.1.0/24}).
\texttt{ipMasq:\ true}
SNATs pod traffic
leaving the node so
pods can reach the
internet.
\texttt{isGateway:\ true}
gives \texttt{cni0}
the first address in
the subnet, making
it the pods' default
gateway.

\textbf{5. The
kubelet config is
\emph{also} run
through \texttt{sed}
for \texttt{SUBNET},
but contains no such
token.} The
prerequisite loop
applies the same
substitution to
\texttt{configs/}\allowbreak{}\texttt{kubelet-}\allowbreak{}\texttt{config.}\allowbreak{}\texttt{yaml},
which as shipped has
no \texttt{SUBNET}
placeholder --- so
the \texttt{sed} is
a no-op and the file
is copied unchanged.
This is a harmless
leftover in the
upstream repo, not a
step you can skip:
the loop still needs
to produce
\texttt{kubelet-}\allowbreak{}\texttt{config.}\allowbreak{}\texttt{yaml}
in the current
directory for the
\texttt{scp} that
follows.

\textbf{6.
\texttt{maxPods:\ 16}
is unusually low.}
The Kubernetes
default is 110.
Sixteen suits a 2 GB
node and keeps each
node's \texttt{/24}
pod subnet
comfortably
oversized. Worth
knowing if you later
try to schedule more
than 32 pods and hit
a scheduling wall
that looks like a
resource problem.

\textbf{7. The
kubelet needs both a
server and a client
identity.}
\texttt{tlsCertFile}
/
\texttt{tlsPrivateKeyFile}
(\texttt{/}\allowbreak{}\texttt{var/}\allowbreak{}\texttt{lib/}\allowbreak{}\texttt{kubelet/}\allowbreak{}\texttt{kubelet.}\allowbreak{}\texttt{crt}
and \texttt{.key})
serve port 10250 for
inbound API server
connections.
\texttt{-\/-kubeconfig}
supplies the client
identity for
outbound calls.
\texttt{authentication.}\allowbreak{}\texttt{x509.}\allowbreak{}\texttt{clientCAFile}
verifies who is
calling in.
\texttt{authorization.}\allowbreak{}\texttt{mode:}\allowbreak{}\texttt{\ Webhook}
delegates every
inbound
authorization
decision back to the
API server --- which
is the mechanism
Chapter 8's RBAC
ClusterRole
satisfies.

\hypertarget{the-work-8}{%
\section{The
work}\label{the-work-8}}

\textbf{Prerequisites
--- run on the
jumpbox.}

Template and push
the per-node config
files:

\begin{cmdlst}
# on: jumpbox
for HOST in node-0 node-1; do
  SUBNET=$(grep ${HOST} machines.txt | cut -d " " -f 4)
  sed "s|SUBNET|$SUBNET|g" \
    configs/10-bridge.conf > 10-bridge.conf

  sed "s|SUBNET|$SUBNET|g" \
    configs/kubelet-config.yaml > kubelet-config.yaml

  scp 10-bridge.conf kubelet-config.yaml \
  root@${HOST}:~/
done
\end{cmdlst}

This is where
\texttt{machines.txt}
formatting matters
most.
\texttt{cut\ -d\ "\ "\ -f\ 4}
splits on a single
literal space; tabs
or double spaces
make \texttt{SUBNET}
empty, the
substitution
produces
\texttt{"subnet":\ ""},
and CNI fails at pod
creation with an
IPAM error --- three
chapters later.

Push the binaries,
configs, and units:

\begin{cmdlst}
# on: jumpbox
for HOST in node-0 node-1; do
  scp \
    downloads/worker/* \
    downloads/client/kubectl \
    configs/99-loopback.conf \
    configs/containerd-config.toml \
    configs/kube-proxy-config.yaml \
    units/containerd.service \
    units/kubelet.service \
    units/kube-proxy.service \
    root@${HOST}:~/
done
\end{cmdlst}

Push the CNI
plugins:

\begin{cmdlst}
# on: jumpbox
for HOST in node-0 node-1; do
  scp \
    downloads/cni-plugins/* \
    root@${HOST}:~/cni-plugins/
done
\end{cmdlst}

\begin{beyondbox}\calloutlabel{teal}{BEYOND THE SOURCE}

this \texttt{scp}
copies many files
into
\texttt{\textasciitilde{}/cni-plugins/},
which must already
exist ---
\texttt{scp} with
multiple sources
requires a directory
destination. If it
fails with
\texttt{Not\ a\ directory}
or
\texttt{No\ such\ file\ or\ directory},
run
\texttt{for\ HOST\ in\ node-}\allowbreak{}\texttt{0\ node-}\allowbreak{}\texttt{1;\ do\ ssh\ -}\allowbreak{}\texttt{n\ root@\$\{HOST\}\ mkdir\ -}\allowbreak{}\texttt{p\ cni-}\allowbreak{}\texttt{plugins;\ done}
first, then repeat
the \texttt{scp}.

\end{beyondbox}

\textbf{Per-node
provisioning --- run
on \texttt{node-0},
then repeat every
step on
\texttt{node-1}.}

\begin{cmdlst}
# on: jumpbox
ssh root@node-0
\end{cmdlst}

Install OS
dependencies:

\begin{cmdlst}
# on: node-0 (then repeat on node-1)
{
  apt-get update
  apt-get -y install socat conntrack ipset kmod
}
\end{cmdlst}

Each of these has a
specific job:
\texttt{socat} backs
\texttt{kubectl\ port-}\allowbreak{}\texttt{forward};
\texttt{conntrack}
lets
\texttt{kube-proxy}
manage connection
tracking entries;
\texttt{ipset}
supports efficient
iptables rule sets;
\texttt{kmod}
provides
\texttt{modprobe},
needed by the
containerd unit's
\texttt{ExecStartPre}.

Check and disable
swap:

\begin{cmdlst}
# on: node-0 (then repeat on node-1)
swapon --show
\end{cmdlst}

If output is empty
then swap is
disabled. If swap is
enabled:

\begin{cmdlst}
# on: node-0 (then repeat on node-1)
swapoff -a
\end{cmdlst}

\begin{notebox}\calloutlabel{steel}{FROM THE SOURCE}

To ensure swap
remains off after
reboot consult your
Linux distro
documentation.

\end{notebox}

On Debian 12 that
usually means
commenting out the
swap line in
\texttt{/etc/fstab}.

Create the
installation
directories:

\begin{cmdlst}
# on: node-0 (then repeat on node-1)
mkdir -p \
  /etc/cni/net.d \
  /opt/cni/bin \
  /var/lib/kubelet \
  /var/lib/kube-proxy \
  /var/lib/kubernetes \
  /var/run/kubernetes
\end{cmdlst}

\texttt{/etc/cni/net.d}
and
\texttt{/opt/cni/bin}
are the conventional
CNI paths, and they
must match
\texttt{containerd-}\allowbreak{}\texttt{config.}\allowbreak{}\texttt{toml}'s
\texttt{conf\_dir}
and
\texttt{bin\_dir}.
\texttt{/}\allowbreak{}\texttt{var/}\allowbreak{}\texttt{lib/}\allowbreak{}\texttt{kubelet}
and
\texttt{/}\allowbreak{}\texttt{var/}\allowbreak{}\texttt{lib/}\allowbreak{}\texttt{kube-}\allowbreak{}\texttt{proxy}
already exist from
Chapters 4 and 5;
\texttt{mkdir\ -p}
makes that harmless.

Install the worker
binaries:

\begin{cmdlst}
# on: node-0 (then repeat on node-1)
{
  mv crictl kube-proxy kubelet runc \
    /usr/local/bin/
  mv containerd containerd-shim-runc-v2 containerd-stress /bin/
  mv cni-plugins/* /opt/cni/bin/
}
\end{cmdlst}

Note containerd's
binaries go to
\texttt{/bin}, not
\texttt{/usr/local/bin}
--- because
\texttt{containerd.}\allowbreak{}\texttt{service}
specifies
\texttt{ExecStart=/}\allowbreak{}\texttt{bin/}\allowbreak{}\texttt{containerd}
literally. Put it
elsewhere and the
unit fails to start
with
\texttt{No\ such\ file\ or\ directory}.

\textbf{Configure
CNI networking:}

\begin{cmdlst}
# on: node-0 (then repeat on node-1)
mv 10-bridge.conf 99-loopback.conf /etc/cni/net.d/
\end{cmdlst}

The numeric prefixes
set evaluation
order:
\texttt{10-bridge.conf}
before
\texttt{99-}\allowbreak{}\texttt{loopback.}\allowbreak{}\texttt{conf}.

Load
\texttt{br-netfilter}:

\begin{cmdlst}
# on: node-0 (then repeat on node-1)
{
  modprobe br-netfilter
  echo "br-netfilter" >> /etc/modules-load.d/modules.conf
}
\end{cmdlst}

First line loads it
now; second makes it
persist across
reboot.

\begin{cmdlst}
# on: node-0 (then repeat on node-1)
{
  echo "net.bridge.bridge-nf-call-iptables = 1" \
    >> /etc/sysctl.d/kubernetes.conf
  echo "net.bridge.bridge-nf-call-ip6tables = 1" \
    >> /etc/sysctl.d/kubernetes.conf
  sysctl -p /etc/sysctl.d/kubernetes.conf
}
\end{cmdlst}

These sysctls only
exist once
\texttt{br-netfilter}
is loaded --- which
is why the
\texttt{modprobe}
comes first.
Reversing the order
gives
\texttt{sysctl:}\allowbreak{}\texttt{\ cannot\ stat\ /}\allowbreak{}\texttt{proc/}\allowbreak{}\texttt{sys/}\allowbreak{}\texttt{net/}\allowbreak{}\texttt{bridge/}\allowbreak{}\texttt{bridge-}\allowbreak{}\texttt{nf-}\allowbreak{}\texttt{call-}\allowbreak{}\texttt{iptables}.

\textbf{Configure
containerd:}

\begin{cmdlst}
# on: node-0 (then repeat on node-1)
{
  mkdir -p /etc/containerd/
  mv containerd-config.toml /etc/containerd/config.toml
  mv containerd.service /etc/systemd/system/
}
\end{cmdlst}

Note the rename to
\texttt{config.toml}
--- containerd looks
for that exact path.

\textbf{Configure
the kubelet:}

\begin{cmdlst}
# on: node-0 (then repeat on node-1)
{
  mv kubelet-config.yaml /var/lib/kubelet/
  mv kubelet.service /etc/systemd/system/
}
\end{cmdlst}

\textbf{Configure
kube-proxy:}

\begin{cmdlst}
# on: node-0 (then repeat on node-1)
{
  mv kube-proxy-config.yaml /var/lib/kube-proxy/
  mv kube-proxy.service /etc/systemd/system/
}
\end{cmdlst}

\textbf{Start the
services:}

\begin{cmdlst}
# on: node-0 (then repeat on node-1)
{
  systemctl daemon-reload
  systemctl enable containerd kubelet kube-proxy
  systemctl start containerd kubelet kube-proxy
}
\end{cmdlst}

\texttt{kubelet.service}
declares
\texttt{Requires=containerd.}\allowbreak{}\texttt{service}
and
\texttt{After=containerd.}\allowbreak{}\texttt{service},
so systemd
guarantees ordering
regardless of the
order on this
command line.

Check the kubelet:

\begin{cmdlst}
# on: node-0 (then repeat on node-1)
systemctl is-active kubelet
\end{cmdlst}

\begin{outlst}
active
\end{outlst}

\textbf{Be sure to
complete every step
in this section on
both \texttt{node-0}
and \texttt{node-1}
before continuing.}

\hypertarget{verification-8}{%
\section{Verification}\label{verification-8}}

From the jumpbox:

\begin{cmdlst}
# on: jumpbox
ssh root@server \
  "kubectl get nodes \
  --kubeconfig admin.kubeconfig"
\end{cmdlst}

\begin{outlst}
NAME     STATUS   ROLES    AGE    VERSION
node-0   Ready    <none>   1m     v1.32.3
node-1   Ready    <none>   10s    v1.32.3
\end{outlst}

\texttt{ROLES} shows
\texttt{\textless{}none\textgreater{}}
because role labels
are a convention
applied by
installers, not
something Kubernetes
sets. \texttt{Ready}
means the kubelet
registered, passed
its own health
checks, and found a
usable CNI
configuration --- a
node with no CNI
config reports
\texttt{NotReady}
with
\texttt{network\ plugin\ is\ not\ ready}.

\begin{beyondbox}\calloutlabel{teal}{BEYOND THE SOURCE}

on each node,
\texttt{ip\ addr\ show\ cni0}
will show nothing
until the first pod
is scheduled --- the
bridge is created on
demand by the CNI
plugin, not at
kubelet startup. Its
absence now is not a
fault.

\end{beyondbox}

\hypertarget{when-it-goes-wrong-8}{%
\section{When it
goes
wrong}\label{when-it-goes-wrong-8}}

\begin{longtable}[]{@{}
  >{\raggedright\arraybackslash}p{(\columnwidth - 4\tabcolsep) * \real{0.3200}}
  >{\raggedright\arraybackslash}p{(\columnwidth - 4\tabcolsep) * \real{0.2400}}
  >{\raggedright\arraybackslash}p{(\columnwidth - 4\tabcolsep) * \real{0.4400}}@{}}
\toprule\noalign{}
\begin{minipage}[b]{\linewidth}\raggedright
Symptom
\end{minipage} &
\begin{minipage}[b]{\linewidth}\raggedright
Cause
\end{minipage} &
\begin{minipage}[b]{\linewidth}\raggedright
Diagnosis
\end{minipage} \\
\midrule\noalign{}
\endhead
\bottomrule\noalign{}
\endlastfoot
Node never appears
in
\texttt{kubectl\ get\ nodes}
& kubelet cannot
reach or
authenticate to the
API server &
\texttt{journalctl\ -}\allowbreak{}\texttt{u\ kubelet\ -}\allowbreak{}\texttt{n\ 50}
on the node.
\texttt{no\ such\ host}
\textrightarrow{}
\texttt{/etc/hosts}.
\texttt{Unauthorized}
\textrightarrow{}
kubeconfig.
\texttt{connection\ refused}
\textrightarrow{}
API server. \\
Node appears but
\texttt{NotReady} &
No CNI config found
&
\texttt{ls\ /}\allowbreak{}\texttt{etc/}\allowbreak{}\texttt{cni/}\allowbreak{}\texttt{net.}\allowbreak{}\texttt{d/}\allowbreak{}
should show both
files.
\texttt{journalctl\ -}\allowbreak{}\texttt{u\ kubelet\ \textbackslash{}\textbar{}\ grep\ -}\allowbreak{}\texttt{i\ cni}. \\
\texttt{NotReady},
kubelet logs a CNI
version or plugin
error &
\texttt{10-bridge.conf}
has an empty or
malformed subnet &
\texttt{cat\ /}\allowbreak{}\texttt{etc/}\allowbreak{}\texttt{cni/}\allowbreak{}\texttt{net.}\allowbreak{}\texttt{d/}\allowbreak{}\texttt{10-}\allowbreak{}\texttt{bridge.}\allowbreak{}\texttt{conf}
--- the
\texttt{subnet}
field must contain a
real CIDR. Fix on
the jumpbox and
re-push. \\
kubelet:
\texttt{Exec\ format\ error}
& Architecture
mismatch between
jumpbox and node &
\texttt{file\ /}\allowbreak{}\texttt{usr/}\allowbreak{}\texttt{local/}\allowbreak{}\texttt{bin/}\allowbreak{}\texttt{kubelet}
on the node. Rebuild
with matching
architectures. \\
containerd fails:
\texttt{/}\allowbreak{}\texttt{bin/}\allowbreak{}\texttt{containerd:}\allowbreak{}\texttt{\ no\ such\ file}
& Binary installed
to
\texttt{/usr/local/bin}
instead of
\texttt{/bin} &
\texttt{mv\ /}\allowbreak{}\texttt{usr/}\allowbreak{}\texttt{local/}\allowbreak{}\texttt{bin/}\allowbreak{}\texttt{containerd\ /}\allowbreak{}\texttt{bin/}\allowbreak{},
then restart. \\
Pods stuck in
\texttt{ContainerCreating}
& CNI plugin
binaries missing
from
\texttt{/opt/cni/bin}
&
\texttt{ls\ /}\allowbreak{}\texttt{opt/}\allowbreak{}\texttt{cni/}\allowbreak{}\texttt{bin/}\allowbreak{}
should list
\textasciitilde a
dozen binaries
including
\texttt{bridge},
\texttt{host-local},
\texttt{loopback}.
Re-run the CNI push
and the
\texttt{mv}. \\
Pods run but
ClusterIP
unreachable
\emph{from inside a
pod} &
\texttt{br-netfilter}
not loaded or
sysctls not applied
&
\texttt{lsmod\ \textbackslash{}\textbar{}\ grep\ br\_netfilter};
\texttt{sysctl\ net.}\allowbreak{}\texttt{bridge.}\allowbreak{}\texttt{bridge-}\allowbreak{}\texttt{nf-}\allowbreak{}\texttt{call-}\allowbreak{}\texttt{iptables}
must be
\texttt{1}. \\
kubelet:
\texttt{failed\ to\ run\ Kubelet:}\allowbreak{}\texttt{\ running\ with\ swap\ on}
& Swap enabled and
\texttt{failSwapOn}
not honoured &
\texttt{swapoff\ -a},
then
\texttt{systemctl\ restart\ kubelet}. \\
Node \texttt{Ready}
then flaps to
\texttt{NotReady} &
Node OOM --- the API
server or kubelet
being killed &
\texttt{dmesg\ \textbackslash{}\textbar{}\ grep\ -}\allowbreak{}\texttt{i\ oom}
on the node; check
RAM against the
Chapter 1 table. \\
\end{longtable}

\hypertarget{check-your-understanding-9}{%
\section{Check your
understanding}\label{check-your-understanding-9}}

\textbf{Q1. The
kubelet reads a
kubeconfig at
\texttt{/}\allowbreak{}\texttt{var/}\allowbreak{}\texttt{lib/}\allowbreak{}\texttt{kubelet/}\allowbreak{}\texttt{kubeconfig}
and a config file at
\texttt{/}\allowbreak{}\texttt{var/}\allowbreak{}\texttt{lib/}\allowbreak{}\texttt{kubelet/}\allowbreak{}\texttt{kubelet-}\allowbreak{}\texttt{config.}\allowbreak{}\texttt{yaml}.
What does each one
govern?}

The kubeconfig is
the kubelet's
\textbf{outbound
identity} --- where
the API server is,
which CA to trust,
and which client
certificate to
present when
watching for pods
and reporting
status.
\texttt{kubelet-}\allowbreak{}\texttt{config.}\allowbreak{}\texttt{yaml}
is the kubelet's
\textbf{own
behaviour and
inbound security}
--- cgroup driver,
max pods, the
container runtime
endpoint, and how to
authenticate and
authorize clients
connecting to port
10250. Outbound
versus inbound.

\textbf{Q2. Both
nodes run identical
\texttt{kube-}\allowbreak{}\texttt{proxy-}\allowbreak{}\texttt{config.}\allowbreak{}\texttt{yaml}
with
\texttt{clusterCIDR:}\allowbreak{}\texttt{\ "10.}\allowbreak{}\texttt{200.}\allowbreak{}\texttt{0.}\allowbreak{}\texttt{0/}\allowbreak{}\texttt{16"},
but each has a
different
\texttt{10-bridge.conf}
subnet. Why the
asymmetry?}

\texttt{kube-proxy}
needs to know the
whole pod CIDR so it
can distinguish pod
traffic from
external traffic
when deciding what
to masquerade ---
that knowledge is
cluster-wide and
identical
everywhere. CNI IPAM
needs to know which
slice \emph{this
node} may allocate
from, and that must
be unique per node
or two nodes would
hand out the same
pod IPs. One is a
global fact, the
other is a local
allocation.

\textbf{Q3. You
schedule a pod on
\texttt{node-0}. It
gets
\texttt{10.200.0.5}
and can reach the
internet, but cannot
ping a pod on
\texttt{node-1}. Is
the CNI
configuration
wrong?}

No --- this is the
expected state at
the end of Chapter
9, and it is exactly
what Chapter 11
fixes. The bridge
gives pods on the
same node L2
connectivity to each
other, and
\texttt{ipMasq:\ true}
gives them outbound
internet access via
the node's own IP.
Nothing yet tells
\texttt{node-0}'s
kernel that
\texttt{10.200.1.0/24}
is reachable via
\texttt{node-1}'s
address. That is a
routing gap, not a
CNI gap.

\hypertarget{chap-10}{%
\chapter{Configuring
kubectl for Remote
Access}\label{chap-10}}

\upstreamlab{docs/10-configuring-kubectl.md}

\hypertarget{what-youre-building-10}{%
\section{What you're
building}\label{what-youre-building-10}}

A kubeconfig at the
default location
\texttt{\textasciitilde{}/.kube/config}
on the jumpbox,
authenticating as
\texttt{admin}
against
\texttt{https:}\allowbreak{}\texttt{/}\allowbreak{}\texttt{/}\allowbreak{}\texttt{server.}\allowbreak{}\texttt{kubernetes.}\allowbreak{}\texttt{local:}\allowbreak{}\texttt{6443}.

\hypertarget{why-its-necessary-10}{%
\section{Why it's
necessary}\label{why-its-necessary-10}}

\begin{whybox}

Until now, every
\texttt{kubectl}
invocation has
required either
\texttt{ssh\ root@server}
or an explicit
\texttt{-\/-kubeconfig\ admin.}\allowbreak{}\texttt{kubeconfig}.
This chapter makes
the jumpbox a proper
administrative
client:
\texttt{kubectl\ get\ nodes}
just works, from the
machine where you
have been doing
everything else.

\end{whybox}

The difference from
the
\texttt{admin.}\allowbreak{}\texttt{kubeconfig}
built in Chapter 5
is one field. That
file points at
\texttt{127.0.0.1}
and is used
\emph{on} the
control plane host.
This one points at
\texttt{server.}\allowbreak{}\texttt{kubernetes.}\allowbreak{}\texttt{local}
and is used
\emph{from} the
jumpbox. Same
identity, different
network path.

\hypertarget{key-concepts-10}{%
\section{Key
concepts}\label{key-concepts-10}}

\textbf{1.
\texttt{\textasciitilde{}/.kube/config}
is a well-known
default.}
\texttt{kubectl}
looks there when no
\texttt{-\/-kubeconfig}
flag and no
\texttt{KUBECONFIG}
environment variable
are set. Writing
there is what makes
the flag unnecessary
from now on.

\textbf{2. No
\texttt{-\/-kubeconfig}
flag and no
\texttt{-\/-embed-certs}
here.} Both
omissions are
deliberate. Without
\texttt{-\/-kubeconfig},
\texttt{kubectl\ config}
writes to the
default location ---
that is the whole
point. Without
\texttt{-\/-embed-certs},
the file stores
\emph{file paths} to
\texttt{admin.crt}
and
\texttt{admin.key}
rather than copies
of them. That means
\textbf{this
kubeconfig only
works from the
directory containing
those files}, or
after you edit the
paths to be
absolute. Running
\texttt{kubectl\ get\ nodes}
from \texttt{/root}
instead of
\texttt{/}\allowbreak{}\texttt{root/}\allowbreak{}\texttt{kubernetes-}\allowbreak{}\texttt{the-}\allowbreak{}\texttt{hard-}\allowbreak{}\texttt{way}
fails with a
file-not-found error
on the certificate.

\textbf{3. The
context is named,
not
\texttt{default}.}
Chapter 5 used
\texttt{set-}\allowbreak{}\texttt{context\ default};
here the context is
\texttt{kubernetes-}\allowbreak{}\texttt{the-}\allowbreak{}\texttt{hard-}\allowbreak{}\texttt{way}.
A cosmetic
difference, but it
is why
\texttt{kubectl\ config\ current-}\allowbreak{}\texttt{context}
prints something
different on the
jumpbox than on
\texttt{server}.

\hypertarget{the-work-9}{%
\section{The
work}\label{the-work-9}}

All commands run on
the
\textbf{jumpbox},
from
\texttt{/}\allowbreak{}\texttt{root/}\allowbreak{}\texttt{kubernetes-}\allowbreak{}\texttt{the-}\allowbreak{}\texttt{hard-}\allowbreak{}\texttt{way}.

First confirm the
API server is
reachable by name:

\begin{cmdlst}
# on: jumpbox
curl --cacert ca.crt \
  https://server.kubernetes.local:6443/version
\end{cmdlst}

\begin{outlst}
{
  "major": "1",
  "minor": "32",
  "gitVersion": "v1.32.3",
  "gitCommit": "32cc146f75aad04beaaa245a7157eb35063a9f99",
  "gitTreeState": "clean",
  "buildDate": "2025-03-11T19:52:21Z",
  "goVersion": "go1.23.6",
  "compiler": "gc",
  "platform": "linux/arm64"
}
\end{outlst}

Generate the
kubeconfig:

\begin{cmdlst}
# on: jumpbox
{
  kubectl config set-cluster kubernetes-the-hard-way \
    --certificate-authority=ca.crt \
    --embed-certs=true \
    --server=https://server.kubernetes.local:6443

  kubectl config set-credentials admin \
    --client-certificate=admin.crt \
    --client-key=admin.key

  kubectl config set-context kubernetes-the-hard-way \
    --cluster=kubernetes-the-hard-way \
    --user=admin

  kubectl config use-context kubernetes-the-hard-way
}
\end{cmdlst}

Note the asymmetry:
the \textbf{CA is
embedded}
(\texttt{-\/-embed-}\allowbreak{}\texttt{certs=true}
on
\texttt{set-cluster})
but the
\textbf{client
certificate and key
are referenced by
path} (no
\texttt{-\/-embed-certs}
on
\texttt{set-credentials}).
This is the
directory dependency
described above.

The results of
running the command
above should create
a kubeconfig file in
the default location
\texttt{\textasciitilde{}/.kube/config}
used by the
\texttt{kubectl}
commandline tool.
This also means you
can run the
\texttt{kubectl}
command without
specifying a config.

\hypertarget{verification-9}{%
\section{Verification}\label{verification-9}}

\begin{cmdlst}
# on: jumpbox
kubectl version
\end{cmdlst}

\begin{outlst}
Client Version: v1.32.3
Kustomize Version: v5.5.0
Server Version: v1.32.3
\end{outlst}

The presence of
\texttt{Server\ Version}
is the proof --- it
required an
authenticated round
trip. In Chapter 2
the same command
with
\texttt{-\/-client}
printed only the
first two lines.

\begin{cmdlst}
# on: jumpbox
kubectl get nodes
\end{cmdlst}

\begin{outlst}
NAME     STATUS   ROLES    AGE    VERSION
node-0   Ready    <none>   10m   v1.32.3
node-1   Ready    <none>   10m   v1.32.3
\end{outlst}

\hypertarget{when-it-goes-wrong-9}{%
\section{When it
goes
wrong}\label{when-it-goes-wrong-9}}

\begin{longtable}[]{@{}
  >{\raggedright\arraybackslash}p{(\columnwidth - 4\tabcolsep) * \real{0.3200}}
  >{\raggedright\arraybackslash}p{(\columnwidth - 4\tabcolsep) * \real{0.2400}}
  >{\raggedright\arraybackslash}p{(\columnwidth - 4\tabcolsep) * \real{0.4400}}@{}}
\toprule\noalign{}
\begin{minipage}[b]{\linewidth}\raggedright
Symptom
\end{minipage} &
\begin{minipage}[b]{\linewidth}\raggedright
Cause
\end{minipage} &
\begin{minipage}[b]{\linewidth}\raggedright
Diagnosis
\end{minipage} \\
\midrule\noalign{}
\endhead
\bottomrule\noalign{}
\endlastfoot
\texttt{open\ admin.}\allowbreak{}\texttt{crt:}\allowbreak{}\texttt{\ no\ such\ file\ or\ directory}
& Running from a
different directory;
the paths are
relative &
\texttt{cd\ /}\allowbreak{}\texttt{root/}\allowbreak{}\texttt{kubernetes-}\allowbreak{}\texttt{the-}\allowbreak{}\texttt{hard-}\allowbreak{}\texttt{way},
or re-run
\texttt{set-credentials}
with absolute
paths. \\
\texttt{The\ connection\ to\ the\ server\ .}\allowbreak{}\texttt{.}\allowbreak{}\texttt{.}\allowbreak{}\texttt{\ was\ refused}
& API server down,
or name resolves to
the wrong host &
\texttt{curl\ -\/-cacert\ ca.}\allowbreak{}\texttt{crt\ https:}\allowbreak{}\texttt{/}\allowbreak{}\texttt{/}\allowbreak{}\texttt{server.}\allowbreak{}\texttt{kubernetes.}\allowbreak{}\texttt{local:}\allowbreak{}\texttt{6443/}\allowbreak{}\texttt{version};
\texttt{getent\ hosts\ server.}\allowbreak{}\texttt{kubernetes.}\allowbreak{}\texttt{local}. \\
\texttt{x509:}\allowbreak{}\texttt{\ certificate\ signed\ by\ unknown\ authority}
& Embedded CA does
not match the API
server's issuer &
Re-run this
chapter's block;
confirm you are
using the same
\texttt{ca.crt} that
signed
\texttt{kube-}\allowbreak{}\texttt{api-}\allowbreak{}\texttt{server.}\allowbreak{}\texttt{crt}. \\
\texttt{error:}\allowbreak{}\texttt{\ You\ must\ be\ logged\ in\ to\ the\ server\ (Unauthorized)}
& \texttt{admin.crt}
not presented, or
not signed by the
cluster CA &
\texttt{openssl\ verify\ -}\allowbreak{}\texttt{CAfile\ ca.}\allowbreak{}\texttt{crt\ admin.}\allowbreak{}\texttt{crt}. \\
\texttt{Server\ Version}
missing from
\texttt{kubectl\ version}
& The server round
trip failed but the
command exited 0 &
Check for a warning
line above the
output; run
\texttt{kubectl\ get\ nodes}
for a clearer
error. \\
\texttt{error:}\allowbreak{}\texttt{\ context\ "kubernetes-}\allowbreak{}\texttt{the-}\allowbreak{}\texttt{hard-}\allowbreak{}\texttt{way"\ does\ not\ exist}
& The block was run
partially & Re-run
all four commands
together. \\
\end{longtable}

\hypertarget{check-your-understanding-10}{%
\section{Check your
understanding}\label{check-your-understanding-10}}

\textbf{Q1. Why can
\texttt{kubectl\ version\ -\/-client}
succeed in Chapter 2
but
\texttt{kubectl\ version}
fail here?}

\texttt{-\/-client}
prints the binary's
compiled-in version
and makes no network
call. Without it,
\texttt{kubectl}
queries
\texttt{/version} on
the API server using
the active context's
credentials. The
first is a local
fact; the second is
an end-to-end test
of name resolution,
TLS, authentication,
and authorization.

\textbf{Q2. You copy
\texttt{\textasciitilde{}/.kube/config}
to your laptop. Does
it work?}

No.~The
\texttt{client-}\allowbreak{}\texttt{certificate}
and
\texttt{client-key}
entries are relative
file paths that
exist only on the
jumpbox. To make it
portable, regenerate
with
\texttt{-\/-embed-}\allowbreak{}\texttt{certs=true}
on
\texttt{set-credentials}
--- which also means
the resulting file
contains the
\texttt{admin}
private key and
grants cluster-admin
to anyone who
obtains it. That is
the trade this
chapter makes in the
safe direction,
whether deliberately
or not.

\textbf{Q3. Both
\texttt{\textasciitilde{}/.kube/config}
on the jumpbox and
\texttt{admin.}\allowbreak{}\texttt{kubeconfig}
on \texttt{server}
authenticate as
\texttt{admin}.
Could you use just
one?}

Functionally yes,
but each only works
where it points.
\texttt{admin.}\allowbreak{}\texttt{kubeconfig}
targets
\texttt{127.0.0.1},
so it works only on
\texttt{server};
\texttt{\textasciitilde{}/.kube/config}
targets
\texttt{server.}\allowbreak{}\texttt{kubernetes.}\allowbreak{}\texttt{local},
so it works from
anywhere that
resolves the name.
Keeping both means a
working
\texttt{kubectl} on
the control plane
host even if name
resolution breaks
--- which is exactly
when you most need
to inspect the
cluster.

\hypertarget{chap-11}{%
\chapter{Provisioning
Pod Network
Routes}\label{chap-11}}

\upstreamlab{docs/11-pod-network-routes.md}

\hypertarget{what-youre-building-11}{%
\section{What you're
building}\label{what-youre-building-11}}

Four static routes,
so that traffic
destined for a pod
subnet is forwarded
to the node that
owns it: two on
\texttt{server}, one
on \texttt{node-0},
one on
\texttt{node-1}.

\hypertarget{why-its-necessary-11}{%
\section{Why it's
necessary}\label{why-its-necessary-11}}

\begin{whybox}

Kubernetes mandates
that \textbf{every
pod can reach every
other pod without
NAT}, on a flat
address space. It
does not implement
that requirement ---
it delegates it.
Chapter 9 gave each
node a bridge and a
slice of the pod
CIDR, which produces
working pod-to-pod
connectivity
\emph{within} a
node. Nothing yet
tells any kernel how
to reach the other
node's slice.

\end{whybox}

Without these
routes, a packet
from
\texttt{10.200.0.5}
on \texttt{node-0}
to
\texttt{10.200.1.7}
on \texttt{node-1}
matches no specific
route, falls through
to the default
gateway, and is
dropped by your LAN
router as an
unroutable
destination.
Symptoms: cross-node
pod-to-pod traffic
times out; Services
with endpoints on
the other node
intermittently fail
depending on which
endpoint kube-proxy
selects; a
two-replica
Deployment appears
half-broken.

This is also why
\texttt{server}
needs routes even
though it runs no
pods: the API server
initiates
connections to
kubelets, and some
traffic (including
\texttt{kubectl\ exec}
and
\texttt{port-forward}
paths) originates
there.

\hypertarget{key-concepts-11}{%
\section{Key
concepts}\label{key-concepts-11}}

\textbf{1. This is
the simplest
possible
implementation of
the Kubernetes
network model ---
and its limits are
the reason overlays
exist.} Static L3
routes work here
because all four
machines share one
L2 segment and can
forward directly to
each other. They do
not scale: \emph{n}
nodes need
\emph{n}(\emph{n}\textminus{}1)
routes, every new
node requires
touching every
existing node, and
none of it survives
a node's IP
changing. Production
clusters run a CNI
plugin (Calico,
Cilium, Flannel)
that does this
automatically ---
via BGP, a VXLAN
overlay, or eBPF ---
which is precisely
the job this chapter
does by hand so you
can see what is
being automated.

\textbf{2. The
routes are not
persistent.}
\texttt{ip\ route\ add}
writes to the
running kernel's
routing table only.
\textbf{Reboot any
machine and its
routes are gone},
and cross-node pod
networking silently
breaks. This is the
single most likely
thing to bite you if
you return to this
lab the next day.

\textbf{3. Each node
needs routes only to
the \emph{other}
nodes.}
\texttt{node-0}
needs a route to
\texttt{10.200.1.0/24}
but not to
\texttt{10.200.0.0/24}
--- its own subnet
is already directly
connected via
\texttt{cni0}.
Adding a route for
your own subnet
would shadow the
local one and break
pod networking on
that node entirely.

\textbf{4.
\texttt{via} means
next hop, and the
next hop is the
node's LAN address.}
\texttt{ip\ route\ add\ 10.}\allowbreak{}\texttt{200.}\allowbreak{}\texttt{1.}\allowbreak{}\texttt{0/}\allowbreak{}\texttt{24\ via\ \textless{}node-}\allowbreak{}\texttt{1\textquotesingle{}s\ LAN\ IP\textgreater{}}
says: to reach that
pod subnet, hand the
packet to
\texttt{node-1} on
the local network.
\texttt{node-1}'s
kernel then forwards
to \texttt{cni0} and
on to the pod.

\hypertarget{the-work-10}{%
\section{The
work}\label{the-work-10}}

All commands run
from the
\textbf{jumpbox}.

Gather addresses
from the machine
database:

\begin{cmdlst}
# on: jumpbox
{
  SERVER_IP=$(grep server machines.txt | cut -d " " -f 1)
  NODE_0_IP=$(grep node-0 machines.txt | cut -d " " -f 1)
  NODE_0_SUBNET=$(grep node-0 machines.txt | cut -d " " -f 4)
  NODE_1_IP=$(grep node-1 machines.txt | cut -d " " -f 1)
  NODE_1_SUBNET=$(grep node-1 machines.txt | cut -d " " -f 4)
}
\end{cmdlst}

These must be set in
the \textbf{same
shell} that runs the
commands below ---
the heredocs expand
them locally before
sending. If you open
a new terminal,
re-run this block
first.

\begin{beyondbox}\calloutlabel{teal}{BEYOND THE SOURCE}

\texttt{echo\ "\$NODE\_0\_IP\ \$NODE\_0\_SUBNET\ /}\allowbreak{}\texttt{\ \$NODE\_1\_IP\ \$NODE\_1\_SUBNET"}
before proceeding.
Empty subnets mean
\texttt{machines.txt}
has tabs or multiple
spaces; fix that
file rather than
working around it
here.

\end{beyondbox}

Add routes on
\texttt{server} for
both pod subnets:

\begin{cmdlst}
# on: jumpbox
ssh root@server <<EOF
  ip route add ${NODE_0_SUBNET} via ${NODE_0_IP}
  ip route add ${NODE_1_SUBNET} via ${NODE_1_IP}
EOF
\end{cmdlst}

The unquoted
\texttt{EOF} is
deliberate --- it
allows the shell to
substitute the
variables before the
text reaches
\texttt{ssh}.
Quoting it
(\texttt{\textless{}\textless{}\textquotesingle{}EOF\textquotesingle{}})
would send the
literal
\texttt{\$\{NODE\_0\_SUBNET\}}
text and the command
would fail on the
remote side.

Add a route on
\texttt{node-0} to
reach
\texttt{node-1}'s
pods:

\begin{cmdlst}
# on: jumpbox
ssh root@node-0 <<EOF
  ip route add ${NODE_1_SUBNET} via ${NODE_1_IP}
EOF
\end{cmdlst}

And on
\texttt{node-1} to
reach
\texttt{node-0}'s
pods:

\begin{cmdlst}
# on: jumpbox
ssh root@node-1 <<EOF
  ip route add ${NODE_0_SUBNET} via ${NODE_0_IP}
EOF
\end{cmdlst}

\hypertarget{verification-10}{%
\section{Verification}\label{verification-10}}

\begin{cmdlst}
# on: jumpbox
ssh root@server ip route
\end{cmdlst}

\begin{outlst}
default via XXX.XXX.XXX.XXX dev ens160 
10.200.0.0/24 via XXX.XXX.XXX.XXX dev ens160 
10.200.1.0/24 via XXX.XXX.XXX.XXX dev ens160 
XXX.XXX.XXX.0/24 dev ens160 proto kernel scope link src XXX.XXX.XXX.XXX 
\end{outlst}

\begin{cmdlst}
# on: jumpbox
ssh root@node-0 ip route
\end{cmdlst}

\begin{outlst}
default via XXX.XXX.XXX.XXX dev ens160 
10.200.1.0/24 via XXX.XXX.XXX.XXX dev ens160 
XXX.XXX.XXX.0/24 dev ens160 proto kernel scope link src XXX.XXX.XXX.XXX 
\end{outlst}

\begin{cmdlst}
# on: jumpbox
ssh root@node-1 ip route
\end{cmdlst}

\begin{outlst}
default via XXX.XXX.XXX.XXX dev ens160 
10.200.0.0/24 via XXX.XXX.XXX.XXX dev ens160 
XXX.XXX.XXX.0/24 dev ens160 proto kernel scope link src XXX.XXX.XXX.XXX 
\end{outlst}

Your interface name
may differ from
\texttt{ens160}.
Check three things:
\texttt{server} has
both
\texttt{10.200.x}
routes, each node
has exactly one
(pointing at the
\emph{other} node),
and the \texttt{via}
addresses are LAN
IPs rather than pod
addresses.

Note that the
\texttt{cni0} routes
do not appear yet on
the workers --- that
entry is created
when the bridge
comes up with the
first pod, in
Chapter 12.

\hypertarget{when-it-goes-wrong-10}{%
\section{When it
goes
wrong}\label{when-it-goes-wrong-10}}

\begin{longtable}[]{@{}
  >{\raggedright\arraybackslash}p{(\columnwidth - 4\tabcolsep) * \real{0.3200}}
  >{\raggedright\arraybackslash}p{(\columnwidth - 4\tabcolsep) * \real{0.2400}}
  >{\raggedright\arraybackslash}p{(\columnwidth - 4\tabcolsep) * \real{0.4400}}@{}}
\toprule\noalign{}
\begin{minipage}[b]{\linewidth}\raggedright
Symptom
\end{minipage} &
\begin{minipage}[b]{\linewidth}\raggedright
Cause
\end{minipage} &
\begin{minipage}[b]{\linewidth}\raggedright
Diagnosis
\end{minipage} \\
\midrule\noalign{}
\endhead
\bottomrule\noalign{}
\endlastfoot
\texttt{Error:}\allowbreak{}\texttt{\ any\ valid\ prefix\ is\ expected}
& A subnet variable
is empty &
\texttt{echo\ \$NODE\_0\_SUBNET}.
Empty means
\texttt{machines.txt}
field separation is
wrong. \\
\texttt{RTNETLINK\ answers:}\allowbreak{}\texttt{\ File\ exists}
& The route is
already present &
Harmless.
\texttt{ip\ route\ del\ \textless{}subnet\textgreater{}}
first if you need to
change it. \\
\texttt{RTNETLINK\ answers:}\allowbreak{}\texttt{\ Network\ is\ unreachable}
& The \texttt{via}
address is not on a
directly connected
subnet & The next
hop must be a LAN IP
reachable at L2 from
that machine. \\
Cross-node pod
traffic still fails
with routes present
& The hypervisor is
dropping packets
with unknown
source/destination
IPs & Disable MAC/IP
anti-spoofing or
enable promiscuous
mode on the virtual
switch. Common on
VMware and some
cloud VPCs. \\
Routes disappear
after reboot &
\texttt{ip\ route\ add}
is not persistent &
Re-run this chapter.
To persist, use
\texttt{/}\allowbreak{}\texttt{etc/}\allowbreak{}\texttt{network/}\allowbreak{}\texttt{interfaces}
post-up hooks or a
systemd unit. \\
Chapter 12: NodePort
works from one node
but not the other &
One route is missing
or reversed &
\texttt{ssh\ root@node-}\allowbreak{}\texttt{0\ ip\ route}
and confirm it
points at
\texttt{node-1}, not
itself. \\
\end{longtable}

\hypertarget{check-your-understanding-11}{%
\section{Check your
understanding}\label{check-your-understanding-11}}

\textbf{Q1. Why does
\texttt{server} need
routes to both pod
subnets when it runs
no pods?}

Because the control
plane initiates
connections into the
cluster. The API
server proxies
\texttt{kubectl\ logs},
\texttt{exec}, and
\texttt{port-forward}
to kubelets, and
those paths can
involve pod
addresses. Without
routes, those
operations fail with
timeouts while the
cluster otherwise
looks healthy --- a
confusing failure
that presents as
``logs are broken''
rather than
``routing is
broken.''

\textbf{Q2. You add
a third worker,
\texttt{node-2},
with
\texttt{10.200.2.0/24}.
How many routes must
you add, and where?}

Five. Two on
\texttt{node-2} (to
\texttt{10.200.0.0/24}
and
\texttt{10.200.1.0/24}),
one on
\texttt{node-0} (to
\texttt{10.200.2.0/24}),
one on
\texttt{node-1}
(same), and one on
\texttt{server}
(same). Note that
adding one node
means touching every
existing machine ---
the count grows as
\emph{n}(\emph{n}\textminus{}1)
plus one per node
for the control
plane, which is why
nobody does this by
hand past about
three nodes.

\textbf{Q3. What
would change if
\texttt{node-0} and
\texttt{node-1} were
on different subnets
behind a router?}

The \texttt{via}
next hop would have
to be the router
rather than the
node, and the router
itself would need
routes for both pod
CIDRs pointing at
the correct nodes.
At that point you
are configuring pod
routing in network
infrastructure you
may not control,
which is exactly the
situation overlay
networks (VXLAN,
Geneve) were
designed for: they
tunnel pod traffic
inside ordinary
node-to-node
packets, so the
intervening network
never needs to know
pod addresses exist.

\hypertarget{chap-12}{%
\chapter{Smoke
Test}\label{chap-12}}

\upstreamlab{docs/12-smoke-test.md}

\hypertarget{what-youre-building-12}{%
\section{What you're
building}\label{what-youre-building-12}}

Nothing new. You are
exercising five
capabilities end to
end: encryption at
rest, Deployments,
port forwarding,
logs, exec, and
Services. Each one
validates a
different chapter's
work.

\hypertarget{why-its-necessary-12}{%
\section{Why it's
necessary}\label{why-its-necessary-12}}

\begin{whybox}

This is the
acceptance test, and
its real value is
diagnostic mapping
--- each test fails
in a way that points
at a specific
earlier chapter:

\end{whybox}

\begin{longtable}[]{@{}
  >{\raggedright\arraybackslash}p{(\columnwidth - 4\tabcolsep) * \real{0.3333}}
  >{\raggedright\arraybackslash}p{(\columnwidth - 4\tabcolsep) * \real{0.3333}}
  >{\raggedright\arraybackslash}p{(\columnwidth - 4\tabcolsep) * \real{0.3333}}@{}}
\toprule\noalign{}
\begin{minipage}[b]{\linewidth}\raggedright
Test
\end{minipage} &
\begin{minipage}[b]{\linewidth}\raggedright
Validates
\end{minipage} &
\begin{minipage}[b]{\linewidth}\raggedright
Fails if you got
wrong
\end{minipage} \\
\midrule\noalign{}
\endhead
\bottomrule\noalign{}
\endlastfoot
Data encryption &
API server
encryption provider
& Chapter 6 \\
Deployment &
Scheduler,
controller manager,
kubelet, containerd,
CNI & Chapters 8,
9 \\
Port forwarding &
API server
\textrightarrow{}
kubelet path,
\texttt{socat} on
the node & Chapters
8 (RBAC), 9 \\
Logs & API server
\textrightarrow{}
kubelet
\texttt{nodes/log}
subresource &
Chapter 8 (RBAC) \\
Exec & API server
\textrightarrow{}
kubelet
\texttt{nodes/proxy}
subresource &
Chapter 8 (RBAC) \\
NodePort Service &
kube-proxy iptables,
routes & Chapters 9,
11 \\
\end{longtable}

Run them in order. A
failure tells you
where to look.

\hypertarget{key-concepts-12}{%
\section{Key
concepts}\label{key-concepts-12}}

\textbf{1. The
hexdump test is the
only direct proof of
encryption.} Reading
a Secret back
through
\texttt{kubectl}
tells you nothing
--- the API server
decrypts
transparently.
Reading the raw etcd
value bypasses the
API server entirely
and shows you what
is actually on disk.

\textbf{2.
\texttt{port-forward},
\texttt{logs}, and
\texttt{exec} all
traverse the same
path.} Client
\textrightarrow{}
API server
\textrightarrow{}
kubelet on port
10250
\textrightarrow{}
container. All three
depend on Chapter
8's RBAC ClusterRole
and on the kubelet's
server certificate.
If all three fail
identically with
\texttt{Forbidden},
it is one cause, not
three.

\textbf{3. There is
no cluster DNS in
this build.} No
CoreDNS is
installed, so
Services are
reachable only by
ClusterIP and
NodePort, never by
name.
\texttt{nginx.}\allowbreak{}\texttt{default.}\allowbreak{}\texttt{svc.}\allowbreak{}\texttt{cluster.}\allowbreak{}\texttt{local}
will not resolve
from inside a pod.
This is why the
smoke test reaches
the Service by node
hostname and port
rather than by name.

\textbf{4.
\texttt{LoadBalancer}
is unavailable, by
design.} No cloud
provider integration
means no external IP
allocation.
\texttt{NodePort} is
the ceiling here.

\hypertarget{the-work-11}{%
\section{The
work}\label{the-work-11}}

All commands run
from the
\textbf{jumpbox}
unless noted.

\textbf{Data
encryption.} Create
a secret:

\begin{cmdlst}
# on: jumpbox
kubectl create secret generic kubernetes-the-hard-way \
  --from-literal="mykey=mydata"
\end{cmdlst}

Read the raw bytes
out of etcd:

\begin{cmdlst}
# on: jumpbox
ssh root@server \
    'etcdctl get /registry/secrets/default/kubernetes-the-hard-way | hexdump -C'
\end{cmdlst}

\begin{outlst}
00000000  2f 72 65 67 69 73 74 72  79 2f 73 65 63 72 65 74  |/registry/secret|
00000010  73 2f 64 65 66 61 75 6c  74 2f 6b 75 62 65 72 6e  |s/default/kubern|
00000020  65 74 65 73 2d 74 68 65  2d 68 61 72 64 2d 77 61  |etes-the-hard-wa|
00000030  79 0a 6b 38 73 3a 65 6e  63 3a 61 65 73 63 62 63  |y.k8s:enc:aescbc|
00000040  3a 76 31 3a 6b 65 79 31  3a 4f 1b 80 d8 89 72 f4  |:v1:key1:O....r.|
00000050  60 8a 2c a0 76 1a e1 dc  98 d6 00 7a a4 2f f3 92  |`.,.v......z./..|
\end{outlst}

\emph{(Output
truncated; yours
continues for
roughly 350 bytes of
ciphertext.)}

The etcd key should
be prefixed with
\texttt{k8s:}\allowbreak{}\texttt{enc:}\allowbreak{}\texttt{aescbc:}\allowbreak{}\texttt{v1:}\allowbreak{}\texttt{key1},
which indicates the
\texttt{aescbc}
provider was used to
encrypt the data
with the
\texttt{key1}
encryption key.

\begin{riskbox}\calloutlabel{brick}{RISK}

Read the hexdump
carefully. The
\textbf{key path} is
plaintext --- anyone
with etcd access can
enumerate which
Secrets exist and in
which namespaces.
Only the
\textbf{value} is
encrypted. If you
can read
\texttt{mydata}
anywhere in the
ASCII column,
encryption is not
working and Chapter
6 is the place to
look.

\end{riskbox}

\textbf{Deployments.}
Create one:

\begin{cmdlst}
# on: jumpbox
kubectl create deployment nginx \
  --image=nginx:latest
\end{cmdlst}

\begin{cmdlst}
# on: jumpbox
kubectl get pods -l app=nginx
\end{cmdlst}

\begin{cmdlst}
NAME                     READY   STATUS    RESTARTS   AGE
nginx-56fcf95486-c8dnx   1/1     Running   0          8s
\end{cmdlst}

Reaching
\texttt{Running}
proves a long chain
worked: the
controller manager
created a ReplicaSet
and Pod, the
scheduler assigned a
node, that node's
kubelet saw the
assignment,
containerd pulled
\texttt{nginx:latest}
from Docker Hub, CNI
allocated an IP from
the node's subnet
and attached the pod
to \texttt{cni0},
and runc started the
container.

A pod stuck in
\texttt{Pending}
means the scheduler
is not running or no
node is
\texttt{Ready}.
Stuck in
\texttt{ContainerCreating}
means CNI or image
pull ---
\texttt{kubectl\ describe\ pod\ \textless{}name\textgreater{}}
gives the specific
reason.

\textbf{Port
forwarding.} Capture
the pod name:

\begin{cmdlst}
# on: jumpbox
POD_NAME=$(kubectl get pods -l app=nginx \
  -o jsonpath="{.items[0].metadata.name}")
\end{cmdlst}

\begin{cmdlst}
# on: jumpbox
kubectl port-forward $POD_NAME 8080:80
\end{cmdlst}

\begin{outlst}
Forwarding from 127.0.0.1:8080 -> 80
Forwarding from [::1]:8080 -> 80
\end{outlst}

This blocks.
\textbf{In a second
terminal on the
jumpbox:}

\begin{cmdlst}
# on: jumpbox (second terminal)
curl --head http://127.0.0.1:8080
\end{cmdlst}

\begin{outlst}
HTTP/1.1 200 OK
Server: nginx/1.27.4
Date: Sun, 06 Apr 2025 17:17:12 GMT
Content-Type: text/html
Content-Length: 615
Last-Modified: Wed, 05 Feb 2025 11:06:32 GMT
Connection: keep-alive
ETag: "67a34638-267"
Accept-Ranges: bytes
\end{outlst}

Return to the first
terminal and stop it
with
\texttt{Ctrl-C}:

\begin{outlst}
Forwarding from 127.0.0.1:8080 -> 80
Forwarding from [::1]:8080 -> 80
Handling connection for 8080
^C
\end{outlst}

The tunnel runs
jumpbox
\textrightarrow{}
API server
\textrightarrow{}
kubelet
\textrightarrow{}
\texttt{socat} on
the node
\textrightarrow{}
pod. It is the one
test that proves
\texttt{socat} was
installed in Chapter
9, and it uses no
pod routing at all
--- which makes it a
useful isolation
test when NodePort
fails but
port-forward works.

\textbf{Logs:}

\begin{cmdlst}
# on: jumpbox
kubectl logs $POD_NAME
\end{cmdlst}

\begin{outlst}
...
127.0.0.1 - - [06/Apr/2025:17:17:12 +0000] "HEAD / HTTP/1.1" 200 0 "-" "curl/7.88.1" "-"
\end{outlst}

The entry you see is
from the
port-forward request
you just made ---
the source shows as
\texttt{127.0.0.1}
because the traffic
arrived through the
forwarding tunnel.

\textbf{Exec:}

\begin{cmdlst}
# on: jumpbox
kubectl exec -ti $POD_NAME -- nginx -v
\end{cmdlst}

\begin{outlst}
nginx version: nginx/1.27.4
\end{outlst}

\textbf{Services.}
Expose the
deployment:

\begin{cmdlst}
# on: jumpbox
kubectl expose deployment nginx \
  --port 80 --type NodePort
\end{cmdlst}

\begin{notebox}\calloutlabel{steel}{FROM THE SOURCE}

The LoadBalancer
service type can not
be used because your
cluster is not
configured with
cloud provider
integration.

\end{notebox}

Get the allocated
node port:

\begin{cmdlst}
# on: jumpbox
NODE_PORT=$(kubectl get svc nginx \
  --output=jsonpath='{range .spec.ports[0]}{.nodePort}')
\end{cmdlst}

It will fall in
\texttt{30000-32767},
the range set by
\texttt{-\/-service-}\allowbreak{}\texttt{node-}\allowbreak{}\texttt{port-}\allowbreak{}\texttt{range}
on the API server.

Find which node the
pod landed on:

\begin{cmdlst}
# on: jumpbox
NODE_NAME=$(kubectl get pods \
  -l app=nginx \
  -o jsonpath="{.items[0].spec.nodeName}")
\end{cmdlst}

\begin{cmdlst}
# on: jumpbox
curl -I http://${NODE_NAME}:${NODE_PORT}
\end{cmdlst}

\begin{outlst}
Server: nginx/1.27.4
Date: Sun, 06 Apr 2025 17:18:36 GMT
Content-Type: text/html
Content-Length: 615
Last-Modified: Wed, 05 Feb 2025 11:06:32 GMT
Connection: keep-alive
ETag: "67a34638-267"
Accept-Ranges: bytes
\end{outlst}

This works because
the jumpbox resolves
\texttt{node-0} /
\texttt{node-1} from
Chapter 3's
\texttt{/etc/hosts}
entries.

\begin{beyondbox}\calloutlabel{teal}{BEYOND THE SOURCE}

a NodePort is opened
on \textbf{every}
node, not just the
one running the pod.
Try
\texttt{curl\ -}\allowbreak{}\texttt{I\ http:}\allowbreak{}\texttt{/}\allowbreak{}\texttt{/}\allowbreak{}\texttt{\textless{}the\ other\ node\textgreater{}:}\allowbreak{}\texttt{\$\{NODE\_PORT\}}
--- it should also
return 200, with
kube-proxy
forwarding across
nodes. \textbf{This
is the real test of
Chapter 11's
routes.} If the
pod's own node works
and the other does
not, a route is
missing or reversed.

\end{beyondbox}

\hypertarget{when-it-goes-wrong-11}{%
\section{When it
goes
wrong}\label{when-it-goes-wrong-11}}

\begin{longtable}[]{@{}
  >{\raggedright\arraybackslash}p{(\columnwidth - 4\tabcolsep) * \real{0.3200}}
  >{\raggedright\arraybackslash}p{(\columnwidth - 4\tabcolsep) * \real{0.2400}}
  >{\raggedright\arraybackslash}p{(\columnwidth - 4\tabcolsep) * \real{0.4400}}@{}}
\toprule\noalign{}
\begin{minipage}[b]{\linewidth}\raggedright
Symptom
\end{minipage} &
\begin{minipage}[b]{\linewidth}\raggedright
Cause
\end{minipage} &
\begin{minipage}[b]{\linewidth}\raggedright
Diagnosis
\end{minipage} \\
\midrule\noalign{}
\endhead
\bottomrule\noalign{}
\endlastfoot
Hexdump shows
readable
\texttt{mydata} &
Encryption config
not loaded, or
\texttt{identity}
provider listed
first & On
\texttt{server}:
\texttt{cat\ /}\allowbreak{}\texttt{var/}\allowbreak{}\texttt{lib/}\allowbreak{}\texttt{kubernetes/}\allowbreak{}\texttt{encryption-}\allowbreak{}\texttt{config.}\allowbreak{}\texttt{yaml};
check
\texttt{-\/-encryption-}\allowbreak{}\texttt{provider-}\allowbreak{}\texttt{config}
in the unit. \\
\texttt{etcdctl\ get}
returns nothing &
Wrong key path, or
the Secret was
created in another
namespace &
\texttt{ssh\ root@server\ "etcdctl\ get\ /}\allowbreak{}\texttt{registry/}\allowbreak{}\texttt{secrets/}\allowbreak{}\texttt{\ -\/-prefix\ -\/-keys-}\allowbreak{}\texttt{only"}. \\
Pod \texttt{Pending}
& Scheduler down, or
no \texttt{Ready}
node &
\texttt{kubectl\ get\ nodes};
\texttt{ssh\ root@server\ systemctl\ is-}\allowbreak{}\texttt{active\ kube-}\allowbreak{}\texttt{scheduler}. \\
Pod
\texttt{ContainerCreating}
& CNI failure or
image pull failure &
\texttt{kubectl\ describe\ pod\ \$POD\_NAME}
--- the Events
section names the
cause. \\
Pod
\texttt{ImagePullBackOff}
& Nodes lack
outbound internet or
DNS & On the node:
\texttt{curl\ -}\allowbreak{}\texttt{I\ https:}\allowbreak{}\texttt{/}\allowbreak{}\texttt{/}\allowbreak{}\texttt{registry-}\allowbreak{}\texttt{1.}\allowbreak{}\texttt{docker.}\allowbreak{}\texttt{io/}\allowbreak{}\texttt{v2/}\allowbreak{}. \\
\texttt{logs}/\texttt{exec}/\texttt{port-forward}
all give
\texttt{Forbidden\ (user=kubernetes.}\allowbreak{}\texttt{.}\allowbreak{}\texttt{.}\allowbreak{}\texttt{)}
& Chapter 8's RBAC
ClusterRole never
applied & Re-run
\texttt{kubectl\ apply\ -}\allowbreak{}\texttt{f\ kube-}\allowbreak{}\texttt{apiserver-}\allowbreak{}\texttt{to-}\allowbreak{}\texttt{kubelet.}\allowbreak{}\texttt{yaml}
on
\texttt{server}. \\
\texttt{port-forward}:
\texttt{unable\ to\ do\ port\ forwarding:}\allowbreak{}\texttt{\ socat\ not\ found}
& \texttt{socat}
missing on that node
&
\texttt{apt-}\allowbreak{}\texttt{get\ -}\allowbreak{}\texttt{y\ install\ socat}
on the node. \\
\texttt{logs}/\texttt{exec}
fail with a TLS
error to :10250 &
kubelet serving
certificate wrong,
or
\texttt{-\/-kubelet-}\allowbreak{}\texttt{certificate-}\allowbreak{}\texttt{authority}
mismatched &
\texttt{journalctl\ -}\allowbreak{}\texttt{u\ kubelet}
on the node; check
\texttt{/}\allowbreak{}\texttt{var/}\allowbreak{}\texttt{lib/}\allowbreak{}\texttt{kubelet/}\allowbreak{}\texttt{kubelet.}\allowbreak{}\texttt{crt}. \\
NodePort works on
the pod's node,
fails on the other &
Missing or reversed
route &
\texttt{ssh\ root@node-}\allowbreak{}\texttt{X\ ip\ route}
--- Chapter 11. \\
NodePort fails on
both nodes &
kube-proxy not
running, or
\texttt{br-netfilter}
sysctls unset &
\texttt{systemctl\ is-}\allowbreak{}\texttt{active\ kube-}\allowbreak{}\texttt{proxy};
\texttt{iptables\ -}\allowbreak{}\texttt{t\ nat\ -}\allowbreak{}\texttt{L\ KUBE-}\allowbreak{}\texttt{SERVICES\ -}\allowbreak{}\texttt{n\ \textbackslash{}\textbar{}\ head}. \\
\texttt{NODE\_PORT}
is empty & The
Service was not
created &
\texttt{kubectl\ get\ svc\ nginx}. \\
\end{longtable}

\hypertarget{check-your-understanding-12}{%
\section{Check your
understanding}\label{check-your-understanding-12}}

\textbf{Q1. The etcd
hexdump shows the
key path in
cleartext but the
value encrypted.
What does that tell
an attacker with
read access to
etcd?}

Complete metadata:
every Secret's name
and namespace, and
by extension the
shape of the whole
cluster --- which
applications exist,
which namespaces are
in use, how things
are organized. They
cannot read Secret
\emph{contents}
without the
encryption key, but
``there is a Secret
called
\texttt{prod-}\allowbreak{}\texttt{db-}\allowbreak{}\texttt{credentials}
in namespace
\texttt{payments}''
is itself useful
intelligence.
Encryption at rest
protects values, not
the existence of
objects.

\textbf{Q2.
\texttt{kubectl\ port-}\allowbreak{}\texttt{forward}
succeeds but
\texttt{curl\ http:}\allowbreak{}\texttt{/}\allowbreak{}\texttt{/}\allowbreak{}\texttt{\$\{NODE\_NAME\}:}\allowbreak{}\texttt{\$\{NODE\_PORT\}}
times out. What is
the most likely
cause, and what does
the combination rule
out?}

Most likely a
missing pod network
route or a
kube-proxy problem.
The combination
rules out a great
deal: port-forward
proves the API
server can reach the
kubelet, that the
RBAC grant is in
place, that the pod
is running, and that
the container is
serving on port 80.
The only things left
in the NodePort path
that port-forward
does not exercise
are kube-proxy's
iptables rules and
the routing from
Chapter 11.

\textbf{Q3. You
delete the nginx
pod. It comes back
with a new name and
a new IP. Which
components did
what?}

The kubelet notices
the container is
gone and reports it.
The ReplicaSet
controller (inside
\texttt{kube-}\allowbreak{}\texttt{controller-}\allowbreak{}\texttt{manager})
sees one replica
instead of the
desired one and
creates a new Pod
object. The
scheduler assigns it
a node. That node's
kubelet pulls the
image (cached) and
starts it, and CNI
allocates a fresh IP
from the node's
\texttt{host-local}
range --- the old
one is released back
to the pool. The
Service's endpoint
list is updated by
the endpoints
controller, and
kube-proxy rewrites
its iptables rules.
Five independent
controllers, no
central coordinator,
all reconciling
against etcd.

\hypertarget{chap-13}{%
\chapter{Cleaning
Up}\label{chap-13}}

\upstreamlab{docs/13-cleanup.md}

\hypertarget{what-youre-building-13}{%
\section{What you're
building}\label{what-youre-building-13}}

Nothing. You are
removing the four
machines.

\hypertarget{why-its-necessary-13}{%
\section{Why it's
necessary}\label{why-its-necessary-13}}

\begin{whybox}

Two reasons beyond
reclaiming
resources. This
cluster has
\texttt{PermitRootLogin\ yes}
on three hosts, an
unauthenticated
etcd, and an
unencrypted CA
private key with a
ten-year validity
sitting on both the
jumpbox and
\texttt{server}.
Leaving it running
on any network you
share with anything
else is a genuine
risk.

\end{whybox}

The second reason is
pedagogical. The
build is fully
reproducible from
\texttt{machines.txt}
plus these thirteen
chapters. Tearing it
down and rebuilding
it from scratch ---
ideally without
looking at the
commands --- is the
single most
effective way to
convert this from
something you
followed into
something you know.

\hypertarget{key-concepts-13}{%
\section{Key
concepts}\label{key-concepts-13}}

\textbf{1. There is
nothing to clean up
in a cloud
provider.}
\emph{Previous
versions of this
guide made use of
GCP resources for
various aspects of
compute and
networking. The
current version is
agnostic, and all
configuration is
performed on the
\texttt{jumpbox},
\texttt{server}, or
nodes.} Every
artifact lives on
one of the four
machines. No load
balancers, firewall
rules, or managed
disks are left
behind.

\textbf{2. Destroy
the jumpbox too.} It
holds
\texttt{ca.key},
every component
private key,
\texttt{admin.}\allowbreak{}\texttt{kubeconfig},
and
\texttt{encryption-}\allowbreak{}\texttt{config.}\allowbreak{}\texttt{yaml}.
It is the
highest-value target
in the build --- a
working
cluster-admin
credential and the
ability to mint
more.

\textbf{3. Nothing
is persistent by
design.} No
PersistentVolumes
were created, no
external storage
attached. Deleting
the VMs deletes
everything.

\hypertarget{the-work-12}{%
\section{The
work}\label{the-work-12}}

Clean up is as
simple as deleting
all virtual machines
you created for this
exercise.

Delete
\texttt{jumpbox},
\texttt{server},
\texttt{node-0}, and
\texttt{node-1}
through whatever
provisioned them ---
your hypervisor UI,
\texttt{virsh\ destroy\ \&\&\ virsh\ undefine},
\texttt{multipass\ delete\ -\/-purge},
or your cloud
console.

\begin{beyondbox}\calloutlabel{teal}{BEYOND THE SOURCE}

if you are keeping
the machines to
rebuild on, remove
the
\texttt{/etc/hosts}
entries you appended
in Chapter 3 first.
A stale
\texttt{\#\ Kubernetes\ The\ Hard\ Way}
block pointing at
reassigned IPs
produces name
resolution failures
on the next build
that look exactly
like network faults.
Also delete
\texttt{\textasciitilde{}/.kube/config}
on the jumpbox and
remove
\texttt{/}\allowbreak{}\texttt{root/}\allowbreak{}\texttt{kubernetes-}\allowbreak{}\texttt{the-}\allowbreak{}\texttt{hard-}\allowbreak{}\texttt{way/}\allowbreak{}\texttt{*.}\allowbreak{}\texttt{key}
before reusing that
host for anything
else.

\end{beyondbox}

See
\protect\hyperlink{app-teardown}{Appendix
D} for a per-machine
in-place teardown if
you want to reset
without destroying
the VMs.

\hypertarget{verification-11}{%
\section{Verification}\label{verification-11}}

No machine responds
to \texttt{ping} or
\texttt{ssh}. No
stale entries remain
in your DHCP
reservations or DNS.

\hypertarget{when-it-goes-wrong-12}{%
\section{When it
goes
wrong}\label{when-it-goes-wrong-12}}

\begin{longtable}[]{@{}
  >{\raggedright\arraybackslash}p{(\columnwidth - 4\tabcolsep) * \real{0.3200}}
  >{\raggedright\arraybackslash}p{(\columnwidth - 4\tabcolsep) * \real{0.2400}}
  >{\raggedright\arraybackslash}p{(\columnwidth - 4\tabcolsep) * \real{0.4400}}@{}}
\toprule\noalign{}
\begin{minipage}[b]{\linewidth}\raggedright
Symptom
\end{minipage} &
\begin{minipage}[b]{\linewidth}\raggedright
Cause
\end{minipage} &
\begin{minipage}[b]{\linewidth}\raggedright
Fix
\end{minipage} \\
\midrule\noalign{}
\endhead
\bottomrule\noalign{}
\endlastfoot
Next build:
\texttt{REMOTE\ HOST\ IDENTIFICATION\ HAS\ CHANGED}
& New VMs reusing
old IPs; stale host
keys &
\texttt{ssh-}\allowbreak{}\texttt{keygen\ -}\allowbreak{}\texttt{R\ \textless{}hostname\textgreater{}}
and
\texttt{ssh-}\allowbreak{}\texttt{keygen\ -}\allowbreak{}\texttt{R\ \textless{}ip\textgreater{}},
or delete
\texttt{\textasciitilde{}/}\allowbreak{}\texttt{.}\allowbreak{}\texttt{ssh/}\allowbreak{}\texttt{known\_hosts}. \\
Next build: names
resolve to the wrong
machines & Old
\texttt{/etc/hosts}
block still present
& Remove the
\texttt{\#\ Kubernetes\ The\ Hard\ Way}
block before
re-running Chapter
3. \\
Next build: nodes
will not register &
Old
\texttt{admin.}\allowbreak{}\texttt{kubeconfig}
or certificates from
the previous build &
Start from a clean
clone; regenerate
all certificates
from Chapter 4. \\
\end{longtable}

\hypertarget{check-your-understanding-13}{%
\section{Check your
understanding}\label{check-your-understanding-13}}

\textbf{Q1. Which
single file, if it
leaked, would most
compromise this
cluster?}

\texttt{ca.key}.
With it, anyone can
issue a certificate
with
\texttt{CN=admin,\ O=system:}\allowbreak{}\texttt{masters}
and hold permanent
cluster-admin --- or
impersonate any
node, the scheduler,
or the API server
itself. It cannot be
revoked (there is no
CRL or OCSP in this
build) and it is
valid for ten years.
\texttt{encryption-}\allowbreak{}\texttt{config.}\allowbreak{}\texttt{yaml}
is a close second,
but it only decrypts
Secrets;
\texttt{ca.key}
grants full control.

\textbf{Q2. You want
to rebuild but keep
the machines. What
is the minimum
reset?}

On \texttt{server}
and both workers:
stop and disable
every unit, remove
\texttt{/}\allowbreak{}\texttt{etc/}\allowbreak{}\texttt{systemd/}\allowbreak{}\texttt{system/}\allowbreak{}\texttt{\{etcd,kube-}\allowbreak{}\texttt{*,containerd\}.}\allowbreak{}\texttt{service},
and delete
\texttt{/}\allowbreak{}\texttt{var/}\allowbreak{}\texttt{lib/}\allowbreak{}\texttt{\{etcd,kubernetes,kubelet,kube-}\allowbreak{}\texttt{proxy\}},
\texttt{/etc/kubernetes},
\texttt{/etc/etcd},
\texttt{/etc/cni/net.d},
and
\texttt{/opt/cni/bin}.
On the jumpbox:
delete the generated
\texttt{*.crt},
\texttt{*.key},
\texttt{*.csr},
\texttt{*.kubeconfig},
\texttt{encryption-}\allowbreak{}\texttt{config.}\allowbreak{}\texttt{yaml},
and
\texttt{\textasciitilde{}/.kube/config}.
Then remove the
\texttt{/etc/hosts}
block on all four.
Appendix D has this
as a runnable
sequence.

\textbf{Q3. What
have you built that
would be worth
keeping?}

\texttt{machines.txt}
and your notes on
what broke. The
certificates and
kubeconfigs are ten
seconds of
regeneration; the
record of which
failure mode you hit
and what the log
line actually said
is the part that
does not reproduce
itself.

\appendix

\hypertarget{app-glossary}{%
\chapter{Glossary}\label{app-glossary}}

\textbf{Admission
plugin} --- A
component that
inspects an
authorized API
request before it is
persisted, and may
mutate or reject it.
This build enables
\texttt{NamespaceLifecycle},
\texttt{NodeRestriction},
\texttt{LimitRanger},
\texttt{ServiceAccount},
\texttt{DefaultStorageClass},
and
\texttt{ResourceQuota}.

\textbf{aescbc} ---
The AES-256-CBC
encryption provider
used for Secrets at
rest, requiring a
32-byte key. Named
as \texttt{key1} in
\texttt{encryption-}\allowbreak{}\texttt{config.}\allowbreak{}\texttt{yaml}.

\textbf{Authentication}
--- Establishing
\emph{who} a caller
is. Here,
exclusively via
X.509 client
certificates: the
certificate's CN
becomes the username
and its O becomes
the group.

\textbf{Authorization}
--- Establishing
whether an
authenticated caller
\emph{may} perform
an action. Here,
\texttt{Node,RBAC}
--- two modes
evaluated in order,
either of which can
allow.

\textbf{cni0} ---
The Linux bridge
created on each
worker by the CNI
\texttt{bridge}
plugin. Pods attach
to it; it holds the
first address of the
node's pod subnet
and acts as their
default gateway.
Created on demand
when the first pod
starts.

\textbf{CNI
(Container Network
Interface)} --- The
plugin specification
the kubelet uses to
attach pods to a
network.
Configuration lives
in
\texttt{/etc/cni/net.d},
plugin binaries in
\texttt{/opt/cni/bin}.

\textbf{containerd}
--- The container
runtime. Pulls
images, manages
snapshots, and
invokes
\texttt{runc} to
create containers.
Reached by the
kubelet over the CRI
socket at
\texttt{unix:}\allowbreak{}\texttt{/}\allowbreak{}\texttt{/}\allowbreak{}\texttt{/}\allowbreak{}\texttt{var/}\allowbreak{}\texttt{run/}\allowbreak{}\texttt{containerd/}\allowbreak{}\texttt{containerd.}\allowbreak{}\texttt{sock}.

\textbf{Cluster
CIDR} ---
\texttt{10.200.0.0/16}.
The address space
from which all pod
IPs are allocated,
subdivided per node
(\texttt{10.200.0.0/24}
for \texttt{node-0},
\texttt{10.200.1.0/24}
for
\texttt{node-1}).

\textbf{ClusterIP}
--- A virtual IP
from the service
range, assigned to a
Service. Not bound
to any interface;
realized entirely by
kube-proxy's
iptables rules.

\textbf{crictl} ---
A CRI-level
debugging client
that talks directly
to containerd,
bypassing
Kubernetes.
Installed on workers
for diagnosis.

\textbf{CSR
(Certificate Signing
Request)} --- An
unsigned certificate
request produced by
\texttt{openssl\ req\ -}\allowbreak{}\texttt{new},
carrying a subject
and requested
extensions,
submitted to a CA
for signing.

\textbf{etcd} ---
The distributed
key-value store
holding all cluster
state. Single member
here, on loopback,
over plain HTTP.

\textbf{\texttt{-\/-embed-}\allowbreak{}\texttt{certs=true}}
--- A
\texttt{kubectl\ config}
flag that copies
certificate and key
\emph{contents} into
a kubeconfig rather
than referencing
file paths. Makes
the file portable
and as sensitive as
the key inside it.

\textbf{Jumpbox} ---
The administrative
host. Not a cluster
member; holds the
CA, all generated
credentials, and the
binary cache.

\textbf{kube-apiserver}
--- The only
component that reads
and writes etcd.
Authenticates and
authorizes every
request. The hub of
the architecture.

\textbf{kube-controller-manager}
--- Runs the
reconciliation
loops: deployments,
replicasets,
endpoints, node
lifecycle, service
accounts, CSR
signing.

\textbf{kube-proxy}
--- Programs
iptables rules on
each node so
ClusterIPs and
NodePorts route to
pod IPs. Does not
run containers or
carry traffic
itself.

\textbf{kube-scheduler}
--- Assigns pending
pods to nodes by
writing
\texttt{spec.nodeName}.
Does not contact the
kubelet.

\textbf{kubeconfig}
--- A YAML
credential bundle: a
cluster (server URL
+ CA), a user
(client
certificate), and a
context joining
them.

\textbf{kubelet} ---
The node agent.
Watches the API
server for pods
assigned to its
node, instructs
containerd, invokes
CNI, and reports
status. Serves its
own API on port
10250.

\textbf{Node
authorizer} --- A
dedicated
authorization mode
that scopes each
kubelet to objects
related to its own
node. Requires
username
\texttt{system:}\allowbreak{}\texttt{node:}\allowbreak{}\texttt{\textless{}nodeName\textgreater{}}
and group
\texttt{system:nodes}.

\textbf{NodePort}
--- A Service type
that opens the same
high port
(30000--32767) on
\emph{every} node,
forwarding to the
Service's endpoints
wherever they run.

\textbf{\texttt{nodes/proxy},
\texttt{nodes/log},
\texttt{nodes/stats},
\texttt{nodes/spec},
\texttt{nodes/metrics}}
--- Kubelet API
subresources.
Granted to user
\texttt{kubernetes}
by
\texttt{kube-}\allowbreak{}\texttt{apiserver-}\allowbreak{}\texttt{to-}\allowbreak{}\texttt{kubelet.}\allowbreak{}\texttt{yaml};
required for
\texttt{kubectl\ logs},
\texttt{exec}, and
\texttt{port-forward}.

\textbf{runc} ---
The low-level OCI
runtime that
actually creates
namespaces and
cgroups. Invoked by
containerd, never
directly by
Kubernetes.

\textbf{SAN (Subject
Alternative Name)}
--- The certificate
extension listing
valid hostnames and
IPs for a server
certificate.
Preserved through
signing only because
of
\texttt{-}\allowbreak{}\texttt{copy\_extensions\ copyall}.

\textbf{ServiceAccount
token} --- A JWT
mounted into pods,
signed by the
\texttt{service-}\allowbreak{}\texttt{accounts}
key pair, with
\texttt{iss} set to
\texttt{https:}\allowbreak{}\texttt{/}\allowbreak{}\texttt{/}\allowbreak{}\texttt{server.}\allowbreak{}\texttt{kubernetes.}\allowbreak{}\texttt{local:}\allowbreak{}\texttt{6443}.

\textbf{Service
CIDR} ---
\texttt{10.32.0.0/24}.
The ClusterIP
allocation range;
\texttt{10.32.0.1}
is reserved for the
\texttt{kubernetes}
Service fronting the
API server.

\textbf{\texttt{system:masters}}
--- A group with a
built-in
ClusterRoleBinding
to
\texttt{cluster-admin}.
The \texttt{admin}
certificate's O
field.

\textbf{\texttt{system:nodes}
/
\texttt{system:}\allowbreak{}\texttt{node:}\allowbreak{}\texttt{\textless{}name\textgreater{}}}
--- The group and
username format the
Node authorizer
requires from
kubelets.

\textbf{\texttt{system:}\allowbreak{}\texttt{node-}\allowbreak{}\texttt{proxier}}
--- The group
granting kube-proxy
its permissions, via
a built-in
ClusterRoleBinding.

\textbf{Webhook
authorization
(kubelet)} ---
\texttt{authorization.}\allowbreak{}\texttt{mode:}\allowbreak{}\texttt{\ Webhook}
in
\texttt{kubelet-}\allowbreak{}\texttt{config.}\allowbreak{}\texttt{yaml},
delegating every
inbound
authorization
decision to the API
server via
SubjectAccessReview.

\hypertarget{app-commands}{%
\chapter{Command
Reference}\label{app-commands}}

\hypertarget{host-cheat-sheet}{%
\section{Host cheat
sheet}\label{host-cheat-sheet}}

\begin{longtable}[]{@{}
  >{\raggedright\arraybackslash}p{(\columnwidth - 2\tabcolsep) * \real{0.5000}}
  >{\raggedright\arraybackslash}p{(\columnwidth - 2\tabcolsep) * \real{0.5000}}@{}}
\toprule\noalign{}
\begin{minipage}[b]{\linewidth}\raggedright
Chapter
\end{minipage} &
\begin{minipage}[b]{\linewidth}\raggedright
Runs on
\end{minipage} \\
\midrule\noalign{}
\endhead
\bottomrule\noalign{}
\endlastfoot
1 & all four \\
2 & jumpbox \\
3 & jumpbox (plus
one \texttt{sshd}
edit per cluster
machine) \\
4, 5, 6 & jumpbox \\
7, 8 & jumpbox to
copy, then
\texttt{server} \\
9 & jumpbox to copy,
then \texttt{node-0}
\textbf{and}
\texttt{node-1} \\
10, 11, 12 &
jumpbox \\
13 & hypervisor \\
\end{longtable}

\hypertarget{diagnosis-first-commands}{%
\section{Diagnosis
--- first
commands}\label{diagnosis-first-commands}}

\begin{cmdlst}
# Is the service alive?
systemctl is-active <unit>
systemctl status <unit>

# Why did it fail? (always the second command)
journalctl -u <unit> -n 50 --no-pager

# Is it listening?
ss -tlnp | grep -E '6443|2379|10250|10256'

# Control plane reachable, unauthenticated?
curl --cacert ca.crt https://server.kubernetes.local:6443/version

# etcd healthy?                             (on: server)
etcdctl endpoint health
etcdctl member list

# Node and pod state                        (on: jumpbox)
kubectl get nodes -o wide
kubectl get pods -A -o wide
kubectl describe pod <name>        # Events section names most failures

# Container runtime, below Kubernetes       (on: a worker)
crictl ps -a
crictl logs <container-id>

# Pod networking                            (on: any cluster machine)
ip route
ip addr show cni0
lsmod | grep br_netfilter
sysctl net.bridge.bridge-nf-call-iptables   # must be 1

# Service rules                             (on: a worker)
iptables -t nat -L KUBE-SERVICES -n | head -20
\end{cmdlst}

\hypertarget{certificate-inspection}{%
\section{Certificate
inspection}\label{certificate-inspection}}

\begin{cmdlst}
openssl x509 -in <cert>.crt -noout -subject          # CN -> username, O -> group
openssl x509 -in <cert>.crt -noout -ext subjectAltName
openssl x509 -in <cert>.crt -noout -dates
openssl verify -CAfile ca.crt <cert>.crt

# What is actually embedded in a kubeconfig?
kubectl config view --kubeconfig=<file> --raw \
  -o jsonpath='{.users[0].user.client-certificate-data}' \
  | base64 -d | openssl x509 -noout -subject
\end{cmdlst}

\hypertarget{service-and-file-locations}{%
\section{Service and
file
locations}\label{service-and-file-locations}}

\begin{longtable}[]{@{}
  >{\raggedright\arraybackslash}p{(\columnwidth - 6\tabcolsep) * \real{0.1250}}
  >{\raggedright\arraybackslash}p{(\columnwidth - 6\tabcolsep) * \real{0.1875}}
  >{\raggedright\arraybackslash}p{(\columnwidth - 6\tabcolsep) * \real{0.2500}}
  >{\raggedright\arraybackslash}p{(\columnwidth - 6\tabcolsep) * \real{0.4375}}@{}}
\toprule\noalign{}
\begin{minipage}[b]{\linewidth}\raggedright
Host
\end{minipage} &
\begin{minipage}[b]{\linewidth}\raggedright
Unit
\end{minipage} &
\begin{minipage}[b]{\linewidth}\raggedright
Binary
\end{minipage} &
\begin{minipage}[b]{\linewidth}\raggedright
Key configuration
\end{minipage} \\
\midrule\noalign{}
\endhead
\bottomrule\noalign{}
\endlastfoot
server &
\texttt{etcd} &
\texttt{/}\allowbreak{}\texttt{usr/}\allowbreak{}\texttt{local/}\allowbreak{}\texttt{bin/}\allowbreak{}\texttt{etcd}
& flags in unit;
data at
\texttt{/var/lib/etcd} \\
server &
\texttt{kube-apiserver}
&
\texttt{/}\allowbreak{}\texttt{usr/}\allowbreak{}\texttt{local/}\allowbreak{}\texttt{bin/}\allowbreak{}\texttt{kube-}\allowbreak{}\texttt{apiserver}
& flags in unit;
credentials in
\texttt{/}\allowbreak{}\texttt{var/}\allowbreak{}\texttt{lib/}\allowbreak{}\texttt{kubernetes/}\allowbreak{} \\
server &
\texttt{kube-}\allowbreak{}\texttt{controller-}\allowbreak{}\texttt{manager}
&
\texttt{/}\allowbreak{}\texttt{usr/}\allowbreak{}\texttt{local/}\allowbreak{}\texttt{bin/}\allowbreak{}\texttt{kube-}\allowbreak{}\texttt{controller-}\allowbreak{}\texttt{manager}
&
\texttt{/}\allowbreak{}\texttt{var/}\allowbreak{}\texttt{lib/}\allowbreak{}\texttt{kubernetes/}\allowbreak{}\texttt{kube-}\allowbreak{}\texttt{controller-}\allowbreak{}\texttt{manager.}\allowbreak{}\texttt{kubeconfig} \\
server &
\texttt{kube-scheduler}
&
\texttt{/}\allowbreak{}\texttt{usr/}\allowbreak{}\texttt{local/}\allowbreak{}\texttt{bin/}\allowbreak{}\texttt{kube-}\allowbreak{}\texttt{scheduler}
&
\texttt{/}\allowbreak{}\texttt{etc/}\allowbreak{}\texttt{kubernetes/}\allowbreak{}\texttt{config/}\allowbreak{}\texttt{kube-}\allowbreak{}\texttt{scheduler.}\allowbreak{}\texttt{yaml} \\
worker &
\texttt{containerd}
&
\texttt{/bin/containerd}
&
\texttt{/}\allowbreak{}\texttt{etc/}\allowbreak{}\texttt{containerd/}\allowbreak{}\texttt{config.}\allowbreak{}\texttt{toml} \\
worker &
\texttt{kubelet} &
\texttt{/}\allowbreak{}\texttt{usr/}\allowbreak{}\texttt{local/}\allowbreak{}\texttt{bin/}\allowbreak{}\texttt{kubelet}
&
\texttt{/}\allowbreak{}\texttt{var/}\allowbreak{}\texttt{lib/}\allowbreak{}\texttt{kubelet/}\allowbreak{}\texttt{kubelet-}\allowbreak{}\texttt{config.}\allowbreak{}\texttt{yaml},
\texttt{/}\allowbreak{}\texttt{var/}\allowbreak{}\texttt{lib/}\allowbreak{}\texttt{kubelet/}\allowbreak{}\texttt{kubeconfig} \\
worker &
\texttt{kube-proxy}
&
\texttt{/}\allowbreak{}\texttt{usr/}\allowbreak{}\texttt{local/}\allowbreak{}\texttt{bin/}\allowbreak{}\texttt{kube-}\allowbreak{}\texttt{proxy}
&
\texttt{/}\allowbreak{}\texttt{var/}\allowbreak{}\texttt{lib/}\allowbreak{}\texttt{kube-}\allowbreak{}\texttt{proxy/}\allowbreak{}\texttt{kube-}\allowbreak{}\texttt{proxy-}\allowbreak{}\texttt{config.}\allowbreak{}\texttt{yaml},
\texttt{/}\allowbreak{}\texttt{var/}\allowbreak{}\texttt{lib/}\allowbreak{}\texttt{kube-}\allowbreak{}\texttt{proxy/}\allowbreak{}\texttt{kubeconfig} \\
\end{longtable}

\hypertarget{ports}{%
\section{Ports}\label{ports}}

\begin{longtable}[]{@{}
  >{\raggedright\arraybackslash}p{(\columnwidth - 4\tabcolsep) * \real{0.3333}}
  >{\raggedright\arraybackslash}p{(\columnwidth - 4\tabcolsep) * \real{0.3333}}
  >{\raggedright\arraybackslash}p{(\columnwidth - 4\tabcolsep) * \real{0.3333}}@{}}
\toprule\noalign{}
\begin{minipage}[b]{\linewidth}\raggedright
Port
\end{minipage} &
\begin{minipage}[b]{\linewidth}\raggedright
Component
\end{minipage} &
\begin{minipage}[b]{\linewidth}\raggedright
Exposure
\end{minipage} \\
\midrule\noalign{}
\endhead
\bottomrule\noalign{}
\endlastfoot
6443 &
kube-apiserver &
\texttt{0.0.0.0} \\
2379 & etcd client &
\texttt{127.0.0.1}
only \\
2380 & etcd peer &
\texttt{127.0.0.1}
only \\
10250 & kubelet API
&
\texttt{0.0.0.0} \\
30000--32767 &
NodePort range &
every node \\
\end{longtable}

\hypertarget{address-spaces}{%
\section{Address
spaces}\label{address-spaces}}

\begin{longtable}[]{@{}
  >{\raggedright\arraybackslash}p{(\columnwidth - 2\tabcolsep) * \real{0.5000}}
  >{\raggedright\arraybackslash}p{(\columnwidth - 2\tabcolsep) * \real{0.5000}}@{}}
\toprule\noalign{}
\begin{minipage}[b]{\linewidth}\raggedright
Range
\end{minipage} &
\begin{minipage}[b]{\linewidth}\raggedright
Purpose
\end{minipage} \\
\midrule\noalign{}
\endhead
\bottomrule\noalign{}
\endlastfoot
Your LAN & Node
addresses \\
\texttt{10.200.0.0/16}
& Cluster CIDR (all
pods) \\
\texttt{10.200.0.0/24}
& node-0 pod
subnet \\
\texttt{10.200.1.0/24}
& node-1 pod
subnet \\
\texttt{10.32.0.0/24}
& Service
ClusterIPs \\
\texttt{10.32.0.1} &
The
\texttt{kubernetes}
Service (API
server) \\
\end{longtable}

\hypertarget{restart-sequence-after-a-credential-change}{%
\section{Restart
sequence after a
credential
change}\label{restart-sequence-after-a-credential-change}}

\begin{cmdlst}
# 1. Regenerate on the jumpbox (certificate AND every kubeconfig embedding it)
# 2. Re-distribute with the relevant Chapter 4 / 5 scp commands
# 3. Restart the consumer:
ssh root@server   'systemctl restart kube-apiserver kube-controller-manager kube-scheduler'
ssh root@node-0   'systemctl restart kubelet kube-proxy'
ssh root@node-1   'systemctl restart kubelet kube-proxy'
\end{cmdlst}

\begin{tipbox}\calloutlabel{teal}{RULE OF THUMB}

None of these
components watch
their credential
files. A restart is
always required.

\end{tipbox}

\hypertarget{app-next}{%
\chapter{What This
Cluster Lacks, and
What to Learn
Next}\label{app-next}}

You now have a
working Kubernetes
cluster. Here is the
honest distance
between it and
something you would
run in production,
roughly in order of
how much it matters.

\hypertarget{availability}{%
\section{Availability}\label{availability}}

\textbf{Single
control plane node.}
One \texttt{server},
one etcd member, one
API server. Losing
that host loses the
cluster's entire
state and its
control plane
simultaneously.
Production runs
three or five etcd
members across
failure domains and
at least three API
servers behind a
load balancer.
\emph{Next step: add
two more etcd
members and convert
\texttt{-\/-initial-}\allowbreak{}\texttt{cluster}
to a three-member
configuration. It is
the single most
instructive
follow-up exercise
in this list.}

\textbf{Pre-release
component versions.}
\texttt{containerd-}\allowbreak{}\texttt{2.}\allowbreak{}\texttt{1.}\allowbreak{}\texttt{0-}\allowbreak{}\texttt{beta.}\allowbreak{}\texttt{0},
\texttt{etcd-}\allowbreak{}\texttt{v3.}\allowbreak{}\texttt{6.}\allowbreak{}\texttt{0-}\allowbreak{}\texttt{rc.}\allowbreak{}\texttt{3},
and
\texttt{runc\ v1.}\allowbreak{}\texttt{3.}\allowbreak{}\texttt{0-}\allowbreak{}\texttt{rc.}\allowbreak{}\texttt{1}
are beta and RC
builds. Pin to GA
releases before
trusting anything to
this stack.

\hypertarget{security}{%
\section{Security}\label{security}}

\textbf{etcd is
unauthenticated
plaintext on
loopback.} Shell
access on
\texttt{server}
equals full
read/write to every
object in the
cluster. \emph{Next
step: enable
\texttt{-\/-client-}\allowbreak{}\texttt{cert-}\allowbreak{}\texttt{auth}
on etcd and the
matching
\texttt{-\/-etcd-cafile}/\texttt{-\/-etcd-certfile}/\texttt{-\/-etcd-keyfile}
flags on the API
server. The
certificates are
already staged in
\texttt{/etc/etcd/}.}

\textbf{One CA for
everything.}
Production separates
the cluster CA, the
etcd CA, and the
front-proxy CA so
that compromising
one does not
compromise all.

\textbf{Ten-year
certificates with no
rotation and no
revocation.} There
is no CRL, no OCSP,
and no automated
renewal. A leaked
key is valid until
2035. \emph{Next
step: explore the
CSR API and kubelet
certificate rotation
--- the controller
manager is already
configured to sign
with
\texttt{-\/-cluster-}\allowbreak{}\texttt{signing-}\allowbreak{}\texttt{cert-}\allowbreak{}\texttt{file}.}

\textbf{Root SSH
enabled on every
cluster machine.} An
explicit convenience
tradeoff. Not
defensible outside a
lab.

\textbf{No
NetworkPolicy
enforcement.} Every
pod can reach every
other pod and every
Service. The
\texttt{bridge} CNI
plugin does not
implement
NetworkPolicy ---
that requires
Calico, Cilium, or
similar. \emph{Next
step: install Calico
and write a
default-deny
policy.}

\textbf{No Pod
Security admission.}
\texttt{-\/-allow-}\allowbreak{}\texttt{privileged=true}
with no Pod Security
Standards means any
pod can request
privileged mode and
escape to the node.
\emph{Next step:
enable Pod Security
admission at the
namespace level with
\texttt{restricted}
enforcement.}

\textbf{No audit
policy.} Audit
logging is enabled
(\texttt{-\/-audit-}\allowbreak{}\texttt{log-}\allowbreak{}\texttt{path})
but with the default
policy, which
records metadata
only. A real policy
would define which
requests to log at
which level.

\hypertarget{networking}{%
\section{Networking}\label{networking}}

\textbf{No cluster
DNS.} No CoreDNS, so
Services are
unreachable by name
---
\texttt{nginx.}\allowbreak{}\texttt{default.}\allowbreak{}\texttt{svc.}\allowbreak{}\texttt{cluster.}\allowbreak{}\texttt{local}
does not resolve.
This is the most
visible gap versus
any real cluster,
and the one most
worth closing first.
\emph{Next step:
deploy CoreDNS and
set
\texttt{clusterDNS}
and
\texttt{clusterDomain}
in
\texttt{kubelet-}\allowbreak{}\texttt{config.}\allowbreak{}\texttt{yaml}.}

\textbf{Static
routes instead of a
CNI network plugin.}
Four hand-written
routes that do not
survive a reboot and
scale as
\emph{n}\textsuperscript{2}.
\emph{Next step:
replace them with
Flannel or Calico
and watch the same
routes appear
automatically.}

\textbf{No Ingress.}
NodePort is the
ceiling. No TLS
termination, no
host- or path-based
routing, no single
entry point.
\emph{Next step:
deploy ingress-nginx
as a DaemonSet with
\texttt{hostNetwork}.}

\textbf{No
LoadBalancer
support.} No cloud
provider
integration.
\emph{Next step:
MetalLB in L2 mode
gives you real
LoadBalancer
Services on a
bare-metal LAN.}

\hypertarget{workload-capabilities}{%
\section{Workload
capabilities}\label{workload-capabilities}}

\textbf{No storage.}
No CSI driver, no
StorageClass, no
PersistentVolumes.
Every workload here
is necessarily
stateless.
\emph{Next step:
deploy the
local-path
provisioner for a
minimal dynamic
StorageClass.}

\textbf{No metrics.}
No metrics-server,
so
\texttt{kubectl\ top}
fails and
HorizontalPodAutoscaler
cannot function.

\textbf{No
observability.} No
log aggregation, no
Prometheus, no
tracing.
\texttt{journalctl}
on each machine is
the entire story.

\textbf{\texttt{maxPods:\ 16}.}
Fine for a 2 GB
node; a hard
scheduling ceiling
at 32 pods
cluster-wide.

\hypertarget{operations}{%
\section{Operations}\label{operations}}

\textbf{No backup.}
\texttt{/var/lib/etcd}
is never
snapshotted.
\emph{Next step:
\texttt{etcdctl\ snapshot\ save},
then practice a
restore. Restoring
etcd is a skill
worth having before
you need it.}

\textbf{No upgrade
path.} Binaries were
placed by hand.
Upgrading means
repeating Chapters
7--9 with new
versions and no
coordination or
drain.

\textbf{Nothing
survives reboot
cleanly.} Pod routes
are ephemeral; swap
may return via
\texttt{/etc/fstab}.

\hypertarget{a-suggested-order}{%
\section{A suggested
order}\label{a-suggested-order}}

\begin{enumerate}
\def\labelenumi{\arabic{enumi}.}
\tightlist
\item
  \textbf{Rebuild it
  from scratch
  without the
  guide.} Everything
  else assumes the
  build is in your
  hands, not your
  notes.
\item
  \textbf{Deploy
  CoreDNS.} Closes
  the most visible
  functional gap and
  teaches you what
  \texttt{clusterDNS}
  in the kubelet
  config actually
  does.
\item
  \textbf{Break it
  deliberately.}
  Stop the scheduler
  and create a
  Deployment.
  Corrupt a
  certificate
  subject. Delete a
  route. Predict the
  symptom before you
  look. This is
  where the
  diagnostic tables
  in this guide stop
  being reference
  material and start
  being intuition.
\item
  \textbf{Add two
  etcd members.}
  Quorum, Raft, and
  the real reason
  odd member counts
  matter.
\item
  \textbf{Enable TLS
  between the API
  server and etcd.}
  A small,
  well-scoped change
  with a clear
  security payoff.
\item
  \textbf{Replace
  the static routes
  with a real CNI
  plugin.} Then read
  its config and
  find the routes it
  created.
\item
  \textbf{Snapshot
  etcd and restore
  it.} Then reason
  about what
  happened to the
  pods that were
  running.
\end{enumerate}

\hypertarget{app-teardown}{%
\chapter{Teardown}\label{app-teardown}}

The upstream lab's
teardown is one
line: delete the
four VMs. Use that
if you can --- it is
complete and leaves
nothing behind.

\begin{beyondbox}\calloutlabel{teal}{BEYOND THE SOURCE}

the sequences below
reset the machines
in place, for when
rebuilding the VMs
is not convenient.
They are additions,
not upstream steps.

\end{beyondbox}

\hypertarget{reset-the-workers}{%
\section{Reset the
workers}\label{reset-the-workers}}

\begin{cmdlst}
# on: node-0, node-1
{
  systemctl stop kubelet kube-proxy containerd
  systemctl disable kubelet kube-proxy containerd
  rm -f /etc/systemd/system/{kubelet,kube-proxy,containerd}.service
  systemctl daemon-reload
  rm -rf /var/lib/kubelet /var/lib/kube-proxy /var/lib/kubernetes \
         /var/run/kubernetes /etc/cni/net.d /opt/cni/bin /etc/containerd
  rm -f /usr/local/bin/{kubelet,kube-proxy,runc,crictl,kubectl}
  rm -f /bin/{containerd,containerd-shim-runc-v2,containerd-stress}
  ip link delete cni0 2>/dev/null
  rm -f /etc/sysctl.d/kubernetes.conf
}
\end{cmdlst}

\texttt{ip\ link\ delete\ cni0}
removes the bridge;
without it a rebuild
with a different pod
subnet inherits a
stale bridge address
and pods fail to get
usable IPs.

\hypertarget{reset-the-control-plane}{%
\section{Reset the
control
plane}\label{reset-the-control-plane}}

\begin{cmdlst}
# on: server
{
  systemctl stop kube-apiserver kube-controller-manager kube-scheduler etcd
  systemctl disable kube-apiserver kube-controller-manager kube-scheduler etcd
  rm -f /etc/systemd/system/{kube-apiserver,kube-controller-manager,kube-scheduler,etcd}.service
  systemctl daemon-reload
  rm -rf /var/lib/etcd /var/lib/kubernetes /etc/etcd /etc/kubernetes
  rm -f /usr/local/bin/{kube-apiserver,kube-controller-manager,kube-scheduler,kubectl,etcd,etcdctl}
  rm -f /var/log/audit.log
}
\end{cmdlst}

\hypertarget{reset-the-jumpbox}{%
\section{Reset the
jumpbox}\label{reset-the-jumpbox}}

\begin{cmdlst}
# on: jumpbox
{
  cd /root/kubernetes-the-hard-way
  rm -f *.crt *.key *.csr *.kubeconfig ca.srl \
        encryption-config.yaml hosts 10-bridge.conf kubelet-config.yaml
  rm -f ~/.kube/config
}
\end{cmdlst}

Keep
\texttt{machines.txt}
and the
\texttt{downloads/}
directory --- the
inventory is still
accurate and
re-downloading 500
MB is pointless.
Everything else
regenerates in under
a minute.

\hypertarget{remove-the-shared-hosts-entries}{%
\section{Remove the
shared hosts
entries}\label{remove-the-shared-hosts-entries}}

\begin{cmdlst}
# on: jumpbox -- removes the block from all four machines
{
  sed -i '/# Kubernetes The Hard Way/,+3d' /etc/hosts
  for host in server node-0 node-1; do
    ssh -n root@${host} "sed -i '/# Kubernetes The Hard Way/,+3d' /etc/hosts"
  done
}
\end{cmdlst}

\texttt{+3d} deletes
the header comment
plus the three
machine lines.
Verify with
\texttt{cat\ /etc/hosts}
before relying on it
--- if you appended
the block twice, run
it twice.

Do this
\textbf{before}
re-running Chapter
3, or the second
build appends a
duplicate block and
name resolution
becomes
order-dependent in a
way that is
genuinely hard to
diagnose.

\hypertarget{remove-stale-ssh-host-keys}{%
\section{Remove
stale SSH host
keys}\label{remove-stale-ssh-host-keys}}

If the machines will
be rebuilt and reuse
the same addresses:

\begin{cmdlst}
# on: jumpbox
{
  for h in server node-0 node-1 \
           server.kubernetes.local node-0.kubernetes.local node-1.kubernetes.local; do
    ssh-keygen -R ${h}
  done
  while read IP FQDN HOST SUBNET; do ssh-keygen -R ${IP}; done < machines.txt
}
\end{cmdlst}

\hypertarget{verify-the-reset}{%
\section{Verify the
reset}\label{verify-the-reset}}

\begin{cmdlst}
# on: jumpbox
{
  for h in server node-0 node-1; do
    echo "--- ${h}"
    ssh -n root@${h} "systemctl list-units --type=service --state=running | grep -E 'kube|etcd|containerd' || echo clean"
  done
}
\end{cmdlst}

\texttt{clean} from
all three means you
are ready to start
again at Chapter 3.

\end{document}
